% Plantilla LaTeX para Trabajo de Integración Curricular - FIEE/EPN
% Compilar con: XeLaTeX -> Biber -> XeLaTeX -> XeLaTeX
% Mantener el archivo referencias_ieee.bib en la misma carpeta de este .tex.

\documentclass[11pt,a4paper]{article}

% -------------------- IDIOMA Y TIPOGRAFÍA --------------------
% Para cumplir estrictamente con Arial se debe compilar con XeLaTeX o LuaLaTeX.
% Si se usa pdfLaTeX, se mantiene Helvetica solo como alternativa de emergencia.
\usepackage{iftex}
\ifPDFTeX
    \usepackage[T1]{fontenc}
    \usepackage[utf8]{inputenc}
    \usepackage{helvet}
    \renewcommand{\familydefault}{\sfdefault}
\else
    \usepackage{fontspec}
    \setmainfont{Arial}
    \setsansfont{Arial}
    \renewcommand{\familydefault}{\sfdefault}
\fi
\usepackage[spanish,es-tabla]{babel}
\usepackage{csquotes}
\usepackage{ragged2e}
\setlength{\JustifyingParindent}{0pt} % evita que \justifying reintroduzca la sangría
\ifPDFTeX
    \DeclareUnicodeCharacter{2212}{-}
    \DeclareUnicodeCharacter{2192}{$\to$}
    \DeclareUnicodeCharacter{2500}{-}
    \DeclareUnicodeCharacter{2502}{|}
    \DeclareUnicodeCharacter{2514}{`-}
    \DeclareUnicodeCharacter{251C}{|-}
    \DeclareUnicodeCharacter{2190}{$\leftarrow$}
\fi
\usepackage{graphicx}

% -------------------- PÁGINA Y ESPACIADO --------------------
\usepackage[
    a4paper,
    top=2.82cm,
    left=2.75cm,
    right=1.75cm,
    bottom=1.62cm,
    footskip=0.80cm
]{geometry}
\usepackage{setspace}
\setstretch{1.5}
\setlength{\parindent}{0pt}
\setlength{\parskip}{0pt}
\AtBeginDocument{\setlength{\parindent}{0pt}}
\setlength{\emergencystretch}{3em}
\hfuzz=1000pt
\hbadness=10000
\AtBeginDocument{\justifying}

% -------------------- PAQUETES GENERALES --------------------
\usepackage{amsmath}
\usepackage{amssymb}
\usepackage{graphicx}
\usepackage{array}
\usepackage{booktabs}
\usepackage{enumitem}
\usepackage{float}
\usepackage{fancyvrb}
\usepackage{caption}
\usepackage{chngcntr}
\usepackage{etoolbox}
\usepackage{titlesec}
\usepackage{tocloft}
\usepackage{fancyhdr}
\usepackage[hidelinks]{hyperref}


% -------------------- NUMERACIÓN DE PÁGINA --------------------
\newcounter{frontmatterlastpage}
\pagestyle{fancy}
\fancyhf{}
\fancyfoot[C]{\fontsize{11}{13}\selectfont\thepage}
\renewcommand{\headrulewidth}{0pt}
\renewcommand{\footrulewidth}{0pt}
\fancypagestyle{plain}{%
  \fancyhf{}%
  \fancyfoot[C]{\fontsize{11}{13}\selectfont\thepage}%
  \renewcommand{\headrulewidth}{0pt}%
  \renewcommand{\footrulewidth}{0pt}%
}

% -------------------- TÍTULOS: NIVELES 1 A 5 --------------------
\setcounter{secnumdepth}{5}
\setcounter{tocdepth}{4}

\titleformat{\section}[block]
  {\normalfont\bfseries\fontsize{16}{19}\selectfont}
  {\thesection}{1em}{\MakeUppercase}
\titlespacing*{\section}{0pt}{0pt}{12pt}

\titleformat{\subsection}[block]
  {\normalfont\bfseries\fontsize{14}{17}\selectfont}
  {\thesubsection}{1em}{\MakeUppercase}
\titlespacing*{\subsection}{0pt}{12pt}{6pt}

\titleformat{\subsubsection}[block]
  {\normalfont\bfseries\fontsize{12}{15}\selectfont}
  {\thesubsubsection}{1em}{\MakeUppercase}
\titlespacing*{\subsubsection}{0pt}{12pt}{6pt}

\titleformat{\paragraph}[block]
  {\normalfont\bfseries\fontsize{12}{15}\selectfont}
  {\theparagraph}{1em}{}
\titlespacing*{\paragraph}{0pt}{12pt}{6pt}

\titleformat{\subparagraph}[block]
  {\normalfont\itshape\fontsize{11}{13}\selectfont}
  {\thesubparagraph}{1em}{}
\titlespacing*{\subparagraph}{0pt}{12pt}{6pt}

% -------------------- ÍNDICE DE CONTENIDO --------------------
\addto\captionsspanish{\renewcommand{\contentsname}{ÍNDICE DE CONTENIDO}}
\addto\captionsspanish{\renewcommand{\listtablename}{ÍNDICE DE TABLAS}}
\addto\captionsspanish{\renewcommand{\listfigurename}{ÍNDICE DE FIGURAS}}
\renewcommand{\cfttoctitlefont}{\hfill\bfseries\fontsize{16}{19}\selectfont}
\renewcommand{\cftaftertoctitle}{\hfill\vspace{12pt}}
\renewcommand{\cftsecfont}{\bfseries}
\renewcommand{\cftsecpagefont}{\bfseries}
\renewcommand{\cftsecleader}{\cftdotfill{\cftdotsep}}
\setlength{\cftbeforesecskip}{2pt}
\setlength{\cftbeforesubsecskip}{0pt}

% -------------------- TABLAS, FIGURAS Y ECUACIONES --------------------
\counterwithout{table}{section}
\counterwithout{figure}{section}
\counterwithin{equation}{section}
\captionsetup{font=normalsize,labelfont=normalfont,labelsep=period,justification=centering,singlelinecheck=false,hypcap=false}
\captionsetup[table]{position=above}
\captionsetup[figure]{position=below}
\DeclareCaptionType{listing}[Listado][ÍNDICE DE LISTADOS]
\AtBeginEnvironment{table}{\singlespacing}
\AtBeginEnvironment{tabular}{\singlespacing}

% -------------------- REFERENCIAS IEEE --------------------
\usepackage[
    backend=biber,
    style=ieee,
    sorting=none,
    giveninits=true,
    doi=true,
    url=true,
    isbn=false
]{biblatex}
\addbibresource{referencias_ieee.bib}

% -------------------- DATOS EDITABLES DE PORTADA --------------------
\newcommand{\institucion}{ESCUELA POLITÉCNICA NACIONAL}
\newcommand{\facultad}{FACULTAD DE INGENIERÍA ELÉCTRICA Y ELECTRÓNICA}
\newcommand{\tituloTrabajo}{ANÁLISIS DE HERRAMIENTAS SOFTWARE PARA LA GESTIÓN DE ETIQUETADO DE VIDEOS}
\newcommand{\subtituloTrabajo}{DESARROLLO DE UN MÓDULO (FUNCIONALIDAD B) PARA UNA HERRAMIENTA SOFTWARE QUE GESTIONA LAS ANOTACIONES DE VIDEOS PARA LA LENGUA DE SEÑAS ECUATORIANA (LSEc)}
\newcommand{\tipoTrabajo}{TRABAJO DE INTEGRACIÓN CURRICULAR PRESENTADO COMO REQUISITO PARA LA OBTENCIÓN DEL TÍTULO DE INGENIERO EN TECNOLOGÍAS DE LA INFORMACIÓN}
\newcommand{\autorUno}{Francisco Sebastián Imbaquinga Espinosa}
\newcommand{\correoAutorUno}{francisco.imbaquinga@epn.edu.ec}
\newcommand{\directorTIC}{Xavier Alexander Calderón Hinojosa}
\newcommand{\correoDirector}{xavier.calderon@epn.edu.ec}
\newcommand{\lugarFecha}{DMQ, JULIO 2026}

% -------------------- COMANDOS AUXILIARES --------------------
\newcommand{\frontsection}[1]{%
    \clearpage
    \phantomsection
    \addcontentsline{toc}{section}{#1}
    {\centering\bfseries\fontsize{16}{19}\selectfont #1\par}
    \vspace{18pt}
}

\newcommand{\firma}[2]{%
    \vspace{1.8cm}
    \begin{center}
        \rule{7cm}{0.4pt}\\
        #1\\
        #2
    \end{center}
}

\newcommand{\firmaDerecha}[2]{%
    \vspace{2.5cm}
    \begin{flushright}
        \begin{minipage}{0.45\textwidth}
            \centering
            \rule{\linewidth}{0.4pt}\\[-0.1cm]
            #1\\
            \textbf{#2}
        \end{minipage}
    \end{flushright}
}

% ============================================================
% Ajustes de formato solicitados en la revisión del tutor.
%
% CÓMO USARLO: agregar la línea
%     % ============================================================
% Ajustes de formato solicitados en la revisión del tutor.
%
% CÓMO USARLO: agregar la línea
%     % ============================================================
% Ajustes de formato solicitados en la revisión del tutor.
%
% CÓMO USARLO: agregar la línea
%     \input{preambulo_correcciones}
% en el preámbulo del documento principal (main.tex), después de
% cargar los paquetes habituales (caption, fancyvrb, float, etc.)
% y ANTES de \begin{document}.
% ============================================================

% ------------------------------------------------------------
% 1. Numeración de tablas y figuras por capítulo:
%    Tabla 2.1, Tabla 2.2, ..., Figura 3.1, Figura 3.2, ...
%    (el número de capítulo corresponde al contador de \section)
% ------------------------------------------------------------
\usepackage{chngcntr}
\counterwithin{table}{section}
\counterwithin{figure}{section}

% ------------------------------------------------------------
% 2. Tipo flotante "Código", numerado por capítulo (Código 2.1,
%    Código 3.1, ...). Reemplaza al antiguo "Listado".
%    Se usa en el texto con: \captionof{codigo}{Descripción}
%    Requiere el paquete caption (si el documento principal ya lo
%    carga, la línea \usepackage{caption} puede eliminarse).
% ------------------------------------------------------------
\usepackage{caption}
\DeclareCaptionType[within=section]{codigo}[Código][Índice de códigos]

% ------------------------------------------------------------
% 3. TikZ, usado por la figura del ciclo de desarrollo y por el
%    diagrama de clases del middleware (capítulo de Metodología).
% ------------------------------------------------------------
\usepackage{tikz}
\usetikzlibrary{arrows.meta, positioning, shapes.geometric, shapes.multipart, calc, fit}

% en el preámbulo del documento principal (main.tex), después de
% cargar los paquetes habituales (caption, fancyvrb, float, etc.)
% y ANTES de \begin{document}.
% ============================================================

% ------------------------------------------------------------
% 1. Numeración de tablas y figuras por capítulo:
%    Tabla 2.1, Tabla 2.2, ..., Figura 3.1, Figura 3.2, ...
%    (el número de capítulo corresponde al contador de \section)
% ------------------------------------------------------------
\usepackage{chngcntr}
\counterwithin{table}{section}
\counterwithin{figure}{section}

% ------------------------------------------------------------
% 2. Tipo flotante "Código", numerado por capítulo (Código 2.1,
%    Código 3.1, ...). Reemplaza al antiguo "Listado".
%    Se usa en el texto con: \captionof{codigo}{Descripción}
%    Requiere el paquete caption (si el documento principal ya lo
%    carga, la línea \usepackage{caption} puede eliminarse).
% ------------------------------------------------------------
\usepackage{caption}
\DeclareCaptionType[within=section]{codigo}[Código][Índice de códigos]

% ------------------------------------------------------------
% 3. TikZ, usado por la figura del ciclo de desarrollo y por el
%    diagrama de clases del middleware (capítulo de Metodología).
% ------------------------------------------------------------
\usepackage{tikz}
\usetikzlibrary{arrows.meta, positioning, shapes.geometric, shapes.multipart, calc, fit}

% en el preámbulo del documento principal (main.tex), después de
% cargar los paquetes habituales (caption, fancyvrb, float, etc.)
% y ANTES de \begin{document}.
% ============================================================

% ------------------------------------------------------------
% 1. Numeración de tablas y figuras por capítulo:
%    Tabla 2.1, Tabla 2.2, ..., Figura 3.1, Figura 3.2, ...
%    (el número de capítulo corresponde al contador de \section)
% ------------------------------------------------------------
\usepackage{chngcntr}
\counterwithin{table}{section}
\counterwithin{figure}{section}

% ------------------------------------------------------------
% 2. Tipo flotante "Código", numerado por capítulo (Código 2.1,
%    Código 3.1, ...). Reemplaza al antiguo "Listado".
%    Se usa en el texto con: \captionof{codigo}{Descripción}
%    Requiere el paquete caption (si el documento principal ya lo
%    carga, la línea \usepackage{caption} puede eliminarse).
% ------------------------------------------------------------
\usepackage{caption}
\DeclareCaptionType[within=section]{codigo}[Código][Índice de códigos]

% ------------------------------------------------------------
% 3. TikZ, usado por la figura del ciclo de desarrollo y por el
%    diagrama de clases del middleware (capítulo de Metodología).
% ------------------------------------------------------------
\usepackage{tikz}
\usetikzlibrary{arrows.meta, positioning, shapes.geometric, shapes.multipart, calc, fit}

% ============================================================
% DOCUMENTO
% ============================================================
\begin{document}

% -------------------- PORTADA --------------------
\hypersetup{pageanchor=false}
\begin{titlepage}
    \thispagestyle{empty}
    \begin{center}
        \bfseries
        \institucion\\[2.2cm]
        \facultad\\[2.2cm]

        \MakeUppercase{\tituloTrabajo}\\[0.8cm]
        \MakeUppercase{\subtituloTrabajo}\\[1.8cm]

        \MakeUppercase{\tipoTrabajo}\\[2.1cm]

        \normalfont\bfseries
        \autorUno\\
        \correoAutorUno\\[1.5cm]

        \directorTIC\\
        \correoDirector\\[0.5cm]

        \vfill
        \lugarFecha
    \end{center}
\end{titlepage}
\hypersetup{pageanchor=true}

% -------------------- PRELIMINARES CON NÚMEROS ROMANOS --------------------
\pagenumbering{Roman}
\setcounter{page}{1}

\frontsection{CERTIFICACIONES}
Yo, \autorUno, declaro que el trabajo de integración curricular aquí descrito es de mi autoría; que no ha sido previamente presentado para ningún grado o calificación profesional; y, que he consultado las referencias bibliográficas que se incluyen en este documento.

Asimismo, declaro que he utilizado la herramienta de inteligencia artificial ChatGPT y Claude únicamente como apoyo en [USO ESPECÍFICO], sin atribuirle autoría, y que todo el contenido derivado ha sido revisado, validado y es de mi exclusiva responsabilidad.

\firma{\autorUno}{AUTOR}
\vspace{2\baselineskip}

Certifico que el presente trabajo de integración curricular fue desarrollado por \autorUno, bajo mi supervisión.

\firmaDerecha{\directorTIC}{DIRECTOR}

\frontsection{DECLARACIÓN DE AUTORÍA}
A través de la presente declaración, afirmamos que el trabajo de integración curricular aquí descrito, así como el (los) producto(s) resultante(s) del mismo, son públicos y estarán a disposición de la comunidad a través del repositorio institucional de la Escuela Politécnica Nacional; sin embargo, la titularidad de los derechos patrimoniales nos corresponde a los autores que hemos contribuido en el desarrollo del presente trabajo; observando para el efecto las disposiciones establecidas por el órgano competente en propiedad intelectual, la normativa interna y demás normas.\\
\\

Francisco Imbaquinga
\vspace{1\baselineskip}

Xavier Calderón


\frontsection{DEDICATORIA}
(Opcional)

\frontsection{AGRADECIMIENTO}
(Opcional)

\clearpage
\phantomsection
\addcontentsline{toc}{section}{ÍNDICE DE CONTENIDO}
\tableofcontents

\clearpage
\phantomsection
\addcontentsline{toc}{section}{ÍNDICE DE TABLAS}
\listoftables

\clearpage
\phantomsection
\addcontentsline{toc}{section}{ÍNDICE DE FIGURAS}
\listoffigures

\clearpage
\phantomsection
\addcontentsline{toc}{section}{ÍNDICE DE COMANDOS}
\listofcomandos

\clearpage
\phantomsection
\addcontentsline{toc}{section}{ÍNDICE DE CÓDIGOS}
\listofcodigos

\frontsection{RESUMEN}
La anotación de videos en Lengua de Señas Ecuatoriana (LSEc) exige herramientas que respeten el criterio lingüístico del investigador y, al mismo tiempo, permitan aprovechar modelos modernos de Inteligencia Artificial. ELAN, una de las herramientas más utilizadas para la anotación multimodal, facilita segmentar y etiquetar eventos sobre una línea de tiempo. No obstante, su arquitectura basada en Java y su mecanismo de extensión dificultan la conexión directa con servicios de inferencia desarrollados en entornos actuales, principalmente Python, con dependencias especializadas. Esto limita la incorporación de procesos automáticos, como la segmentación temporal y la clasificación de glosas, al flujo habitual de anotación.

Para resolver esta limitación, en este trabajo se desarrolló un middleware que funciona como un intermediario entre ELAN y los modelos de Inteligencia Artificial, sin necesidad de modificar a fondo la herramienta. La idea central es sencilla: ELAN actúa como cliente y envía la solicitud de análisis, mientras que el middleware recibe esa solicitud, verifica que los datos sean correctos, ubica el modelo que debe usarse y coordina su ejecución en un entorno aislado. Finalmente, el resultado regresa a ELAN convertido en segmentos de tiempo que el investigador puede aceptar, ajustar o eliminar. De este modo, la revisión final siempre queda en manos del especialista.

El desarrollo se llevó a cabo por etapas, comprobando cada avance antes de continuar con el siguiente. Primero se estudió cómo funciona ELAN por dentro para entender la forma correcta de conectarle un proceso externo. Con esa base se definió una organización en tres niveles: ELAN como herramienta de anotación, el middleware como coordinador local y los modelos como servicios independientes que procesan los videos. A partir de ahí se fueron sumando capacidades: la comunicación básica, la instalación de modelos, su ejecución controlada, el seguimiento de cada tarea y, por último, la conexión visual dentro de ELAN.

Los resultados confirmaron que la integración funciona y es estable. El sistema permitió instalar modelos, comprobar que estén disponibles, ejecutar análisis sobre los videos, informar en qué estado se encuentra cada tarea y distinguir con claridad entre un error real y una demora por tiempo agotado. Las pruebas realizadas mostraron que el middleware responde de forma ordenada tanto en situaciones normales como ante problemas frecuentes, por ejemplo un video que no existe, un modelo desactivado o un contenedor detenido, sin que ello afecte a ELAN. En conjunto, el trabajo entregó una solución local, ordenada y ampliable que acerca los modelos de Inteligencia Artificial al trabajo diario de anotación, mantiene los datos dentro del equipo del investigador y deja el camino preparado para incorporar nuevos modelos en el futuro sin tener que construir otra herramienta.

\textbf{PALABRAS CLAVE:} ELAN, middleware, orquestación de servicios, Inteligencia Artificial, Lengua de Señas Ecuatoriana.

\frontsection{ABSTRACT}
Video annotation in Ecuadorian Sign Language (LSEc) requires tools that preserve the researcher's linguistic judgment while making it possible to use modern Artificial Intelligence models. ELAN is widely used for multimodal annotation because it supports segmenting and labeling events on a timeline. However, its Java-based architecture and extension mechanism make direct integration difficult with inference services developed in current environments, especially Python, with specialized dependencies. This limits bringing automatic processes, such as temporal segmentation and gloss classification, into the regular annotation workflow.

To address this limitation, this work developed a middleware that acts as an intermediary between ELAN and Artificial Intelligence models, without having to deeply modify the tool. The core idea is simple: ELAN works as a client and sends the analysis request, while the middleware receives that request, checks that the data is correct, locates the model to be used, and coordinates its execution in an isolated environment. Finally, the result returns to ELAN as time segments that the researcher can accept, adjust, or delete. In this way, the final review always remains in the hands of the specialist.

The development was carried out in stages, moving step by step and verifying each advance before continuing with the next one. First, the inner workings of ELAN were studied to understand the correct way to connect an external process to it. On that basis, a three-level organization was defined: ELAN as the annotation tool, the middleware as the local coordinator, and the models as independent services that process the videos. From there, capabilities were added progressively: basic communication, model installation, their controlled execution, the tracking of each task, and, finally, the visual connection within ELAN.

The results confirmed that the integration works and is stable. The system made it possible to install models, verify that they are available, run analyses on the videos, report the state of each task, and clearly distinguish a real error from a delay due to a timeout. The tests showed that the middleware responds in an orderly way both in normal situations and when facing common problems, such as a video that does not exist, a disabled model, or a stopped container, without affecting ELAN. Overall, the work delivered a local, orderly, and extensible solution that brings Artificial Intelligence models closer to the daily annotation work, keeps the data inside the researcher's own machine, and leaves the path prepared to incorporate new models in the future without having to build another tool.

\textbf{KEYWORDS:} ELAN, middleware, service orchestration, Artificial Intelligence, Ecuadorian Sign Language.
\setcounter{frontmatterlastpage}{\value{page}}

% -------------------- CUERPO CON NÚMEROS ARÁBIGOS --------------------
\clearpage
\pagenumbering{arabic}
\setcounter{page}{1}

% El argumento opcional define cómo aparece el título en el índice
% (los capítulos van en mayúsculas; los subtítulos quedan como están).
\section[INTRODUCCIÓN]{Introducción}
La Lengua de Señas Ecuatoriana (LSEc) representa una parte fundamental del patrimonio lingüístico y cultural del país, siendo de suma importancia para la comunicación de personas con dificultades auditivas y su integración con la sociedad. El estudio sistemático de la LSEc ha ganado relevancia en estos últimos años, especialmente en el contexto de investigaciones en Inteligencia Artificial y el procesamiento de lenguaje natural. Para el análisis profundo de videos de LSEc, gran parte de la comunidad científica ha adoptado herramientas de anotación multimodal como ELAN (EUDICO Linguistic Annotator) que permite segmentar y etiquetar eventos lingüísticos sobre la línea de tiempo. Sin embargo, la arquitectura cerrada y obsoleta de ELAN basada en Java dificulta la integración con modelos de Inteligencia Artificial, ya que necesitan plataformas modernas como Python y una gran cantidad de potencia de cálculo.

Este trabajo de integración curricular tiene como objetivo cerrar la brecha entre ELAN y los servicios de Inteligencia Artificial desarrollando un middleware de orquestación que actúa como intermediario. Se trata de un software independiente de ELAN, no de una librería incorporada a la herramienta: ELAN cumple el rol de cliente y el middleware, junto con los modelos que administra, el de servidor local de inferencias. Esta solución está diseñada para un entorno local, lo que permite que ELAN envíe solicitudes de análisis de video a modelos externos sin requerir modificaciones profundas en la arquitectura base del software. El middleware proporciona una interfaz REST, valida las solicitudes del cliente Java al que ELAN está conectado, consulta la lista local de todos los modelos disponibles, coordina la comunicación entre los backends que se ejecutan en contenedores Docker y devuelve un segmento de línea de tiempo que se puede convertir en anotaciones editables. La arquitectura utilizada garantiza la soberanía de los datos al ejecutarse en el entorno local sin depender de servicios en la nube, preserva la estabilidad de ELAN y mantiene aisladas las dependencias propias de cada backend de inferencia.

La implementación del middleware se estructuró en fases iterativas que partieron de una API funcional, continuaron con el registro e instalación local de modelos, la coordinación de backends mediante Docker, el control de trabajos de inferencia y la integración final con ELAN mediante un reconocedor Java. El sistema desarrollado resuelve la incompatibilidad entre la arquitectura cerrada de ELAN basada en Java y entornos modernos de servicios en Python, permitiendo que los investigadores mantengan su flujo de anotación en ELAN y accedan a modelos externos de Inteligencia Artificial sin depender de aplicaciones separadas ni de servicios remotos. El alcance del trabajo se centra en la integración, orquestación y recuperación de resultados, no en el diseño interno, entrenamiento o pipeline específico de procesamiento de cada modelo.


\subsection{Objetivo general}
Desarrollo de un middleware de orquestación de servicios de Inteligencia Artificial diseñado para la herramienta de anotación ELAN.

\subsection{Objetivos específicos}

\begin{enumerate}
    \item Analizar el mecanismo de extensión AVATecH de ELAN mediante el estudio de la documentación oficial e ingeniería inversa del código fuente para definir los contratos de interfaz REST y el esquema de instalación de modelos que permitan la integración de servicios externos de inferencias en el flujo de anotación.
    \item Codificar el núcleo del orquestador de servicios utilizando Python y tecnologías de virtualización ligera (Docker), haciendo uso de mecanismos de control de concurrencia y gestión del ciclo de vida de los contenedores para asegurar la estabilidad del sistema durante la inferencia simultánea de modelos.
    \item Validar el desempeño técnico del middleware mediante pruebas de carga en escenarios de anotación de LSEc, midiendo latencia de comunicación y la correcta recuperación de errores presentados durante fallos de los modelos.

\end{enumerate}

\subsection{Alcance}

El alcance del trabajo está atado a la definición, desarrollo y validación de la capa de infraestructura lógica y se ejecutará a través de una metodología iterativa dividida en fases consolidadas.

En la fase de planteamiento se realizará el estudio sobre soluciones ya implementadas para optimización de procesos en LSEc así como, estudios de la especificación técnica de ELAN con el fin de proponer soluciones de alto impacto.

Por otro lado, en la fase de implementación se procederá con la codificación del núcleo del middleware orquestador. Esta fase se centrará en la configuración del canal de comunicación (HTTP/REST) entre la extensión Java de ELAN y el servidor Python del middleware. También en asegurar que la gestión de dependencias sea agnóstica al sistema operativo del usuario final mediante la ejecución de los servicios externos en contenedores. 

Por último, en la fase de validación se desplegará el sistema en un entorno controlado para realizar pruebas de caja blanca y caja negra. Se verificará la correcta orquestación de un modelo de inteligencia artificial sobre videos LSEc, registrando métricas de latencia, consumo de recursos del sistema y documentando la arquitectura final para garantizar la mantenibilidad y escalabilidad del sistema desarrollado. Es importante recalcar que el proyecto no contempla la creación o entrenamiento de modelos de Inteligencia Artificial, ni la ampliación a herramientas de anotación distintas de ELAN.

\subsection{Marco teórico}

El marco teórico reúne los conceptos necesarios para comprender la integración entre herramientas de anotación multimodal, servicios de inferencia y mecanismos de orquestación local. Se abordan los fundamentos de la anotación temporal en ELAN, la documentación computacional de la Lengua de Señas Ecuatoriana (LSEc), el reconocimiento automático como servicio de apoyo, los contratos REST, el uso de contenedores, la orquestación de servicios de inteligencia artificial y los criterios de interoperabilidad, trazabilidad y mantenibilidad aplicables a este tipo de arquitectura.

\subsubsection{Herramientas de anotación multimodal y ELAN}

La anotación multimodal consiste en asociar descripciones simbólicas con intervalos de tiempo de una señal audiovisual. En un video de lengua de señas no basta con registrar una traducción general de lo observado; también se requiere indicar cuándo inicia y termina una seña, qué mano interviene, qué elementos no manuales aparecen, si existe una pausa, si hay un comentario de revisión o si se incluye una traducción libre del segmento. Por esta razón, la anotación multimodal se apoya en tres componentes principales: el archivo multimedia, los niveles de anotación y las marcas temporales que conectan cada anotación con un fragmento específico de audio o video \cite{wittenburg2006elan}.

ELAN es una herramienta especializada para este tipo de análisis y se ha utilizado ampliamente en estudios de multimodalidad, documentación lingüística, análisis conversacional y lenguas de señas. La herramienta ofrece un entorno visual para reproducir material audiovisual, navegar por la línea de tiempo, crear anotaciones y organizar la información en niveles de anotación. En ELAN, estos niveles se denominan \textit{tiers}; en este documento se utiliza el término \textbf{nivel de anotación} para referirse de forma unificada a esa unidad de organización. Cada nivel de anotación agrupa un tipo de información, por ejemplo glosas, comentarios de revisión o traducciones, de modo que distintos datos lingüísticos puedan sincronizarse con el mismo material audiovisual sin mezclarse entre sí \cite{wittenburg2006elan}.

La unidad básica de trabajo en ELAN es la anotación de un intervalo. Cada intervalo posee un tiempo de inicio, un tiempo de fin, generalmente expresados en milisegundos, y un valor textual que describe el evento observado. Esta estructura exige que cualquier resultado automático destinado a ELAN conserve una semántica temporal precisa: no basta con producir una etiqueta, sino que dicha etiqueta debe estar asociada a un intervalo válido de la línea de tiempo. En otras palabras, el valor de la
inferencia automática depende de la calidad del modelo externo y de la capacidad del sistema
de integración para mantener la semántica temporal que se supone que ELAN debe proporcionar \cite{wittenburg2006elan}.

El formato nativo de ELAN es EAF (Formato de Anotación ELAN), un archivo XML que almacena la estructura del documento de anotación, los niveles de anotación, los archivos multimedia, las marcas de tiempo y las anotaciones\cite{wittenburg2006elan}. Los resultados de un servicio de inferencia deben respetar esta lógica temporal antes de que puedan ser utilizados para la anotación. Una predicción externa debe desglosarse en segmentos con un inicio, un final y un valor textual para que ELAN pueda agregarlos a la línea de tiempo sin convertir manualmente el archivo XML.

ELAN también considera la posibilidad de interactuar con procesos externos a través de reconocedores. En particular, AVATecH define una especificación para incorporar componentes de reconocimiento que reciben parámetros de ELAN, procesan una señal (audio o video) y devuelven segmentos ya catalogados. Este enfoque muestra que la anotación semiautomática forma parte de la evolución natural de ELAN, por lo tanto la herramienta no se limita a la anotación manual, sino que puede apoyarse en detectores o reconocedores externos siempre que estos respeten contratos de entrada y salida comprensibles para el entorno de anotación \cite{auer2010elan,auer2014avatech}.

Desde esta perspectiva, AVATecH constituye un punto de entrada entre ELAN y procesos de reconocimiento externos. La herramienta mantiene su lógica de anotación y el archivo EAF conserva su estructura, mientras el reconocedor actúa como intermediario para solicitar procesamiento, recibir segmentos candidatos y entregarlos al flujo habitual de revisión. De esta manera, la comunicación HTTP, la validación de solicitudes, la ejecución del servicio de inferencia y el manejo de errores pueden ubicarse fuera del núcleo de ELAN, sin alterar la semántica del documento de anotación \cite{auer2014avatech}.

\subsubsection{Lengua de Señas Ecuatoriana: contexto y documentación computacional}

La Lengua de Señas Ecuatoriana, comúnmente referida como LSEc, es una lengua visual-gestual utilizada por las personas sordas en Ecuador. Al igual que otras lenguas de señas, no constituye una transcripción manual del español, sino un sistema lingüístico con recursos expresivos propios, producidos mediante configuraciones manuales, movimiento, orientación, ubicación espacial y aspectos no manuales. La documentación de esta lengua es relevante para la investigación lingüística, las tecnologías de accesibilidad, la educación y la inclusión social \cite{ureta2022lsec_inclusion}.

La documentación de una lengua de señas presenta dificultades particulares porque el fenómeno lingüístico ocurre en el tiempo y el espacio. A diferencia del texto escrito, el video puede contener múltiples rasgos simultáneos: movimiento de las manos, expresión facial, dirección de la mirada, postura corporal y contexto comunicativo. Por lo mismo, herramientas como ELAN resultan adecuadas para documentar la LSEc, ya que permiten distribuir la información en niveles de anotación y mantener la relación de cada observación con el instante exacto del material audiovisual \cite{wittenburg2006elan,ureta2022lsec_inclusion}.

En Ecuador existen trabajos sobre el reconocimiento computacional de LSEc y prototipos para traducir signos en texto o voz bajo condiciones controladas. Estos estudios demuestran que el procesamiento de visión por computadora y el aprendizaje automático en LSEc es un campo activo, especialmente para la clasificación, interpretación o apoyo a la comunicación \cite{guevara2022lsec_interprete,cruz2025prototipo}. En esta cadena tecnológica, la interoperabilidad de los modelos de reconocimiento y las herramientas de anotación es un problema complementario dado que los resultados de un modelo necesitan transformarse en evidencia que pueda ser revisada y ajustada por especialistas.

Hay una diferencia entre el reconocimiento automático y la integración con herramientas de anotación. La creación de corpus de entrenamiento, arquitecturas de aprendizaje profundo, ajuste de hiperparámetros y evaluación lingüística del modelo son parte del campo especial del reconocimiento automático. La interoperabilidad se refiere a cómo transferir resultados temporales, etiquetas y metadatos a un entorno de revisión. Este enfoque mantiene al lingüista como responsable humano en el proceso de validación y aceptación de las anotaciones.

El LSEc también introduce necesidades éticas y operativas. Los videos utilizados en procesos de documentación podrían contener personas identificables y datos sensibles de una comunidad lingüística particular. Por lo tanto, una arquitectura local donde la herramienta de anotación, el middleware y los servicios de inferencia se ejecuten en una estación de trabajo controlada o infraestructura institucional minimizará la exposición de datos a servicios remotos. Esto no sustituye el consentimiento informado o las políticas responsables de gestión de datos, pero reducirá la transferencia innecesaria de material audiovisual a plataformas externas\cite{shi2016edge}.

\subsubsection{Reconocimiento automático de gestos y señas como servicio de apoyo}

El reconocimiento automático de gestos y signos puede verse como un servicio externo que recibe un video y genera resultados candidatos para la anotación. Estos resultados pueden incluir marcos de tiempo donde se detecta actividad relevante, la etiqueta principal para el segmento y, en algunos casos, diferentes resultados con puntuaciones o probabilidades. Para que estos resultados sean útiles en una herramienta como ELAN, el servicio de inferencia necesita tener un contrato de entrada y salida que le permita interpretar la respuesta y convertirla en anotaciones compatibles con la línea de tiempo \cite{auer2014avatech}.

Desde un punto de vista funcional, el reconocimiento de video generalmente se divide en dos tareas generales. La primera es la segmentación temporal, que se encarga de determinar dónde comienza y termina una unidad particular de interés. La segunda es la clasificación, que asigna una etiqueta y un glosa al segmento. Esta separación es útil para las herramientas de anotación, ya que ELAN requiere una etiqueta asociada a un intervalo específico. Si un servicio devuelve una clasificación sin tiempos, el resultado no puede colocarse correctamente en la línea de tiempo; si devuelve una clasificación sin etiqueta, el investigador aún necesita completar la descripción del evento \cite{guevara2022lsec_interprete,cruz2025prototipo}.

Las arquitecturas de aprendizaje profundo pueden variar entre los servicios de inferencia. Un backend puede ser bibliotecas de visión por computadora, modelos secuenciales, redes neuronales profundas o combinaciones multimodales. Pero para una arquitectura de integración, esos detalles deben permanecer detrás de una interfaz estable. Esta separación ayuda a evitar acoplar la herramienta de anotación a una técnica específica y permite reemplazar o actualizar servicios futuros siempre que mantengan el contrato de comunicación \cite{fielding2000rest}.

La salida de un servicio de reconocimiento debe ser postprocesada antes de ser una anotación útil. Los tiempos deben darse en unidades compatibles con ELAN, los intervalos deben tener una duración válida, las etiquetas deben ser serializadas como texto y las respuestas deben contener suficiente información para separar un buen resultado de un error. Cuando el servicio proporciona una predicción diferente, debe tratarse como información de apoyo para la revisión humana, no como una anotación final. En este sentido, la anotación asistida sigue siendo semiautomática: la automatización será menos repetitiva, pero esto no reemplaza la validación experta \cite{auer2010elan}.

El reconocimiento como servicio también requiere robustez. Un modelo puede tardar más de lo esperado, fallar por falta de recursos, recibir un archivo inexistente o devolver una respuesta malformada. Si una herramienta de anotación se comunicara directamente con cada backend, necesitaría implementar una lógica diferente para manejar todos estos casos. Un middleware puede normalizar estos comportamientos a través de estados de trabajo, trazas, mensajes de error y respuestas coherentes, por lo que la integración no depende de que todos los servicios de inferencia se comporten de manera idéntica \cite{fastapi2026features,pydantic2026models}.

\subsubsection{Patrones de integración y rol del middleware}

La integración de ELAN y los servicios actuales de inteligencia artificial no es uniforme. ELAN es una aplicación de escritorio desarrollada en Java con su propia interfaz gráfica y mecanismos de extensión. Los servicios de inferencia suelen ejecutarse como procesos independientes con diferentes dependencias, diferentes entornos de ejecución y, por lo tanto, actualizaciones más frecuentes. Vincular estas dependencias al código de ELAN, sin embargo, haría que fueran más difíciles de distribuir y requeriría una modificación de la herramienta de anotación para cada cambio de backend \cite{dragoni2017microservices}.

Antes de describir el papel de esta capa conviene precisar el término. En ingeniería de software, \textit{middleware} designa la capa de software que se ubica entre aplicaciones independientes y les proporciona servicios de comunicación e intercambio de datos, de modo que ninguna dependa de los detalles internos de la otra \cite{bernstein1996middleware}. Si bien el término se originó en el contexto de los sistemas distribuidos y de las redes de comunicación, su uso consolidado en arquitectura de software es más amplio: abarca cualquier servicio intermedio que conecta programas construidos con tecnologías distintas, traduce formatos de datos, coordina recursos compartidos y oculta la heterogeneidad de los extremos \cite{bernstein1996middleware}.

Conviene además distinguir este concepto de otros términos cercanos con los que suele confundirse. Un módulo es una unidad interna de un programa: comparte su proceso, su lenguaje y su ciclo de vida, y no existe fuera de la aplicación que lo contiene. Un componente, por su parte, es una unidad de composición con interfaces definidas que puede desarrollarse y distribuirse de forma separada, pero que se integra dentro de una aplicación y se despliega y ejecuta junto con ella \cite{szyperski2002component}. El middleware se diferencia de ambos porque no pertenece a ninguna de las aplicaciones que conecta: es un proceso independiente, con instalación, despliegue y ciclo de vida propios, que presta servicios de comunicación a los extremos a través de una interfaz definida \cite{bernstein1996middleware}.

El middleware resuelve el problema de integración al ser una capa intermediaria entre el cliente de anotación y los servicios de inferencia. Su función no es solo reenviar mensajes: también puede validar solicitudes, consultar el registro de modelos disponibles, elegir el mecanismo de ejecución, gestionar el ciclo de vida de los contenedores, manejar colas de trabajo, registrar errores y traducir respuestas técnicas en estructuras que el cliente pueda entender. Esta mediación reduce el acoplamiento ya que ELAN solo necesita conocer la API del middleware y no los detalles internos de cada servicio.

Esta decisión es aplicable a arquitecturas basadas en servicios y microservicios, donde los componentes independientes se comunican a través de interfaces y pueden evolucionar sin afectar al resto del sistema \cite{dragoni2017microservices}. Una integración local no es necesariamente una plataforma distribuida a gran escala, pero conserva la misma separación de responsabilidades: la herramienta de anotación actúa como cliente, el middleware coordina la comunicación y los servicios de inferencia realizan el procesamiento.

Esta relación se ilustra con el patrón sidecar. En este patrón, un componente adicional acompaña y añade funcionalidad a la aplicación principal sin modificar su lógica interna \cite{burns2016container_patterns}. Con esta idea, el middleware puede ejecutarse localmente, mostrar una interfaz HTTP, gestionar recursos de inferencia y producir resultados de vuelta a la herramienta de anotación. ELAN sigue siendo el principal punto de interacción desde el punto de vista del usuario mientras el middleware absorbe la complejidad de la comunicación, ejecución y control de errores.

El middleware también favorece el mantenimiento: si cambia la forma de ejecutar un backend, la estrategia de almacenamiento del registro o la política de concurrencia, estos cambios pueden introducirse en el orquestador sin modificar la interfaz de ELAN. De manera similar, un servicio de inferencia puede actualizarse siempre que mantenga el contrato de comunicación. Esta independencia es particularmente útil en contextos de investigación donde las herramientas de anotación y los modelos computacionales evolucionan a diferentes ritmos \cite{iso25010}.

\subsubsection{Servicios REST y contratos de comunicación}

La comunicación entre una herramienta de anotación y un middleware puede organizarse mediante una API HTTP basada en REST. REST, propuesto por Fielding, estructura los sistemas distribuidos alrededor de recursos y representaciones transferidas mediante una interfaz uniforme \cite{fielding2000rest}. En una integración local, el uso de HTTP ofrece ventajas prácticas: separa procesos, facilita pruebas con herramientas externas, evita dependencias directas entre Java y Python, y convierte las operaciones del middleware en endpoints documentables.

El uso de HTTP también permite asignar significado a las operaciones del sistema. Una consulta de estado puede realizarse mediante una operación de lectura, una solicitud de inferencia puede representarse como la creación de un trabajo y una consulta de métricas puede exponerse mediante un endpoint específico. Esta organización establece un límite público y verificable entre componentes. En entornos locales, el uso de la dirección de loopback permite que la comunicación ocurra dentro de la misma máquina sin exponer necesariamente el servicio a la red externa; al mismo tiempo, por tratarse de una relación cliente-servidor, el mismo servicio podría publicarse en una interfaz de red y atender a varios clientes sin cambiar el contrato de comunicación.

JSON se utiliza con frecuencia como formato de intercambio porque representa objetos estructurados de forma legible y portable. Una solicitud puede incluir el identificador del modelo, la versión, la ruta o referencia del video, parámetros de configuración y el nivel de anotación destino. Una respuesta puede incluir el identificador del trabajo, el estado de ejecución, los segmentos detectados, sus tiempos de inicio y fin, las etiquetas producidas y metadatos como duración o trazas de diagnóstico. Esta estructura resulta flexible para expresar los datos requeridos por ELAN y suficientemente formal para ser validada por el middleware \cite{pydantic2026models}.

Un contrato de comunicación describe qué campos son aceptados, qué tipos de datos son válidos, qué valores deben respetarse y qué errores pueden ocurrir. Los contratos son necesarios cuando los componentes tienen diferentes lenguajes y responsabilidades: las clases de Java construyen solicitudes desde la herramienta de anotación, los modelos Pydantic validan datos en Python, los servicios de middleware ejecutan la lógica de negocio y los backends devuelven resultados de inferencia. Si alguno de estos componentes interpreta un campo temporal, una etiqueta o un identificador de modelo de manera diferente, la integración se vuelve frágil \cite{fielding2000rest}.

FastAPI y Pydantic permiten que estos contratos se realicen en el middleware. FastAPI define puntos finales y genera documentación OpenAPI a partir de anotaciones de tipo, mientras que Pydantic valida cuerpos de solicitud y respuesta a través de modelos declarativos \cite{fastapi2026features,pydantic2026models}. Esto evita que se inicien procesos costosos cuando la solicitud contiene errores, como un modelo inexistente, un campo requerido faltante o un tipo de dato incompatible. En lugar de fallar en el servicio de inferencia, el middleware puede responder con un error estructurado y comprensible.

La gestión de errores es parte del contrato de comunicación dado que el cliente necesita diferenciar entre un modelo no registrado, un modelo deshabilitado, un archivo de video inválido, un tiempo de espera del backend o una falla interna de Docker. Por ello, un middleware debe utilizar códigos semánticos, estados de trabajo y mensajes claros que permitan diagnosticar el problema sin inspeccionar manualmente cada componente. Esta estructura también facilita la definición de pruebas funcionales asociadas a respuestas esperadas y reproducibles \cite{fastapi2026features,pydantic2026models}.

\subsubsection{Contenedores Docker y aislamiento de entornos de inferencia}

Los servicios de inferencia pueden requerir dependencias distintas de las utilizadas por ELAN o por el middleware. Un backend puede necesitar bibliotecas específicas de procesamiento de video, paquetes de cálculo numérico, archivos de pesos, variables de entorno o configuraciones de hardware particulares. Instalar todas esas dependencias directamente en la máquina del usuario o dentro del proceso principal del middleware incrementaría los conflictos de versión y dificultaría la reproducción del entorno. El uso de contenedores permite aislar esas dependencias en unidades de ejecución independientes \cite{docker2026overview}.

Docker permite empaquetar una aplicación y sus dependencias en una imagen ejecutable. La imagen describe el entorno necesario para iniciar el servicio, mientras que el contenedor corresponde a una instancia en ejecución de esa imagen \cite{docker2026overview}. En una arquitectura de inferencia, el backend puede tratarse como un servicio independiente donde el middleware no instala sus dependencias internas, sino que inicia el contenedor correspondiente y le envía mensajes a través de un contrato HTTP. ELAN, el middleware y el backend están trabajando por separado.

El uso de contenedores ofrece portabilidad y reproducibilidad. La portabilidad permite ejecutar un backend en múltiples máquinas con menos dependencia del sistema operativo anfitrión, siempre que Docker esté disponible. La reproducibilidad permite reconstruir el entorno de ejecución a partir de una definición conocida y minimiza las diferencias entre instalaciones. Estas son propiedades importantes para los investigadores, donde la configuración manual de entornos complejos puede obstaculizar el uso de solución \cite{docker2026overview,gupta2025performance}.

El aislamiento de contenedores también delimita responsabilidades. El backend procesa la entrada y devuelve una respuesta de acuerdo con el contrato; el middleware construye o inicia el contenedor, verifica su disponibilidad, envía la solicitud de inferencia, recibe la respuesta y registra el resultado. La verificación de disponibilidad es una solicitud, generalmente al endpoint GET /health, que confirma que el servicio está activo y listo para recibir inferencias. Si el contenedor no se inicia, no responde a esta verificación o excede el tiempo máximo de inferencia, entonces el middleware puede clasificar el error e informarlo de manera uniforme al cliente \cite{docker2026overview}.

La comunicación entre contenedores generalmente se realiza a través de una red común gestionada por Docker. El middleware puede comunicarse con un backend por el nombre del servicio de la red sin usar la dirección IP o exponer todos los puertos al sistema anfitrión. Para datos persistentes, se pueden usar directorios montados en el anfitrión para que el registro de modelos y los manifiestos de respaldo no dependan de la vida útil de un contenedor. Esto se hace para recuperar el estado después de un reinicio o reinstalaciones parciales \cite{docker2026overview}.

La literatura sobre contenedores en entornos de borde y servicios distribuidos muestra que tales tecnologías reducen la complejidad operativa del despliegue, aunque no eliminan la necesidad de monitorear el consumo de recursos, la latencia y la disponibilidad \cite{gupta2025performance}. Docker debe verse como un mecanismo para el aislamiento y la gestión de servicios, no como una solución completa a todos los problemas de ejecución. Es mejor cuando se combina con construcción controlada, inicio explícito, verificación de disponibilidad, apagado ordenado y observabilidad.

\subsubsection{Orquestación de servicios de inteligencia artificial}

La orquestación de servicios de inteligencia artificial comprende los mecanismos que coordinan una solicitud de inferencia desde su recepción hasta la entrega de resultados estructurados. En este contexto, orquestar no significa entrenar un modelo ni modificar su arquitectura, sino gestionar su uso dentro de un flujo controlado. La orquestación incluye el registro de modelos disponibles, la validación de solicitudes, la selección del backend adecuado, el control de ejecución, el manejo de concurrencia, la recolección de métricas y la provisión de respuestas comprensibles para el cliente \cite{dragoni2017microservices}.

El registro de modelos almacena la información necesaria para determinar qué servicios pueden ejecutarse. Un modelo instalado puede identificarse por nombre, versión, descripción, estado, tipo de ejecutor y configuración de despliegue. Esta información puede mantenerse en un registro persistente y complementarse con manifiestos de respaldo. La separación entre registro y ejecución permite que el cliente solicite un modelo por su identificador, sin conocer la ubicación exacta de sus artefactos, la imagen Docker asociada ni los endpoints internos del backend \cite{docker2026overview}.

El manifiesto del modelo es un contrato de instalación. Describe los datos mínimos que el middleware necesita para registrar y operar un backend: identificadores, versión, configuración de la interfaz, comandos o parámetros de despliegue, y los requisitos que espera del ejecutor. Al validar el manifiesto durante la instalación, el sistema evita asumir paquetes incompletos o inconsistentes. Dicha validación es importante para el usuario final en el sentido de que, para depurar la estructura interna de un paquete ZIP, el middleware debe aceptar o rechazar la instalación con una razón clara \cite{pydantic2026models}.

El ejecutor describe la estrategia de ejecución en términos concretos. Un ejecutor orientado a backends ejecutados en contenedores separa la lógica general de los trabajos de los detalles específicos de Docker. Esta abstracción permite al servicio de trabajos solicitar una inferencia sin conocer todos los pasos de bajo nivel para iniciarla. Si más adelante se añade un segundo tipo de ejecución, la arquitectura puede extenderse siempre que mantenga el contrato de comunicación con el cliente \cite{docker_sdk_python}.

La cola de trabajos es responsable de la concurrencia de la inferencia. Los servicios de inteligencia artificial pueden consumir muchos recursos y no siempre es seguro procesar todas las solicitudes a la vez. La cola FIFO se encarga de los trabajos en orden de llegada, mientras que tiene un límite de concurrencia para evitar saturar la CPU, la memoria o la GPU. Utilizar primitivas de Python como hilos y sincronización permite la ejecución en segundo plano de aplicaciones sin bloquear los puntos finales que necesitan verificar el estado o las métricas \cite{python_threading}.

Los trabajos pueden representarse mediante estados observables. Primero se crea una solicitud, luego se valida, se coloca en cola, se ejecuta contra el backend y finalmente concluye con un resultado exitoso o un fallo clasificado. La diferencia entre los estados de fallo y tiempo de espera es importante porque no todos los errores tienen la misma causa: una falla puede deberse a una respuesta inválida del backend, a un problema de archivo o a un error interno, mientras que un tiempo de espera indica que el servicio no respondió dentro del límite establecido. Separar estos casos mejora el diagnóstico y facilita el diseño de pruebas funcionales precisas.

Las métricas completan la orquestación porque vuelven observable la ejecución. Los contadores de trabajos procesados, trabajos exitosos, fallos, tiempos de espera y latencias ofrecen información sobre el comportamiento del sistema bajo condiciones determinadas. La trazabilidad distribuida, entendida como la capacidad de seguir una solicitud a través de varios componentes, es una práctica estándar para diagnosticar sistemas con múltiples servicios \cite{sigelman2010dapper}. En una arquitectura local compuesta por herramienta de anotación, middleware y backends, cada traza permite identificar si un problema ocurrió en el cliente, en la API, en la cola, en Docker o en el servicio de inferencia.

\subsubsection{Interoperabilidad, trazabilidad y mantenibilidad}

Las decisiones arquitectónicas de una integración no se evalúan únicamente por la existencia de una respuesta correcta en un caso exitoso. Un middleware debe intercambiar información con la herramienta de anotación, reportar fallos de manera comprensible, mantener una estructura modificable y operar con estabilidad en un entorno local. Estas características se relacionan con modelos de calidad de software como ISO/IEC 25010, que categoriza la compatibilidad, mantenibilidad, fiabilidad y seguridad en la evaluación de sistemas de software \cite{iso25010}.

La interoperabilidad es la capacidad de dos o más sistemas para intercambiar información y utilizarla correctamente. En una integración con ELAN, la interoperabilidad se logra cuando una solicitud enviada al middleware produce segmentos que la herramienta puede ubicar en su línea de tiempo. Para ello, los tiempos deben expresarse en unidades compatibles, las etiquetas deben serializarse de forma estable, los identificadores de modelo deben ser válidos y los errores deben seguir una estructura predecible. La interoperabilidad no depende solo de una conexión HTTP; requiere que ambas partes compartan el significado de los datos intercambiados \cite{fielding2000rest}.

La trazabilidad permite reconstruir lo ocurrido durante una solicitud. En una arquitectura con varios componentes, un error observado en ELAN puede originarse en la validación de la solicitud, la selección del modelo, el inicio del contenedor, el backend de inferencia o la transformación de la respuesta. Sin registros y estados observables, todos esos casos se perciben como fallos genéricos. Por ello, estados de trabajo, mensajes de error, códigos semánticos, métricas y trazas ayudan a localizar el punto de fallo. Esta propiedad es esencial para depuración, soporte y mantenimiento posterior.

La mantenibilidad se refiere al grado en que un sistema puede modificarse para corregir errores, adaptarse a nuevos requisitos o agregar capacidades. En un middleware, la mantenibilidad mejora cuando las responsabilidades se separan en capas o módulos: enrutadores para API, esquemas para contratos, servicios para lógica de negocio, ejecutores para ejecución, almacenamiento para estado y componentes específicos para Docker, métricas y cola de trabajos. Esta organización reduce el impacto de los cambios porque cada responsabilidad permanece contenida en una parte del sistema. Si cambia el formato de instalación de modelos, no debería alterarse la interfaz de ELAN; si cambia la política de concurrencia, no debería modificarse el contrato de los segmentos \cite{iso25010}.

La ejecución local agrega una capa de control sobre los datos. Procesar cerca del origen es un principio relacionado con la computación en el borde, orientado a reducir la dependencia de servicios remotos y mantener los datos en un entorno gestionado por el usuario o la institución \cite{shi2016edge}. En el procesamiento de videos de LSEc, esta idea permite que el material audiovisual permanezca en la estación de trabajo o infraestructura controlada, sin enviarse a plataformas externas para la inferencia. Esto mejora la soberanía de los datos y permite operar incluso cuando la conectividad externa no forma parte del flujo de trabajo.

En conjunto, estos conceptos sustentan una arquitectura de integración entre herramientas de anotación y servicios de inferencia. ELAN aporta el entorno de anotación y la semántica temporal del resultado; la LSEc define un dominio de aplicación basado en documentación audiovisual; el reconocimiento automático entrega propuestas revisables; los servicios REST, contratos JSON y contenedores permiten conectar componentes heterogéneos sin acoplar sus dependencias internas; la orquestación organiza registro, ejecución, colas, métricas y errores; y la interoperabilidad, trazabilidad y mantenibilidad ofrecen criterios para evaluar la solución más allá de un caso de prueba exitoso.
\vspace{1\baselineskip}

\section[METODOLOGÍA]{Metodología}

Este capítulo describe la metodología para desarrollar y validar el middleware de orquestación integrado con ELAN. La metodología se organiza en diez apartados: enfoque y tipo de trabajo, técnicas de recolección de información, técnicas de análisis de la información, herramientas y stack tecnológico, proceso de desarrollo (ciclo de vida completo), contratos de comunicación, despliegue y configuración del entorno, estructura de directorios del sistema, guía de replicación y uso ético de herramientas de inteligencia artificial generativa. Para evitar ambigüedades, en este capítulo se utiliza el término \textit{apartado} para referirse a la estructura metodológica general y el término \textit{fase} solo para los incrementos del proceso de desarrollo.

El apartado del proceso de desarrollo recorre el ciclo de vida completo del software: levantamiento de requerimientos, análisis y diseño, implementación en cinco fases (base del middleware, registro de modelos e instalación automática, Docker Runner e inferencia, cola de jobs y métricas, e integración con la interfaz de ELAN), pruebas y validación final, y despliegue. Esta organización permite diferenciar la estructura expositiva del capítulo respecto de los incrementos técnicos ejecutados durante el desarrollo.

El objetivo de esta metodología no es solo enumerar actividades, sino también explicar cómo el problema de integración se convirtió en una arquitectura coherente. El sistema desarrollado conecta una aplicación de escritorio en Java, un middleware en Python, servicios HTTP, contenedores Docker y backends de inferencia; por ello, cada decisión metodológica se justifica según su aporte a la interoperabilidad, la separación de responsabilidades, la validación de contratos y la reducción de dependencias entre componentes. Con el fin de mantener una lectura clara, el capítulo explica las decisiones en un nivel conceptual; el lector que requiera mayor profundidad técnica (esquemas completos, archivos de configuración, procedimientos de compilación y ejemplos íntegros de respuestas) puede dirigirse a los Anexos I, II, III y IV, que se referencian en los puntos correspondientes.

\subsection{Enfoque y tipo de trabajo}

El trabajo es de naturaleza aplicada y tecnológica. No se trata de proponer nuevas ideas en aprendizaje automático, ni de aportar un análisis lingüístico de la LSEc, sino de diseñar, implementar y validar un middleware que resuelve un problema concreto de integración: lograr que ELAN solicite inferencias a servicios de IA externos y reciba segmentos temporales compatibles con su flujo de anotación, sin alterar la arquitectura general de la herramienta y sin obligar al investigador a transcribir manualmente los resultados. Vale aclarar que el término \textit{middleware} se emplea aquí en su acepción de ingeniería de software, es decir, como una capa intermedia e independiente que comunica dos aplicaciones heterogéneas, tal como se fundamentó en el marco teórico; no se trata de un módulo interno de ELAN ni de una librería incorporada a la herramienta.

Esta denominación merece justificarse antes de describir la metodología, tomando como base la distinción entre módulo, componente y middleware presentada en el marco teórico. El software construido en este trabajo no comparte el proceso, el lenguaje ni el ciclo de vida de ELAN, por lo que no puede considerarse un módulo; tampoco se despliega ni se ejecuta junto con la herramienta, condición que caracteriza a un componente. Se trata de un proceso autónomo: se instala y se actualiza por separado, corre en su propio contenedor y atiende solicitudes a través de una interfaz de red, ubicado entre ELAN, escrito en Java, y los servicios de inferencia, escritos en Python, sin que ninguno de los dos extremos necesite conocer su implementación interna. La única pieza del trabajo que sí constituye un componente en sentido estricto es el reconocedor Java añadido a ELAN, porque vive dentro del proceso de la herramienta. Por estas razones, en el resto del documento la solución se nombra como middleware.

Como metodología de desarrollo de software se adoptó el modelo iterativo e incremental, que toma de los marcos ágiles, como Scrum, la idea de ciclos cortos con revisión y adaptación continuas, aunque sin las ceremonias y roles de equipo propios de esos marcos, dado que el TIC fue ejecutado por un solo desarrollador. El ciclo de vida cubrió todas las etapas clásicas del desarrollo: levantamiento de requerimientos, análisis y diseño, implementación, pruebas y despliegue; la diferencia frente a un modelo en cascada es que la implementación y las pruebas se repartieron en incrementos verificables: primero la API base, después el registro de modelos, la ejecución de contenedores, la cola de trabajos, la integración con ELAN y, al final, la validación de escenarios de fallo. Este enfoque resultó adecuado porque al inicio no se conocían todas las restricciones de integración entre ELAN, AVATecH, el cliente Java y el servidor del middleware; cada iteración sirvió para confirmar una decisión técnica, detectar incompatibilidades y ajustar partes ya implementadas sin rediseñar el sistema completo.

Un ejemplo de cambio iterativo es la compatibilidad HTTP entre ELAN y el middleware. Durante las pruebas de integración, el cliente Java necesitaba comunicarse de manera confiable a través de HTTP/1.1, por lo que se ajustó la configuración del servidor ASGI que ejecuta FastAPI para evitar incompatibilidades en las solicitudes. Hallazgos de este tipo muestran que las fases de desarrollo no fueron etapas cerradas, sino ciclos con retroalimentación, en los que cada validación podía obligar a revisar decisiones previas sobre comunicación, configuración o manejo de errores.

\subsection{Técnicas de recolección de información}

Para evaluar los requisitos, las restricciones y las decisiones de implementación, se utilizaron tres métodos de recopilación de información: análisis documental, ingeniería inversa del código fuente de ELAN y revisión de trabajos relacionados. Estos métodos se aplicaron en paralelo. El análisis documental permitió comprender las especificaciones y tecnologías del sistema; la ingeniería inversa reveló los puntos de extensión reales en ELAN; y la revisión de trabajos relacionados permitió analizar cómo se producen los resultados de los sistemas de reconocimiento de señas y cómo deben transformarse para ser útiles en una herramienta de anotación.

\subsubsection{Análisis documental}

La primera técnica consistió en analizar la documentación técnica. Esta revisión no se realizó solo como una actividad descriptiva, sino también como referencia para decidir qué debía implementarse en ELAN, qué debía delegarse al middleware y qué debía resolverse en los backends en contenedores.

El documento clave fue la especificación técnica de AVATecH, que describe cómo ELAN puede integrar componentes de reconocimiento externos \cite{auer2014avatech}. Esta especificación explica cómo un reconocedor obtiene información de configuración, accede a los recursos multimedia del archivo EAF, reporta progreso y devuelve fragmentos candidatos para que ELAN los agregue a los niveles de anotación. De esta revisión se concluyó que la integración debía respetar interfaces Java como \textbf{Recognizer2} o \textbf{IncrementalRecognizer}; así, las decisiones tecnológicas respondieron a las necesidades del sistema y no a preferencias de implementación.

También se revisó la documentación de las tecnologías utilizadas en la capa de orquestación. FastAPI permitió organizar los endpoints, el ciclo de inicialización y la documentación automática de la API \cite{fastapi2026features}; Pydantic habilitó modelos de datos ejecutables para validar solicitudes, respuestas y manifiestos \cite{pydantic2026models}. Docker y su SDK para Python permitieron gestionar imágenes, contenedores, redes y montajes desde el middleware \cite{docker2026overview,docker_sdk_python}; y el estilo REST orientó la definición de operaciones HTTP uniformes para consultas, instalaciones, cambios de estado y consulta de resultados \cite{fielding2000rest}. Esta revisión reforzó el vínculo entre la selección tecnológica y los requisitos del sistema.

\subsubsection{Ingeniería inversa del código fuente de ELAN}

El segundo enfoque fue la ingeniería inversa del código fuente de ELAN. La versión 7.1 se utilizó como referencia para el análisis y las pruebas de integración. La revisión no buscó rehacer ninguna parte de la herramienta, sino identificar los puntos donde ELAN permite incluir reconocedores externos en el marco AVATecH. Por esa razón, se analizó el paquete \textbf{mpi.eudico.client.annotator.recognizer}, donde se ubican las interfaces, clases de soporte y mecanismos de comunicación asociados con los reconocedores \cite{auer2010elan,auer2014avatech}.

El primer elemento identificado fue el contrato que un reconocedor externo debe cumplir. En ELAN, este contrato está representado principalmente por las interfaces \textbf{Recognizer2} e \textbf{IncrementalRecognizer}. Estas interfaces permiten recibir parámetros, iniciar el procesamiento mediante \textbf{start()}, abortar la ejecución y comunicarse con ELAN durante el proceso. La comunicación se realiza mediante \textbf{RecognizerHost}, interfaz que permite informar progreso, reportar errores y enviar segmentaciones al entorno de anotación \cite{auer2014avatech}. Esto dejó claro que el cliente Java debía trabajar dentro del flujo de ELAN y no como una herramienta externa desconectada.

El segundo elemento fue el mecanismo de entrega de parámetros. La revisión mostró que ELAN no envía una estructura única y cerrada, sino pares clave-valor a través de \textbf{setParameterValue()}; por ello se definieron parámetros explícitos para la ruta del video, el modelo seleccionado, el nivel de anotación destino y las opciones de inferencia. De esta manera, el panel de configuración no depende de valores fijos para enviar información al middleware.

El tercer elemento fue el retorno de resultados. ELAN no espera una lista genérica de predicciones, sino un objeto \textbf{Segmentation} asociado a un archivo multimedia y a un nivel de anotación, compuesto por objetos \textbf{Segment} con tiempo de inicio, tiempo de fin y etiqueta. Este formato determinó el contrato de respuesta del middleware: cada resultado del backend debía convertirse en un segmento temporal expresado en milisegundos y asociarse al nivel de anotación indicado por el usuario, pues de lo contrario ELAN no puede ubicar la anotación en la línea de tiempo.

También se revisó la construcción del panel gráfico del reconocedor. En ELAN, los paneles de configuración se basan en \textbf{AbstractRecognizerPanel}; su función es organizar controles de entrada antes de iniciar el procesamiento, no ejecutar la inferencia. Por ello se definió \textbf{AIOrchestrationRecognizerPanel} para exponer selección de modelo, validación del video, nivel de anotación destino y parámetros de inferencia. Con esta separación, el panel recopila configuración y el reconocedor ejecuta la comunicación con el middleware y transforma la respuesta en objetos compatibles con ELAN.

Finalmente, se revisó la comunicación HTTP disponible desde Java. Esta revisión orientó la implementación hacia \textbf{java.net.http.HttpClient}, porque permite construir un cliente dentro del mismo proceso de ELAN sin introducir procesos auxiliares. A partir de esta decisión, las clases de integración se ubicaron en el paquete \textbf{mpi.eudico.client.annotator.recognizer.ai}. La ubicación no responde solo a organización del código: permite reutilizar inicialización, configuración, cancelación, reporte de progreso, manejo de errores y entrega de segmentaciones mediante \textbf{RecognizerHost}, manteniendo compatibilidad con el mecanismo de extensión de AVATecH \cite{auer2014avatech}.

\subsubsection{Revisión de trabajos relacionados y modelos existentes}

La tercera técnica fue la revisión de trabajos previos sobre el reconocimiento automático de LSEc y la integración de detectores en ELAN. Esta revisión ayudó a distinguir entre el resultado de un modelo y una anotación que un investigador puede usar: un modelo puede producir una etiqueta, una predicción o una lista de probabilidades, pero ELAN necesita segmentos con inicio, fin y valor textual para agregarlos al nivel de anotación. El trabajo de Guevara sobre el intérprete de frases LSEc muestra que los sistemas de reconocimiento aprenden información visual y generan una clasificación del signo o frase observada \cite{guevara2022lsec_interprete}. Este antecedente confirmó que la salida de un reconocedor requiere una transformación adicional antes de integrarse en la línea de tiempo de ELAN.

Desde una perspectiva técnica, la revisión orientó el contrato de respuesta del middleware. Dado que los backends procesan video, la respuesta debía conservar la información temporal del archivo y no limitarse a un glosario final. Así, se definió \textbf{media\_info} para los metadatos del video y \textbf{predictions} para transportar alternativas de clasificación con sus probabilidades cuando el backend las proporciona. La revisión final queda en manos del investigador: el sistema propone segmentos y etiquetas, pero no convierte las predicciones en anotaciones definitivas sin validación humana.

La revisión también ayudó a diferenciar tres ideas que suelen confundirse: reconocimiento, segmentación y anotación. El reconocimiento produce una etiqueta candidata; la segmentación determina el intervalo de tiempo en el que ocurre el evento; y la anotación es la inserción del resultado en ELAN, siempre sujeta a revisión. Esta distinción justificó que el middleware no entregara solo una predicción textual, sino segmentos con tiempos, etiquetas y alternativas cuando existan. Además, mantiene el contrato independiente de la arquitectura de IA y de la implementación interna del backend.

Finalmente, se revisaron antecedentes de integración entre ELAN y detectores automáticos. Auer et al. describen extensiones hacia la anotación semiautomática de audio y video \cite{auer2010elan}. Este antecedente confirmó que la conexión entre ELAN y procesos externos ya existía, pero también evidenció que el presente trabajo necesitaba una capa adicional: no bastaba con invocar al detector; también había que gestionar los modelos instalados, los contenedores, los contratos HTTP, los errores y los resultados temporales. De ahí que el middleware se plantee como un orquestador entre ELAN y los backends de inferencia.

\subsection{Técnicas de análisis de la información}

La información recopilada se analizó mediante tres técnicas: descomposición funcional, diseño basado en contratos y validación incremental. Estas técnicas convirtieron la información dispersa en decisiones de diseño verificables. La descomposición determinó qué responsabilidades debía asumir cada componente; el diseño basado en contratos definió cómo se comunicarían componentes escritos en tecnologías distintas; y la validación incremental estableció cómo verificar que cada incremento siguiera funcionando dentro del flujo completo.

\subsubsection{Descomposición funcional}

La descomposición funcional consistió en dividir la integración entre ELAN, el middleware y los backends en responsabilidades pequeñas y verificables. Esta decisión se relaciona con la mantenibilidad del software, ya que un sistema con componentes cohesionados y bajo acoplamiento es más fácil de modificar, probar y extender \cite{iso25010}. En este trabajo, la descomposición se refinó durante el desarrollo; cada fase permitió identificar nuevos puntos de integración, errores de comunicación y necesidades de validación.

En este sentido, se definieron servicios internos del middleware. El \textbf{ModelRegistryService} es responsable de la instalación, validación, registro y consulta de modelos utilizando el registry.json. El \textbf{DockerLifecycleService} es responsable de la creación de imágenes y el inicio inicial durante la instalación, mientras que el \textbf{DockerService} es responsable de ejecutar los contenedores. Esta separación respondió a que instalar un modelo y ejecutar una inferencia son procesos distintos: la instalación crea los artefactos y deja el backend disponible por primera vez, mientras que la inferencia localiza el contenedor activo, envía la solicitud y gestiona el resultado. En ambos casos se emplean contenedores para aislar las aplicaciones y sus dependencias \cite{docker2026overview}.

\textbf{JobService} se definió como coordinador de inferencias: recibe la solicitud, consulta el modelo registrado, prepara la ejecución, invoca el backend y devuelve una respuesta compatible con ELAN. \textbf{JobQueue} gestiona la concurrencia mediante una cola FIFO con un límite de inferencias simultáneas, de modo que los recursos de CPU, memoria y GPU no se saturen. Por su parte, \textbf{MetricsService} mantiene contadores de observabilidad: trabajos completados, fallidos, en cola o con tiempo de espera agotado. Con esta separación es posible evaluar el estado del sistema sin mezclar métricas, ejecución y persistencia.

\subsubsection{Diseño dirigido por contratos}

La segunda técnica se basó en contratos. En este trabajo, un contrato es la definición explícita de los datos que un componente acepta, valida y produce antes de comunicarse con otro componente. Esta decisión fue necesaria porque ELAN está implementado en Java, el middleware en Python y los backends se exponen como servicios HTTP. En una arquitectura de este tipo, asumir estructuras implícitas aumenta el riesgo de errores que son difíciles de diagnosticar. Por lo tanto, se definieron contratos para solicitudes HTTP, respuestas de inferencia y paquetes de instalación usando \textbf{manifest.json}, según el principio REST de interfaces uniformes y mensajes bien definidos \cite{fielding2000rest}.

En el middleware, los contratos se implementaron utilizando modelos Pydantic junto con FastAPI. Esta combinación permitió la validación de cuerpos de solicitud antes de ejecutar la lógica de negocio. Si el cliente Java envía un campo faltante, un tipo incorrecto o una estructura inválida, entonces el middleware puede enviar un error estructurado en lugar de pasar por la falla interna. Por tanto, los contratos no quedaron solo como documentación: se convirtieron en una parte ejecutable del sistema \cite{pydantic2026models,fastapi2026features}.

El diseño por contratos también se aplicó al paquete de cada modelo. El archivo \textbf{manifest.json} describe identificador, versión, artefactos, forma de construcción del backend Docker y endpoints de inferencia. Validarlo de forma temprana permite detectar paquetes incompletos antes de construir o iniciar un contenedor, lo que resulta preferible a descubrir el problema durante una inferencia, cuando el modelo ya figura como registrado y el usuario espera que esté disponible.

\subsubsection{Validación incremental mediante pruebas funcionales}

La tercera técnica fue la validación mediante pruebas funcionales. Cada fase produjo un incremento verificable: API base, registro de modelos, gestión de contenedores, cola de inferencia, integración con ELAN y recuperación ante fallos. Se eligió este tipo de validación porque el sistema puede fallar por causas muy distintas entre sí: errores de contrato HTTP, rutas de video inválidas, contenedores detenidos, tiempos de espera del backend o incompatibilidades entre la respuesta del middleware y el formato esperado por ELAN.

Las pruebas se realizaron sobre las interfaces públicas. En el middleware se enviaron solicitudes HTTP válidas e inválidas y las respuestas se compararon con los contratos definidos; en ELAN se comprobó que los segmentos generados aparecieran en el nivel de anotación configurado, con tiempos de inicio y fin consistentes con la respuesta del backend. De este modo se validó el flujo completo: cliente Java, API REST, servicios internos, backend en contenedor y retorno de segmentaciones a ELAN. Además, cualquier escenario puede repetirse siempre que se conserven la solicitud, el paquete y las condiciones de ejecución.

En conjunto, las tres técnicas llevaron de información técnica dispersa a una arquitectura verificable: la descomposición funcional organizó las responsabilidades, el diseño por contratos definió interfaces comprobables y la validación incremental confirmó que esas decisiones funcionaban en el sistema completo.

\subsection{Herramientas y stack tecnológico}

La Tabla~\ref{tab:stack_tecnologico} resume las herramientas utilizadas y las versiones establecidas como base durante el desarrollo. Estas versiones no son la única combinación posible, pero sí el entorno probado para evitar incompatibilidades entre el cliente Java de ELAN, el middleware, Docker y el backend de inferencia. La selección tecnológica tuvo en cuenta los siguientes puntos: comunicación explícita entre componentes, validación de contratos y aislamiento de dependencias a través de contenedores.

\begin{table}[H]
\centering
\caption{Stack tecnológico del sistema desarrollado}
\label{tab:stack_tecnologico}
\footnotesize
\renewcommand{\arraystretch}{1.18}
\setlength{\tabcolsep}{5pt}

\begin{tabular}{
    >{\raggedright\arraybackslash}p{3.1cm}
    >{\raggedright\arraybackslash}p{3.4cm}
    >{\raggedright\arraybackslash}p{6.2cm}
}
\toprule
\textbf{Componente} & \textbf{Tecnología y versión} & \textbf{Razón principal de selección} \\
\midrule

Middleware &
Python 3.11 &
Compatibilidad con bibliotecas de inteligencia artificial y mejoras de rendimiento del intérprete. \\

Framework API &
FastAPI 0.115 &
Definición de endpoints REST, validación integrada y documentación OpenAPI. \\

Validación de datos &
Pydantic 2.x &
Modelos de datos ejecutables para validar solicitudes, respuestas y manifiestos. \\

Servidor ASGI &
Uvicorn 0.32 &
Ejecución de la aplicación FastAPI sobre el estándar ASGI. \\

Contenerización &
Docker 25.x y Docker Compose &
Aislamiento de modelos, dependencias y redes de ejecución. \\

SDK Docker &
Docker SDK for Python 7.x &
Control programático de imágenes y contenedores desde el middleware. \\

Persistencia local &
JSON sobre bind mounts &
Registro simple y auditable de modelos en un entorno local de usuario único. \\

Cliente ELAN &
Java 11 &
Disponibilidad del paquete \texttt{java.net.http} para comunicación HTTP desde ELAN. \\

Cliente--middleware &
HTTP/1.1 + JSON &
Interfaz simple entre el cliente Java y el middleware REST. \\

Middleware--backend &
HTTP interno en red Docker &
Comunicación aislada entre contenedores sin exponer todos los servicios al anfitrión. \\

\bottomrule
\end{tabular}
\end{table}

Python 3.11 fue seleccionado para implementar el middleware porque permite construir servicios de integración compatibles con el ecosistema utilizado por muchos backends de inferencia. En este proyecto, Python no se usó para entrenar modelos, sino para implementar la API, validar datos, coordinar jobs, invocar servicios externos e interpretar respuestas. La versión 3.11 se fijó como línea base por sus mejoras de rendimiento y por ofrecer una plataforma estable para bibliotecas modernas sin depender de versiones experimentales \cite{python311docs}.

FastAPI fue elegido para implementar la API REST porque permite declarar endpoints HTTP usando tipos de Python y modelos Pydantic. Esto fue importante porque el cliente Java envía solicitudes JSON y el middleware debe rechazar de forma controlada cualquier estructura incompleta o inválida. Además, FastAPI genera documentación OpenAPI y soporta carga de archivos mediante \textbf{UploadFile}, lo que facilitó la instalación de paquetes ZIP de modelos \cite{fastapi2026features}.

Se eligió Pydantic 2.x para representar solicitudes, respuestas y manifiestos como modelos de datos. Con esta decisión, errores de tipo, campos faltantes o estructuras inválidas se detectan antes de ejecutar la lógica de negocio \cite{pydantic2026models}. Por ello, Pydantic se aplicó tanto a endpoints del middleware como al archivo \textbf{manifest.json}, evitando que errores de empaquetado se descubran recién durante una inferencia.

Docker y Docker Compose se seleccionaron para aislar los backends de inferencia junto con sus dependencias, ya que mezclarlas con ELAN o con el entorno base del middleware incrementaría el riesgo de incompatibilidades. Docker empaqueta servicios en contenedores y los conecta mediante redes administradas \cite{docker2026overview}. Sobre esa base, el Docker SDK for Python se incorporó para construir imágenes, crear contenedores, iniciar servicios, inspeccionar estados y detener backends desde el propio middleware, sin depender de comandos manuales de consola \cite{docker_sdk_python}.

Se eligió Java 11 para el cliente integrado en ELAN porque incorpora \textbf{java.net.http}, paquete que proporciona \textbf{HttpClient}, \textbf{HttpRequest} y \textbf{HttpResponse} como API estándar para comunicación HTTP \cite{java11_httpclient}. Esta decisión evitó agregar librerías externas al módulo de integración. Además, \textbf{HttpClient} permitió configurar HTTP/1.1 para mantener compatibilidad con el servidor del middleware. Para instalar modelos desde ELAN se construyó manualmente el cuerpo \textbf{multipart/form-data}, manteniendo un cliente Java ligero y controlado.

La persistencia local se resolvió mediante archivos JSON almacenados en bind mounts. Esta decisión responde a un escenario local de usuario único, donde un investigador ejecuta ELAN, el middleware y los backends en su estación de trabajo. En ese contexto, una base de datos relacional exigiría una configuración innecesaria solo para registrar modelos, verificar su estado y recuperar manifiestos. El archivo \textbf{registry.json} funciona como un registro auditable y, gracias a los bind mounts, los datos permanecen en el sistema de archivos del anfitrión aunque los contenedores se recreen \cite{docker2026overview}.

Uvicorn se eligió como servidor ASGI para ejecutar la aplicación FastAPI. Su papel no es realizar inferencias, sino mantener accesible la API del middleware: mientras una inferencia prolongada ocupa un hilo de trabajo, los endpoints de salud, gestión y métricas siguen respondiendo. Este comportamiento complementa a \textbf{JobQueue}, que limita las inferencias simultáneas sin bloquear el resto del servicio. La versión 0.32 se fijó por compatibilidad con FastAPI y con la configuración utilizada durante el desarrollo \cite{uvicorn_docs}.
\vspace{1\baselineskip}

\subsection{Proceso de desarrollo: ciclo de vida del software}

El desarrollo siguió el ciclo de vida completo del software bajo un modelo iterativo e incremental. La Figura~\ref{fig:ciclo_desarrollo} resume ese ciclo: el levantamiento de requerimientos definió qué debía hacer el sistema; el análisis y diseño convirtió esos requerimientos en una arquitectura; la implementación se dividió en cinco fases que entregaron incrementos operativos; la etapa de pruebas validó el sistema completo (fase 6); y el despliegue dejó la solución instalada y documentada. Las flechas de retorno indican que el proceso no fue lineal: cada validación podía obligar a revisar el diseño o ajustar incrementos anteriores, como ocurrió con la compatibilidad HTTP entre el cliente Java y el servidor del middleware.

\begin{figure}[H]
\centering
\resizebox{0.98\textwidth}{!}{%
\begin{tikzpicture}[
    etapa/.style={rectangle, rounded corners=2pt, draw=black!70, fill=gray!12,
                  text width=2.5cm, minimum height=1.5cm, align=center, font=\footnotesize},
    flecha/.style={-{Stealth[length=2.6mm]}, thick},
    retro/.style={-{Stealth[length=2.2mm]}, thick, dashed, black!60},
    node distance=0.55cm
]
\node[etapa] (req)  {Levantamiento de requerimientos};
\node[etapa, right=of req]  (dis)  {Análisis y diseño de la arquitectura};
\node[etapa, right=of dis]  (imp)  {Implementación incremental\\ (fases 1 a 5)};
\node[etapa, right=of imp]  (val)  {Pruebas y validación final\\ (fase 6)};
\node[etapa, right=of val]  (desp) {Despliegue y documentación};

\draw[flecha] (req) -- (dis);
\draw[flecha] (dis) -- (imp);
\draw[flecha] (imp) -- (val);
\draw[flecha] (val) -- (desp);

\draw[retro] (imp.north) -- ++(0,0.6) -| (dis.north)
      node[pos=0.22, above, font=\scriptsize] {ajustes de diseño};
\draw[retro] (val.south) -- ++(0,-0.6) -| (imp.south)
      node[pos=0.22, below, font=\scriptsize] {retroalimentación de cada iteración};
\end{tikzpicture}%
}
\caption{Ciclo de vida de desarrollo seguido en el trabajo}
\label{fig:ciclo_desarrollo}
\end{figure}

La lógica del avance fue incremental: primero se consiguió que ELAN, el middleware y los backends de inferencia se comunicaran; luego se añadieron los componentes que hacen esa comunicación estable, observable y recuperable ante fallos. El criterio de progreso no era que el código compilara, sino que cada parte cumpliera su porción del contrato dentro de la arquitectura general.

\subsubsection{Levantamiento de requerimientos}

Los requerimientos del sistema se derivaron de las técnicas de recolección de información descritas antes: la especificación de AVATecH y la ingeniería inversa de ELAN definieron las condiciones que la integración debía respetar, mientras que la revisión de trabajos relacionados aclaró qué forma debían tener los resultados para ser útiles al investigador. A partir de ese material se establecieron los requerimientos funcionales y no funcionales de la Tabla~\ref{tab:requerimientos}, que guiaron el diseño de la arquitectura y sirvieron después como criterios de verificación en la validación final.

\begin{table}[H]
\centering
\caption{Requerimientos funcionales y no funcionales del sistema}
\label{tab:requerimientos}
\footnotesize
\renewcommand{\arraystretch}{1.18}
\setlength{\tabcolsep}{5pt}

\begin{tabular}{
    >{\centering\arraybackslash}p{1.4cm}
    >{\raggedright\arraybackslash}p{11.6cm}
}
\toprule
\textbf{Código} & \textbf{Descripción} \\
\midrule
RF-01 & Instalar modelos de IA desde un paquete ZIP que incluya su manifiesto de instalación. \\
RF-02 & Listar los modelos instalados y consultar el detalle de cada uno. \\
RF-03 & Activar y desactivar modelos previamente instalados. \\
RF-04 & Ejecutar inferencias sobre un video y devolver segmentos temporales con etiquetas de glosa. \\
RF-05 & Consultar el resultado de un trabajo de inferencia y las métricas acumuladas del servicio. \\
RF-06 & Gestionar el ciclo de vida de los backends en contenedores: construcción, arranque, verificación de disponibilidad y reinicio. \\
RF-07 & Permitir la configuración y ejecución del reconocedor desde la interfaz gráfica de ELAN, incluida la gestión de modelos. \\
RF-08 & Informar los errores con códigos y mensajes comprensibles para el usuario y para otros clientes HTTP. \\
\midrule
RNF-01 & Ejecución completamente local, sin enviar videos ni anotaciones a servicios externos. \\
RNF-02 & Aislamiento de las dependencias de cada modelo respecto de ELAN y del propio middleware. \\
RNF-03 & Control configurable de la cantidad de inferencias simultáneas para no saturar el equipo. \\
RNF-04 & Persistencia del registro de modelos y capacidad de recuperarlo tras reinicios o pérdidas parciales. \\
RNF-05 & Compatibilidad con ELAN 7.1 y su mecanismo de extensión AVATecH, sin alterar la arquitectura de la herramienta. \\
\bottomrule
\end{tabular}
\end{table}

\subsubsection{Análisis y diseño}

Antes de implementar el sistema se revisaron la especificación de AVATecH, el código fuente de ELAN y las restricciones operativas de los backends de IA. Esta revisión ayudó a delimitar qué parte del flujo debía quedar dentro de ELAN y qué parte correspondía a un servicio externo. AVATecH establece que los reconocedores externos deben implementarse a través de interfaces esperadas por ELAN y devolver resultados convertibles en anotaciones temporales \cite{auer2014avatech}; por lo tanto, la solución debía cubrir la configuración, el reporte de progreso, la cancelación y el retorno de segmentos.

Se analizaron cuatro alternativas de integración antes de definir la arquitectura. La primera fue ejecutar procesos externos desde ELAN mediante línea de comandos o archivos intermedios. Aunque mantenía el modelo fuera de ELAN, complicaba el reporte de progreso, la cancelación y la conversión de resultados a estructuras compatibles con el anotador. La segunda fue empaquetar toda la lógica como una librería JAR incorporada a ELAN, de manera que la herramienta resolviera la inferencia dentro de su propio proceso. Esta opción se descartó por una razón de fondo: los modelos de reconocimiento no son código Java, sino pipelines en Python con dependencias nativas de visión artificial y aprendizaje profundo, pesos de cientos de megabytes y requisitos de hardware propios; nada de eso puede empaquetarse en un JAR. Además, una librería comparte el ciclo de vida de la aplicación que la carga: cada actualización de un modelo obligaría a recompilar y redistribuir ELAN, y un fallo o fuga de memoria del modelo comprometería directamente la estabilidad de la herramienta de anotación. La tercera alternativa fue tender un puente entre Java y las bibliotecas nativas mediante JNI, pero introducía una frontera compleja entre Java, Python, bibliotecas nativas y dependencias del modelo \cite{oracle_jni}, y reducía la portabilidad, porque el usuario tendría que mantener un entorno compatible con todas esas dependencias.

La alternativa adoptada fue implementar en ELAN un reconocedor Java que actúe como cliente HTTP de un servicio local. Con ello el sistema queda organizado como cliente-servidor: el middleware es un software independiente que se instala, arranca y actualiza por separado, y ELAN se limita a consumir su API. La elección de HTTP merece una justificación explícita porque el escenario previsto es local y podría parecer innecesario usar un protocolo de red. En la práctica, el costo es despreciable (la comunicación viaja por la interfaz de loopback de la misma máquina, sin salir a la red) y las ventajas son concretas: los procesos quedan aislados, de modo que un fallo del modelo no afecta a ELAN; la frontera entre Java y Python es explícita y verificable, lo que permitió probar el middleware con la herramienta Postman sin depender de ELAN; los modelos se actualizan sin recompilar la herramienta; y la misma interfaz sirve sin cambios si en el futuro el servidor se ubica en otra máquina de la red. La comunicación HTTP/1.1 con JSON se eligió por ser explícita, verificable y desacoplada, de acuerdo con el estilo REST \cite{fielding2000rest}. Para el cliente Java se usó \textbf{java.net.http.HttpClient}, disponible desde Java 11, evitando añadir librerías externas al código de ELAN \cite{java11_httpclient}.

A partir de esta decisión se definió una arquitectura de tres capas: ELAN como cliente de anotación, el middleware como orquestador local y los backends como servicios contenerizados. Docker permitió empaquetar cada backend con sus dependencias y ejecutarlo de forma aislada del sistema anfitrión \cite{docker2026overview}. Así, el investigador puede seguir trabajando en ELAN, mientras el middleware instala modelos, verifica disponibilidad, crea contenedores y envía solicitudes de inferencia.

\begin{figure}[htbp]
    \centering
    \includegraphics[width=1\textwidth]{Diagramas Finales/General Metodología.drawio.png}
    \caption{Arquitectura general del sistema de orquestación}
    \label{fig:arquitectura_general}
\end{figure}

La Figura~\ref{fig:arquitectura_general} resume esta separación de responsabilidades. ELAN no ejecuta directamente el pipeline de IA; envía una solicitud al middleware. El middleware valida la petición, localiza el modelo, coordina Docker y devuelve segmentos temporales que ELAN puede insertar en el nivel de anotación configurado.

En esta etapa también se definió el contrato de instalación mediante \textbf{manifest.json}. Este archivo acompaña a cada paquete de modelo y declara identidad, versión, tarea, artefactos, runtime, configuración del contenedor, endpoints del backend y opciones visibles en ELAN. Definir el esquema antes de implementar la instalación permitió validar los paquetes de forma temprana y devolver errores estructurados, aprovechando las capacidades de Pydantic y FastAPI \cite{pydantic2026models,fastapi2026features}.

El diseño también estableció la estructura interna de clases del middleware, que se mantuvo estable durante todas las fases de implementación. La Figura~\ref{fig:diagrama_clases} presenta el diagrama de clases de los servicios principales con sus operaciones públicas y sus relaciones de uso. La capa de API recibe las solicitudes y las delega en dos coordinadores: \textbf{ModelRegistryService}, a cargo de la instalación y el registro de modelos, y \textbf{JobService}, a cargo de las inferencias. Ambos se apoyan en servicios especializados: la cola de trabajos, las métricas, la selección del mecanismo de ejecución y los servicios que hablan con Docker. La clase abstracta \textbf{BaseRunner} define el contrato de ejecución y \textbf{DockerRunner} lo implementa para backends en contenedores, de modo que otros mecanismos de ejecución podrían agregarse sin tocar el resto del sistema. Las clases Java integradas en ELAN se describen en el apartado de estructura de directorios y su inventario completo consta en el Anexo~IV.

% El diagrama de clases se elabora en draw.io y se exporta como PNG
% (archivo fuente: Diagramas/diagrama_clases_middleware.drawio).
\begin{figure}[H]
    \centering
    \includegraphics[width=1\textwidth]{Diagramas Finales/Diagrama de clases del middleware.drawio.png}
    \caption{Diagrama de clases de los servicios principales del middleware}
    \label{fig:diagrama_clases}
\end{figure}

\subsubsection{Fase 1 — Base del middleware}

El primer paso fue construir una base del middleware simple, ejecutable y verificable. Antes de abordar la gestión de modelos o la inferencia, se implementó un servicio HTTP local capaz de iniciarse en Docker, cargar su configuración, exponer un endpoint de salud y generar registros de diagnóstico. Comenzar por esta base redujo el riesgo inicial, pues comprobó que Python, FastAPI, Uvicorn, Docker y la red de despliegue funcionaban en conjunto.

Desde el principio, las responsabilidades internas se separaron. La capa \textbf{api/} contiene los routers HTTP; \textbf{services/} concentra la lógica de negocio; \textbf{runners/} agrupa los mecanismos de ejecución; \textbf{schemas/} define los contratos Pydantic; y \textbf{core/} centraliza configuración y logging. Esta organización evita que los handlers acumulen lógica de orquestación y facilita la modificación de una parte del sistema sin afectar a todo el middleware.

La configuración se centralizó en \textbf{app/core/config.py} mediante una clase \textbf{Settings} basada en Pydantic. Las rutas, los límites de concurrencia, los tiempos de espera, el nivel de registro y los parámetros del contenedor se declaran allí. El registro estructurado se implementó en \textbf{app/core/logging\_config.py}, con mensajes que incluyen severidad, marca de tiempo, módulo y descripción. La verbosidad se controla con \textbf{MIDDLEWARE\_LOG\_LEVEL}, siguiendo la práctica de configurar el comportamiento desde el entorno \cite{python_logging}.

El manejo global de errores se definió en \textbf{main.py} a través de manejadores de excepciones de FastAPI para errores de registro, Docker y selección de runners. Todos los errores se convierten en una respuesta uniforme con código semántico y detalle legible.

\begin{center}
\begin{minipage}{0.85\textwidth}
\hrule
\vspace{0.12cm}
\small
\begin{Verbatim}[baselinestretch=0.82]
{"error_code": "CODIGO_DESCRIPTIVO", "detail": "Mensaje legible."}
\end{Verbatim}
\vspace{0.05cm}
\hrule

\captionof{codigo}{Formato uniforme de respuesta de error del middleware}
\label{lst:error_response}
\end{minipage}
\end{center}

Este formato fue importante para la integración con ELAN, porque el cliente Java puede leer siempre el campo \textbf{detail}, sin importar si el fallo proviene del registro de modelos, Docker o el backend de inferencia.

La configuración del contenedor del middleware se definió con un \textbf{Dockerfile} basado en \textbf{python:3.11-slim} y un \textbf{docker-compose.yml} con servicio, bind mounts, variables de entorno y red compartida. La red \textbf{elan-ai-shared} se declaró con nombre fijo para que los contenedores de modelos instalados posteriormente se unan al mismo espacio de comunicación \cite{docker2026overview}.

Durante el arranque, antes de aceptar solicitudes, el middleware configura los logs y escanea \textbf{data/bootstrap\_manifests/} para recuperar modelos instalados en sesiones previas. El criterio de verificación fue levantar el servicio con \textbf{docker compose up --build -d} y confirmar la respuesta de \textbf{GET \symbol{47}health}:

\begin{center}
\begin{minipage}{0.95\textwidth}
\hrule
\vspace{0.12cm}
\small
\begin{Verbatim}[baselinestretch=0.82]
{"status": "ok", "service": "elan-ai-orchestrator", "version": "0.1.0"}
\end{Verbatim}
\vspace{0.05cm}
\hrule

\captionof{codigo}{Respuesta esperada del endpoint de salud del middleware}
\label{lst:health_response}
\end{minipage}
\end{center}

Además de la respuesta HTTP, se verificó que Uvicorn iniciara sin errores y que el endpoint de salud siguiera respondiendo mientras una inferencia larga estaba en ejecución. Con ello se confirmó que la base del middleware era suficiente para el caso de uso local previsto \cite{uvicorn_docs}.

\subsubsection{Fase 2 — Registro de modelos e instalación automática}

La segunda fase abordó la gestión de modelos. El middleware debía recibir un paquete ZIP, validar su contenido, construir la imagen Docker del backend, crear el contenedor, verificar su disponibilidad y registrar el modelo instalado. Dado que esta fase combina sistema de archivos, validación, Docker y persistencia, se diseñó como un flujo secuencial con puntos de control y recuperación.

El proceso comienza con la recepción del ZIP y la localización del \textbf{manifest.json}. El middleware valida el manifiesto contra el esquema \textbf{ModelManifest}, comprueba que los artefactos declarados existan dentro del paquete, rechaza rutas inseguras y verifica que la combinación de \textbf{model\_id} y \textbf{version} no esté ya instalada. Solo entonces extrae el paquete, construye la imagen Docker, crea el contenedor y ejecuta la comprobación de disponibilidad. Si la construcción, el arranque, esa comprobación o la persistencia fallan, se ejecuta un rollback que elimina los archivos extraídos y la imagen construida, de modo que el sistema regresa a su estado anterior.

La validación previa a la extracción actúa como barrera de seguridad y consistencia: el identificador del modelo solo admite caracteres seguros, todos los artefactos declarados deben existir dentro del ZIP, las rutas no pueden contener \textbf{../} ni ser absolutas, y la tarea debe estar entre las soportadas. Estas reglas evitan instalar paquetes incompletos o capaces de escribir fuera del directorio previsto, y se apoyan en la validación temprana descrita en el diseño por contratos.

El \textbf{manifest.json} es el contrato entre quien empaqueta un modelo y el middleware. La Tabla~\ref{tab:manifest_campos} resume los campos que el orquestador necesita para instalar, ejecutar y presentar un modelo sin codificar esos datos dentro del middleware. El esquema completo del archivo, con los valores admitidos en cada objeto y el manifest real del modelo de referencia, se recoge en el Anexo~I.

\begin{table}[H]
\centering
\caption{Campos principales del \texttt{manifest.json}}
\label{tab:manifest_campos}
\small
\renewcommand{\arraystretch}{1.15}
\begin{tabular}{@{}>{\raggedright\arraybackslash}p{3.1cm}>{\raggedright\arraybackslash}p{2.1cm}>{\raggedright\arraybackslash}p{7.2cm}@{}}
\toprule
\textbf{Campo} & \textbf{Tipo} & \textbf{Descripción} \\
\midrule
\texttt{model\_id} & string & Identificador único del modelo. Solo admite caracteres \texttt{[A-Za-z0-9\_.-]}. \\
\texttt{version} & string & Versión semántica del modelo, por ejemplo \texttt{1.0.0}. \\
\texttt{task} & string & Tarea del modelo: \texttt{video\_segmentation} o \newline\texttt{video\_segmentation\_and}\newline\texttt{\_gloss\_classification}. \\
\texttt{runtime.mode} & string & Modo de ejecución. En esta implementación, el único valor soportado es \texttt{docker}. \\
\texttt{artifacts} & object & Rutas relativas dentro del ZIP para cada artefacto declarado; todos los archivos deben existir. \\
\texttt{backend\_config} & object & Nombre de imagen, puerto, ruta de comprobación de disponibilidad y timeout de arranque usados durante la instalación. \\
\texttt{container} & object & Puerto, ruta de comprobación de disponibilidad y ruta de inferencia usados durante cada ejecución. \\
\texttt{ui} & object & Configuración de la interfaz de ELAN: nivel de anotación destino, etiqueta y modo de etiquetado. \\
\bottomrule
\end{tabular}
\end{table}

La estructura esperada del ZIP forma parte del mismo contrato. Cualquier backend nuevo debe respetar esta organización para que el middleware ubique el Dockerfile, las dependencias, los pesos, el vocabulario y la configuración del pipeline sin reglas especiales por modelo:

\begin{center}
\begin{minipage}{0.88\textwidth}
\hrule
\vspace{0.12cm}
\small
\begin{Verbatim}[baselinestretch=0.82]
mi_modelo.zip
|-- manifest.json              (obligatorio, en la raiz)
|-- backend/
|   |-- Dockerfile             (imagen Docker del backend)
|   |-- requirements.txt
|   `-- app/
|       |-- main.py            (endpoints /health e /infer)
|       `-- ...
|-- weights/                   (pesos del modelo)
|-- vocab/                     (vocabulario de glosas)
`-- config/
    `-- pipeline_config.json   (configuracion del pipeline)
\end{Verbatim}
\vspace{0.05cm}
\hrule

\captionof{figure}{Estructura esperada del paquete ZIP de un modelo}
\label{lst:estructura_zip_modelo}
\end{minipage}
\end{center}

Para mejorar la resiliencia se implementaron \textit{bootstrap manifests}. Durante la instalación, el middleware guarda una copia del manifest en \textbf{data/bootstrap\_manifests/} con el nombre \textbf{<model\_id>\_\_<version>.json}. Al iniciar, el middleware escanea ese directorio y vuelve a registrar los modelos si \textbf{registry.json} fue eliminado, siempre que el directorio persistente del anfitrión siga disponible.

El nombre del contenedor Docker se genera de forma determinística a partir del identificador y la versión del modelo:

\begin{center}
\begin{minipage}{0.55\textwidth}
\hrule
\vspace{0.12cm}
\small
\begin{Verbatim}[baselinestretch=0.82]
elan-ai-model-{model_id}-{version}
\end{Verbatim}
\vspace{0.05cm}
\hrule

\captionof{codigo}{Patrón de nombre del contenedor Docker de un modelo}
\label{lst:nombre_contenedor_modelo}
\end{minipage}
\end{center}

Con \textbf{model\_id=lsec\_bio\_gloss\_final\_v1} y \textbf{version=1.0.0}, el nombre generado es \textbf{elan-ai-model-lsec\_bio\_gloss\_final\_v1-1.0.0}. Esta regla evita almacenar un identificador adicional: el middleware puede localizar el contenedor derivando su nombre desde los metadatos. Las operaciones de construcción, creación, inicio e inspección se realizaron con el SDK de Docker para Python \cite{docker_sdk_python}.

La Figura~\ref{fig:flujo_instalacion} muestra el flujo de instalación completo. La ruta principal termina con el modelo registrado; la rama de error representa el rollback cuando falla la construcción, el arranque, la comprobación de disponibilidad o la persistencia.

\begin{figure}[htbp]
    \centering
    \includegraphics[width=1\textwidth]{Diagramas/Proceso de instalacion de modelo.png}
    \caption{Proceso de instalación de un modelo}
    \label{fig:flujo_instalacion}
\end{figure}

También se resolvió un caso frecuente de empaquetado en Windows. Si el ZIP conserva una carpeta raíz, el manifest puede quedar en \textbf{lsec\_bio\_gloss\_final\_v1/manifest.json} en lugar de la raíz del comprimido. El middleware detecta si todas las entradas comparten el mismo prefijo y lo elimina al leer el manifest y extraer archivos. De esta manera, un detalle de empaquetado no se convierte en un error difícil de interpretar para el usuario.

El criterio de verificación fue instalar \textbf{lsec\_bio\_gloss\_final\_v1} versión \textbf{1.0.0} mediante \textbf{POST \symbol{47}api\symbol{47}v1\symbol{47}models\symbol{47}install}. Se comprobó que el modelo apareciera como \textbf{available}, que \textbf{registry.json} almacenara la entrada completa, que existiera el bootstrap manifest y que una instalación duplicada respondiera HTTP 409 con \textbf{MODEL\_ALREADY\_EXISTS}.

\subsubsection{Fase 3 — Docker Runner e inferencia}

La tercera fase introdujo \textbf{DockerRunner}, el componente que ejecuta inferencias en el contenedor del modelo. Su trabajo es preparar el entorno antes de cada solicitud: verificar que el contenedor exista, iniciarlo si está detenido, comprobar su disponibilidad, construir la carga útil, enviar la solicitud al backend y adaptar la respuesta al contrato del middleware. Con esta separación, \textbf{JobService} queda libre de los detalles del SDK de Docker.

La comunicación entre el middleware y el backend ocurre dentro de la red \textbf{elan-ai-shared}. El middleware usa el nombre del contenedor como nombre de host y deja que el DNS interno de Docker lo resuelva. Gracias a ello no es necesario exponer un puerto del anfitrión por cada modelo, se elimina la dependencia de direcciones IP y varios backends pueden convivir en la misma red \cite{docker2026overview}.

El método \textbf{ensure\_container()} controla el ciclo de vida antes de inferir. Si el contenedor ya está ejecutándose, se reutiliza; si existe pero está detenido, se reinicia; si no existe pero la imagen sí, se crea un contenedor nuevo con red, nombre y volúmenes. Si falta la imagen, se lanza \textbf{DockerImageNotFoundError}, que el manejador global transforma en HTTP 404 para indicar que el modelo debe reinstalarse. Con esto, el sistema puede recuperarse de reinicios de Docker Desktop o del equipo anfitrión.

Un punto crítico fue traducir la ruta del video. ELAN selecciona archivos del sistema anfitrión, por ejemplo \textbf{C:/Users/imbaq/Videos/seña.mp4}; el backend, al ejecutarse en Linux dentro del contenedor, accede al archivo mediante \textbf{\symbol{47}data\symbol{47}videos}. Por ello, \textbf{DockerRunner} toma el nombre del archivo recibido y lo antepone a la ruta interna del contenedor, produciendo \textbf{\symbol{47}data\symbol{47}videos\symbol{47}seña.mp4}. La variable \textbf{MIDDLEWARE\_VIDEOS\_DIR} se pasa al SDK de Docker como cadena cruda para conservar correctamente rutas de Windows \cite{docker_sdk_python}.

Antes de llamar al backend, \textbf{DockerRunner} ejecuta polling contra \textbf{GET \symbol{47}health} cada 250 milisegundos. Se diferenciaron dos timeouts: \textbf{backend\_config.startup\_timeout\_sec = 180} durante instalación, porque el backend puede cargar pesos por primera vez; y \textbf{container.startup\_timeout\_sec = 30} durante inferencia, porque el contenedor ya debería estar preparado o reiniciarse con rapidez.

Los estados de un job se resumen en la Figura~\ref{fig:estados_job}. La figura funciona como contrato operativo para interpretar logs, métricas y respuestas: todo job inicia en \textbf{RECEIVED}, pasa por validación, preprocesamiento, cola y ejecución, y termina en \textbf{COMPLETED}, \textbf{FAILED} o \textbf{TIMEOUT}.

\begin{figure}[H]
    \centering
    \includegraphics[width=0.85\textwidth]{Diagramas Finales/Máquina de estados de un job.drawio.png}
    \caption{Máquina de estados de un job de inferencia}
    \label{fig:estados_job}
\end{figure}

Cada estado cumple una función específica: \textbf{VALIDATING} verifica el modelo; \textbf{PREPROCESSING} construye \textbf{InferenceInput} y traduce la ruta del video; \textbf{QUEUED} indica espera por el límite de concurrencia; \textbf{RUNNING} corresponde a la llamada de \textbf{DockerRunner}; \textbf{POSTPROCESSING} valida y adapta la respuesta; y los estados finales clasifican el resultado de la ejecución.

La respuesta exitosa incluye los segmentos y un campo de trazabilidad que registra el runner, la imagen Docker, el contenedor, los tiempos de validación, cola, inferencia y postprocesamiento, además del historial de estados. Esta trazabilidad facilita reconstruir la ejecución durante las pruebas y diagnosticar fallos en una arquitectura compuesta por varios servicios; el Anexo~II describe campo por campo esta información.

La Figura~\ref{fig:flujo_inferencia} relaciona los componentes de la inferencia: ELAN envía la solicitud, el middleware encola el job, DockerRunner prepara el contenedor, el backend procesa el video y el middleware adapta los segmentos para ELAN. El bloque interno del backend resume la canalización definida para el modelo utilizado en la validación; esta canalización pertenece al backend ejecutado en el contenedor y no al núcleo del middleware.

\begin{figure}[H]
    \centering
    \includegraphics[width=1\textwidth]{Diagramas Finales/Diagrama de secuencia del flujo completo de inferencia.drawio.png}
    \caption{Flujo completo de una inferencia desde ELAN hasta las anotaciones}
    \label{fig:flujo_inferencia}
\end{figure}

Los errores de \textbf{DockerRunner} se tradujeron a códigos HTTP con significado propio: \textbf{DockerNotAvailableError} responde 503; \textbf{DockerImageNotFoundError}, 404; \textbf{DockerContainerStartError} y \textbf{DockerContainerHealthcheckFailedError}, 502; \textbf{DockerInferenceError}, 502 con el detalle reportado por el backend; y \textbf{DockerTimeoutError}, 504. De este modo, ELAN, Postman o cualquier otro cliente identifica la causa de una falla sin necesidad de revisar los logs.

El criterio de verificación fue ejecutar una inferencia desde Postman y recibir al menos un segmento con \textbf{start\_ms}, \textbf{end\_ms}, \textbf{label} y \textbf{predictions}. También se comprobó que los logs mostraran la transición de \textbf{RECEIVED} a \textbf{COMPLETED} y que \textbf{trace} incluyera tiempos por etapa e identificador del contenedor.

\subsubsection{Fase 4 — Cola de jobs, métricas y estado TIMEOUT}

La cuarta fase añadió control de concurrencia, métricas y separación explícita entre \textbf{FAILED} y \textbf{TIMEOUT}. Estos componentes no modifican el pipeline de IA; regulan el acceso a recursos compartidos, hacen observable la ejecución y comunican fallos temporales de forma clara.

El primer componente fue \textbf{JobQueue}, la cola de trabajos. Su regla de operación es simple de enunciar: existe un número máximo de inferencias que pueden ejecutarse a la vez y, cuando ese cupo está ocupado, cada nueva solicitud espera su turno en estricto orden de llegada. Se descartó un mecanismo de sincronización básico, como un semáforo, porque limita la concurrencia pero no garantiza el orden de llegada; en su lugar, la cola se construyó con las primitivas de sincronización de hilos de Python \cite{python_threading}, combinando una estructura de espera ordenada con señales que despiertan al trabajo más antiguo cuando se libera un cupo. Así se controla la concurrencia y se garantiza un trato equitativo entre solicitudes.

El límite de concurrencia se expone en la variable \textbf{MIDDLEWARE\_MAX\_CONCURRENT\_JOBS}, con valor predeterminado de 1. Funciona como un control indirecto de recursos: impide que varios backends carguen sus pesos en memoria al mismo tiempo y, en equipos con mayor capacidad, permite ampliar la concurrencia sin modificar el código.

El segundo componente fue \textbf{MetricsService}. Este servicio mantiene contadores protegidos contra accesos simultáneos: jobs recibidos, completados, fallidos, expirados por tiempo de espera y un conteo de errores por código. Los campos \textbf{active\_jobs} y \textbf{queued\_jobs} no son contadores acumulados, sino una fotografía del estado de la cola tomada en el momento de consultar \textbf{GET /api/v1/metrics}.

La integración con el servidor Uvicorn exigió un cuidado adicional: una solicitud de inferencia que espera su turno mantiene ocupado un hilo de atención, por lo que el servidor debe tolerar solicitudes en espera y, al mismo tiempo, seguir respondiendo los endpoints de gestión, como el de salud y el de métricas. Si se amplía el límite de concurrencia, la capacidad del servidor debe crecer en proporción \cite{uvicorn_docs}.

El tercer componente fue el estado \textbf{TIMEOUT}. Antes, toda excepción de inferencia terminaba en \textbf{FAILED}; sin embargo, un error funcional del backend y una inferencia que excede el tiempo máximo no representan el mismo problema. Con \textbf{TIMEOUT}, las métricas separan fallos funcionales de fallos por latencia, facilitando el diagnóstico sin revisar logs.

Cuando una inferencia supera el tiempo máximo permitido, el job termina en \textbf{TIMEOUT}, el contador de métricas correspondiente se incrementa y el middleware responde HTTP 504 con el código \textbf{DOCKER\_TIMEOUT}. Esta distinción se usó en las pruebas de carga y en el escenario de tiempo de espera forzado de la validación final.

\subsubsection{Fase 5 — Integración con la interfaz de ELAN}

La quinta fase trasladó las capacidades del middleware a la interfaz gráfica de ELAN. Hasta ese momento, la instalación y la inferencia podían ejecutarse desde Postman o clientes HTTP; sin embargo, el objetivo del TIC requería que el investigador permaneciera dentro del entorno de anotación. Por ello se implementaron acciones para instalar modelos, consultar conectividad, seleccionar el modelo activo, configurar la inferencia, ejecutar el reconocedor y mostrar errores desde ELAN \cite{auer2014avatech}.

La integración se ubicó en el paquete \textbf{mpi.\allowbreak{}eudico.\allowbreak{}client.\allowbreak{}annotator.\allowbreak{}recognizer.\allowbreak{}ai}. La clase \textbf{AIOrchestration\allowbreak{}MiddlewareClient} concentra la comunicación HTTP mediante \textbf{java.net.http.\allowbreak{}HttpClient} y fija \textbf{HTTP/1.1}, porque Uvicorn no ofrece HTTP/2 por defecto y durante las pruebas se identificaron fallos cuando el cliente intentaba negociar otra versión \cite{java11_httpclient,uvicorn_docs}. Esta clase evita duplicar solicitudes en el panel, el diálogo de modelos y el reconocedor.

El método \textbf{installModel} envía el ZIP en el formato \textbf{multipart/form-data}, el formato estándar con el que un cliente HTTP adjunta archivos \cite{rfc7578}. Como Java 11 no incluye un constructor nativo para ese tipo de cuerpo, el cliente lo arma manualmente respetando la estructura que el middleware espera al recibir archivos:

\begin{center}
\begin{minipage}{0.82\textwidth}
\hrule
\vspace{0.12cm}
\small
\begin{Verbatim}[baselinestretch=0.82]
--FormBoundary<uuid>
Content-Disposition: form-data; name="file"; filename="modelo.zip"
Content-Type: application/octet-stream

[bytes del archivo ZIP]
--FormBoundary<uuid>--
\end{Verbatim}
\vspace{0.05cm}
\hrule

\captionof{codigo}{Cuerpo multipart enviado por el cliente Java para instalar modelos}
\label{lst:multipart_install_model}
\end{minipage}
\end{center}

Un detalle de este formato provocó un error difícil de rastrear durante las pruebas: el delimitador declarado en la cabecera de la solicitud no debe llevar los dos guiones iniciales que sí llevan los delimitadores dentro del cuerpo. El diagnóstico de este problema y su solución quedaron registrados en la tabla de problemas frecuentes del Anexo~IV.

La clase \textbf{AIModel\allowbreak{}ManagerDialog} implementa el diálogo Swing de gestión de modelos. Incluye indicador de conexión, prueba de conectividad, tabla de modelos, instalación de ZIP, activación, desactivación y log de actividad. Las operaciones HTTP se ejecutan con \textbf{SwingWorker} para no bloquear el EDT de Swing; al finalizar, \textbf{done()} actualiza la interfaz en el hilo seguro para componentes gráficos \cite{java_swingworker}.

La clase \textbf{AIOrchestration\allowbreak{}RecognizerPanel} reemplazó parámetros hardcodeados por configuración visible: URL del middleware, modelo instalado, nivel de anotación destino, modo de etiquetado y video. Los modelos se cargan en segundo plano desde \textbf{addNotify()}, porque este método se ejecuta cuando el panel entra en la jerarquía visible y evita llamadas de red durante el constructor. Los valores sugeridos de nivel de anotación y modo se toman del campo \textbf{ui} del manifest.

La clase \textbf{AIOrchestration\allowbreak{}Recognizer} lee en tiempo de ejecución los valores seleccionados en el panel: modelo, nivel de anotación, modo de etiquetado y URL base. Con ello, la solicitud enviada al middleware refleja la configuración visible del usuario y no valores definidos en tiempo de compilación. También se crearon DTOs en \textbf{.ai.dto} para separar comunicación, representación de datos y lógica gráfica.

\textbf{SegmentVideoRequest} incorpora el método de fábrica \textbf{create}, que construye el payload con job, video, modelo, versión, nivel de anotación, etiqueta, modo de etiquetado, timeout de 300 segundos y \textbf{parameters} como objeto vacío. Esta decisión permite que el backend use su propia configuración interna y garantiza que el campo \textbf{parameters} siempre esté presente en la solicitud.

La comunicación entre \textbf{AIModel\allowbreak{}ManagerDialog} y \textbf{AIOrchestration\allowbreak{}RecognizerPanel} se resolvió con un flag booleano. Al cerrar el diálogo, el panel consulta \textbf{isModelListChanged()} y refresca el combo si hubo cambios. Finalmente, \textbf{addNotify()} dispara la carga inicial de modelos mediante \textbf{SwingWorker}; si el middleware no está disponible, el error aparece cuando el usuario visualiza el panel y puede interpretarlo en contexto.

\subsubsection{Fase 6 — Validación final}

La sexta fase verificó el sistema completo: middleware, Docker, backend del modelo y reconocedor integrado en ELAN. La validación combinó pruebas de caja negra y de caja blanca, porque no bastaba con probar endpoints aislados; había que asegurar la instalación, la activación, la selección en ELAN, la inferencia, el manejo de errores, las métricas y las anotaciones temporales. Las pruebas se realizaron con el modelo \textbf{lsec\_bio\_gloss\_final\_v1} y videos de LSEc, manteniendo la arquitectura real del sistema.

Se utilizó el modelo real en lugar de uno ficticio. Un modelo simulado habría validado la serialización y el transporte HTTP, pero no habría ejercitado la construcción de imágenes, el arranque de contenedores, la carga de pesos, la lectura de videos, la ejecución del pipeline ni la generación de segmentos reales. Esta decisión encareció las pruebas, aunque produjo evidencia representativa del escenario para el que se diseñó el middleware.

Los escenarios de pruebas se identifican con códigos E-XX. Estos códigos no corresponden a errores o estados HTTP, sino que son etiquetas internas de trazabilidad. Por ejemplo, E-07 corresponde a una inferencia exitosa desde ELAN; E-08, a la revisión visual de anotaciones en la línea de tiempo; y E-13, al tiempo de espera forzado. Esta nomenclatura relaciona la acción, el resultado esperado, la evidencia y la categoría de cada prueba.

Las pruebas de caja negra exploraron el comportamiento visible desde las interfaces expuestas al usuario final. Se definieron 15 escenarios en cuatro categorías: conectividad, gestión de modelos, inferencia y manejo de errores. Los escenarios cubrieron la comprobación de disponibilidad, la instalación, el listado y detalle de modelos, la activación y desactivación, la inferencia desde ELAN y desde Postman, la inserción de anotaciones en el nivel correcto y los fallos controlados: video inexistente, modelo desactivado, middleware detenido, contenedor detenido y tiempo de espera forzado, junto con la revisión de métricas. En cada caso se registraron la acción, la respuesta esperada y el comportamiento visible en ELAN, comprobando tanto la semántica de las respuestas HTTP como la experiencia del investigador.

Las pruebas de caja blanca revisaron estados internos que no se observan solo desde la API: la persistencia en \textbf{registry.json}, la generación del bootstrap manifest, la recuperación automática del registro, el incremento de métricas, la traducción de la ruta del video en los logs y la configuración del bind mount mediante \textbf{docker inspect}. Estas verificaciones confirmaron que el estado persistente, los montajes y los contadores internos quedaban consistentes después de cada operación.

Para rendimiento se usó \textbf{docker stats} en lugar de integrar \textbf{psutil}. La razón es que el consumo se distribuye entre dos contenedores: el middleware mantiene API, cola, contratos y orquestación; el backend ejecuta el pipeline, carga pesos y procesa video. Medir solo el proceso del middleware habría omitido el costo real de inferencia. \textbf{docker stats} permitió observar ambos contenedores desde el anfitrión y estimar el impacto sobre el hardware del usuario \cite{docker_cli_stats}. Los valores esperados para el equipo de desarrollo fueron menos de 5\% de CPU y menos de 150 MB de RAM para el middleware, y entre 60\% y 200\% de CPU junto con 1 a 4 GB de RAM para el backend.

El escenario \textbf{E-07} validó el flujo completo desde ELAN. El usuario selecciona el modelo, el video, el nivel de anotación destino y ejecuta el reconocedor. ELAN construye la solicitud, el middleware valida el contrato, localiza el modelo, traduce la ruta del video y solicita inferencia al backend. El resultado esperado es progreso visible y anotaciones en \textbf{AUTO\_GLOSS\_SEGMENTS}; sus tiempos deben coincidir con los intervalos visibles de señas. El escenario \textbf{E-08} complementa esta evidencia con la captura de la línea de tiempo donde se observan las anotaciones alineadas.

El escenario \textbf{E-13} validó timeout controlado. Desde Postman se envió una inferencia con \textbf{execution.timeout\_sec: 1}, valor menor al tiempo real de procesamiento. El resultado esperado fue HTTP 504 con \textbf{DOCKER\_TIMEOUT}. Después del error, el middleware debía seguir disponible, sin bloquear la cola, perder el registro ni dejar estados inconsistentes.

En la prueba de carga controlada se enviaron dos inferencias simultáneas con \textbf{MIDDLEWARE\_MAX\_CONCURRENT\_JOBS=1}. El comportamiento esperado era que la primera pasara a \textbf{RUNNING} y la segunda quedara en \textbf{QUEUED}; al terminar la primera, la segunda debía continuar en orden FIFO. Los logs registraron la transición de \textbf{QUEUED} a \textbf{RUNNING} y las métricas finales reportaron \textbf{total\_jobs: 2} con ambos trabajos completados. La prueba confirmó que la cola limita la concurrencia, preserva el orden de llegada y mantiene las métricas consistentes.


\subsection{Contratos de comunicación entre componentes}

El sistema se apoya en tres contratos de comunicación que ordenan el intercambio entre ELAN, el middleware y los backends de modelo. Estos contratos funcionan como una frontera clara entre tecnologías: ELAN mantiene su flujo de reconocedores AVATecH en Java, el middleware concentra la orquestación en Python y cada modelo ejecuta su inferencia dentro de un contenedor aislado. Por esta razón, los intercambios se definieron como cuerpos JSON sobre endpoints HTTP, una decisión coherente con el enfoque REST y con la necesidad de contar con interfaces verificables, documentables y desacopladas \cite{fielding2000rest}.

La estructura general de estos contratos se mantuvo estable durante el desarrollo. Los ajustes realizados en las iteraciones afectaron la validación, la compatibilidad o el manejo de errores, pero nunca la responsabilidad principal de cada componente. Esa estabilidad resultó clave para integrar ELAN, Docker y los backends de IA sin depender de los detalles internos de cada uno.

\subsubsection{Contrato entre ELAN y el Middleware (solicitud de inferencia)}

ELAN inicia el flujo mediante una solicitud \textbf{POST /api/v1/jobs/segment-video}. Esta operación no ejecuta reconocimiento dentro de ELAN; solo entrega al middleware los datos necesarios para identificar el trabajo, ubicar el video, seleccionar el modelo, definir el nivel de anotación destino y establecer un tiempo máximo de ejecución. Antes de iniciar operaciones costosas, el middleware valida la estructura recibida con modelos Pydantic integrados a FastAPI, de modo que errores de formato, campos incompletos o valores incompatibles se detecten en la entrada y no durante la ejecución del contenedor:

\begin{center}
\begin{minipage}{0.82\textwidth}
\hrule
\vspace{0.12cm}
\small
\begin{Verbatim}[baselinestretch=0.82]
{
  "job_id": "identificador-unico-del-job",
  "media": {"path": "C:/Users/.../video.mp4"},
  "model": {"model_id": "nombre_modelo", "version": "1.0.0"},
  "annotation": {
    "target_tier": "NombreTier",
    "default_label": "REGION",
    "label_mode": "gloss_top1"
  },
  "execution": {"timeout_sec": 300}
}
\end{Verbatim}
\vspace{0.05cm}
\hrule

\captionof{codigo}{Solicitud JSON enviada desde ELAN al middleware}
\label{lst:request_elan_middleware}
\end{minipage}
\end{center}

El campo \textbf{job\_id} permite seguir una solicitud en logs, métricas y consultas posteriores mediante \textbf{GET /api/v1/jobs/\{job\_id\}}. El objeto \textbf{media} contiene la ruta del video seleccionada desde ELAN, pero el middleware no la entrega al backend sin procesamiento, porque el contenedor ve una jerarquía de archivos distinta a la del anfitrión. El objeto \textbf{model} representa una versión concreta instalada, mientras que \textbf{annotation} define cómo se transformarán los segmentos en anotaciones: \textbf{target\_tier} indica el nivel de anotación destino, \textbf{default\_label} actúa como etiqueta de respaldo y \textbf{label\_mode} decide entre la glosa principal o una etiqueta simple. Finalmente, \textbf{execution.timeout\_sec} limita el tiempo de espera; si el backend lo supera, el middleware responde \textbf{DOCKER\_TIMEOUT} con HTTP 504.

\subsubsection{Contrato entre Middleware y el Backend de modelo (payload de inferencia)}

Después de validar la solicitud y seleccionar el modelo, el middleware traduce el contrato externo al contrato interno del backend. El cambio más importante ocurre en \textbf{media.path}: la ruta seleccionada en ELAN se convierte en la ruta visible dentro del contenedor según el bind mount de videos. Esta traducción es necesaria porque Docker aísla el sistema de archivos de cada contenedor y expone solo los directorios montados explícitamente \cite{docker2026overview}:

\begin{center}
\begin{minipage}{0.70\textwidth}
\hrule
\vspace{0.12cm}
\footnotesize
\begin{Verbatim}[baselinestretch=0.78]
{
  "job_id": "identificador-unico-del-job",
  "media": {"path": "/data/videos/video.mp4"},
  "annotation": { ... },
  "parameters": {}
}
\end{Verbatim}
\vspace{0.05cm}
\hrule

\captionof{codigo}{Payload interno enviado desde el middleware al backend}
\label{lst:payload_middleware_backend}
\end{minipage}
\end{center}

El campo \textbf{parameters} se incluye siempre, aunque sea un objeto vacío. Con ello, el backend puede leer una estructura uniforme y se evitan casos especiales por ausencia del campo. Al mismo tiempo, esto no significa que ELAN pueda sobrescribir libremente la configuración interna del modelo: el pipeline conserva su configuración base en \textbf{config/pipeline\_config.json}. Por tanto, \textbf{parameters} queda como un espacio de extensión controlado, sin reemplazar los artefactos ni las reglas declaradas en el paquete del modelo.

\subsubsection{Contrato entre el Backend de un modelo y el Middleware (respuesta de inferencia)}

Cuando termina la inferencia, el backend devuelve un JSON con la salida técnica del pipeline. El middleware no solo acepta una respuesta porque el código HTTP sea exitoso, sino que también verifica que el tipo de salida, el tiempo del video y una lista de segmentos con marcas de inicio y fin estén presentes. Esta validación evita que una respuesta incompleta llegue a ELAN y genere anotaciones inconsistentes en la línea de tiempo:

\begin{center}
\begin{minipage}{0.92\textwidth}
\hrule
\vspace{0.12cm}
\footnotesize
\begin{Verbatim}[baselinestretch=0.76]
{
  "output_type": "segments_with_gloss",
  "media_info": {"fps": 59.94, "duration_ms": 8508, "total_frames": 510},
  "segments": [{
    "start_ms": 884, "end_ms": 1518,
    "label": "TRABAJAR", "confidence": 0.164,
    "predictions": [
      {"rank": 1, "gloss_id": 10, "gloss": "TRABAJAR", "probability": 0.164},
      {"rank": 2, "gloss_id": 7, "gloss": "VER", "probability": 0.126}
    ]
  }]
}
\end{Verbatim}
\vspace{0.05cm}
\hrule

\captionof{codigo}{Respuesta JSON del backend con segmentos y glosas}
\label{lst:respuesta_backend_middleware}
\end{minipage}
\end{center}

El campo \textbf{output\_type} declara la forma semántica de la respuesta; en esta implementación, \textbf{segments\_with\_gloss} indica intervalos de tiempo acompañados de una etiqueta principal y de alternativas de predicción. El objeto \textbf{media\_info} aporta contexto del video procesado: \textbf{fps}, \textbf{duration\_ms} y \textbf{total\_frames}. Dentro de cada segmento, \textbf{label} es la etiqueta que se anotará, \textbf{confidence} refleja la confianza reportada por el backend y \textbf{predictions} lista las alternativas ordenadas por \textbf{rank}. La respuesta contiene así lo mínimo necesario para anotar en ELAN y, al mismo tiempo, información adicional útil para la revisión posterior. El Anexo~II presenta un ejemplo completo de esta respuesta, incluida la trazabilidad por etapas del campo \textbf{trace}.

La Tabla~\ref{tab:endpoints_api} resume los endpoints disponibles en la API REST del middleware, agrupados según su responsabilidad: disponibilidad del servicio, gestión de modelos, ejecución y consulta de inferencias, y lectura de métricas. Gracias a esta organización, ELAN consume únicamente los endpoints que su flujo necesita, mientras que herramientas como Postman o \textbf{curl} sirven de apoyo en la verificación durante las pruebas funcionales.

\begin{table}[H]
\centering
\caption{Endpoints de la API REST del middleware}
\label{tab:endpoints_api}
\small
\renewcommand{\arraystretch}{1.18}
\setlength{\tabcolsep}{5pt}

\begin{tabular}{
    >{\centering\arraybackslash}p{1.6cm}
    >{\raggedright\arraybackslash}p{5.0cm}
    >{\raggedright\arraybackslash}p{6.7cm}
}
\toprule
\textbf{Método} & \textbf{Endpoint} & \textbf{Descripción} \\
\midrule

\texttt{GET} &
\texttt{/health} &
Verifica la disponibilidad general del servicio. \\

\texttt{GET} &
\texttt{/api/v1/models} &
Lista todos los modelos instalados en el middleware. \\

\texttt{GET} &
\texttt{/api/v1/models/\{id\}} &
Recupera el detalle completo de un modelo específico. \\

\texttt{POST} &
\texttt{/api/v1/models/install} &
Instala un modelo a partir de un archivo ZIP enviado como \texttt{multipart/form-data}. \\

\texttt{PATCH} &
\texttt{/api/v1/models/\{id\}/status} &
Activa o desactiva un modelo previamente instalado. \\

\texttt{POST} &
\texttt{/api/v1/jobs/segment-video} &
Ejecuta una inferencia sobre un video seleccionado. \\

\texttt{GET} &
\texttt{/api/v1/jobs/\{job\_id\}} &
Recupera el resultado de un job de inferencia ejecutado. \\

\texttt{GET} &
\texttt{/api/v1/metrics} &
Consulta métricas acumuladas desde el inicio del servicio. \\

\bottomrule
\end{tabular}
\end{table}

\subsection{Despliegue y configuración del entorno}

El despliegue del middleware se realiza con Docker Compose porque el sistema necesita más que un proceso aislado: requiere una red compartida, comunicación con contenedores de modelos creados dinámicamente y persistencia entre reinicios. El archivo \textbf{docker-compose.yml} define el servicio principal, los bind mounts y la red común usada por los backends; su contenido completo, junto con las consideraciones de seguridad sobre el socket de Docker, consta en el Anexo~III. Esta decisión simplifica el arranque del sistema, ya que el usuario puede levantar el servicio con un comando único mientras Compose crea la red, publica el puerto y monta los directorios requeridos.

Se eligieron bind mounts para los datos persistentes porque dejan los archivos visibles en el sistema del anfitrión. Esto facilita auditar el registro de modelos, revisar manifests de respaldo y conservar información aunque se detengan o eliminen los contenedores. En cambio, los volúmenes con nombre pueden perderse con operaciones de limpieza de Docker, lo que no resulta conveniente para un entorno local de investigación \cite{docker2026overview}.

El middleware utiliza tres montajes con responsabilidades distintas. \textbf{./data/models\_store} conserva el registro y los archivos extraídos de los modelos instalados. \textbf{./data/bootstrap\_manifests} guarda copias de manifests para recuperar información mínima si el registro principal se pierde o queda inconsistente. Además, el socket de Docker permite que el SDK de Docker (para Python) construya imágenes, inicie contenedores, inspeccione estados y gestione backends desde el propio middleware \cite{docker_sdk_python}. Así, aunque el orquestador se ejecuta dentro de un contenedor, mantiene control local sobre los contenedores de inferencia.

La Tabla~\ref{tab:variables_entorno} describe las variables configurables del middleware. Estas variables separan la configuración del entorno respecto del código fuente, algo necesario porque cada estación de trabajo puede tener una ruta distinta para videos, diferente capacidad de hardware o distintos requisitos de logging.

\begin{table}[H]
\centering
\caption{Variables de entorno del middleware}
\label{tab:variables_entorno}
\scriptsize
\renewcommand{\arraystretch}{1.18}
\setlength{\tabcolsep}{5pt}

\begin{tabular}{
    >{\raggedright\arraybackslash}p{4.2cm}
    >{\raggedright\arraybackslash}p{3.4cm}
    >{\raggedright\arraybackslash}p{5.8cm}
}
\toprule
\textbf{Variable} & \textbf{Valor por defecto} & \textbf{Descripción} \\
\midrule

\texttt{MIDDLEWARE\_HOST} &
\texttt{0.0.0.0} &
Interfaz de escucha dentro del contenedor. \\

\texttt{MIDDLEWARE\_PORT} &
\texttt{8000} &
Puerto expuesto al host. \\

\texttt{MIDDLEWARE\_LOG\_LEVEL} &
\texttt{INFO} &
Nivel de logging estructurado. \\

\texttt{MIDDLEWARE\_MODELS\_}\newline
\texttt{STORE\_DIR} &
\texttt{/app/data/}\newline
\texttt{models\_store} &
Directorio donde se almacena el registry de modelos. \\

\texttt{MIDDLEWARE\_BOOTSTRAP\_}\newline
\texttt{MANIFESTS\_DIR} &
\texttt{/app/data/bootstrap/}\newline
\texttt{manifests} &
Directorio persistente para los manifests de respaldo. \\

\texttt{MIDDLEWARE\_VIDEOS\_DIR} &
\textit{requerido} &
Ruta del host donde se encuentran los videos disponibles para inferencia. \\

\texttt{MIDDLEWARE\_DOCKER\_}\newline
\texttt{NETWORK} &
\texttt{elan-ai-shared} &
Red Docker compartida entre el middleware y los backends de los modelos. \\

\texttt{MIDDLEWARE\_MAX\_}\newline
\texttt{CONCURRENT\_JOBS} &
\texttt{1} &
Límite de inferencias simultáneas permitidas por el middleware. \\

\bottomrule
\end{tabular}
\end{table}

La variable \textbf{MIDDLEWARE\_VIDEOS\_DIR} es la única requerida sin valor por defecto porque depende del equipo donde se despliega el sistema. Debe apuntar al directorio del anfitrión que contiene los videos de LSEc disponibles para inferencia. El middleware no intenta inferir esta ruta: la entrega al SDK de Docker como fuente del bind mount usado por los contenedores de modelo. Esta decisión evita transformaciones ambiguas entre sistemas operativos y se relaciona con el problema resuelto en la Fase 3, donde ELAN trabaja con rutas del anfitrión y el backend con rutas internas del contenedor.

Debemos añadir que el despliegue local para un solo usuario es el escenario previsto para el TIC, pero no un límite de la arquitectura. Por diseño, el puerto del middleware solo está disponible en la interfaz de loopback (127.0.0.1) para que solo el ELAN de la misma máquina pueda usarlo. Sin embargo, dado que es una relación cliente-servidor sobre HTTP, sería suficiente publicar el puerto en la dirección IP de la computadora en una red local para que múltiples estaciones de trabajo con ELAN pudieran enviar solicitudes al mismo middleware, y se establecería la dirección en el campo URL del panel del reconocedor. En ese caso, la cola FIFO atendería las solicitudes en orden de llegada y la variable \textbf{MIDDLEWARE\_MAX\_CONCURRENT\_JOBS} determinaría cuántas inferencias se ejecutan en paralelo según la capacidad del servidor; si el valor predeterminado es 1, las solicitudes de otros clientes esperarían su turno. Este escenario y la posibilidad de mover el middleware y los modelos a un servidor de alto rendimiento o a la nube se revisa en las recomendaciones del capítulo de resultados.

Para levantar el sistema por primera vez, o después de cambios en el código, se ejecutan los comandos mostrados a continuación. El primer comando construye la imagen del middleware y deja el servicio en segundo plano; el segundo verifica que el endpoint de salud responda antes de instalar modelos o ejecutar inferencias:

\begin{center}
\begin{minipage}{0.78\textwidth}
\hrule
\vspace{0.10cm}
\footnotesize
\begin{Verbatim}[baselinestretch=0.78]
# Desde el directorio middleware/
docker compose up --build -d

# Verificar que el servicio está activo
curl http://localhost:8000/health
\end{Verbatim}
\vspace{0.03cm}
\hrule

\captionof{comando}{Comandos iniciales para desplegar y verificar el middleware}
\label{lst:comandos_despliegue_middleware}
\end{minipage}
\end{center}

La integración con ELAN requiere compilar el árbol fuente modificado con Maven e incorporar el JAR resultante a la instalación de la herramienta; el procedimiento paso a paso, incluido el registro del reconocedor mediante el mecanismo SPI de Java, se describe en el Anexo~IV. Esta integración no reemplaza el flujo de anotación existente; añade un reconocedor dentro del mecanismo de extensión previsto por AVATecH \cite{auer2014avatech}. Una vez instalado, el reconocedor \textit{AI Orchestration Middleware} aparece en la lista de reconocedores de ELAN y permite seleccionar modelo, configurar el nivel de anotación destino y ejecutar inferencias sin manipular directamente los contenedores Docker.


\subsection{Estructura de directorios del sistema}

La Figura~\ref{lst:estructura_directorios} muestra cómo se organizó el middleware para reflejar la separación de responsabilidades definida en el diseño. La capa \textbf{api/} cepta solicitudes HTTP y serializa las respuestas; \textbf{core/} es la capa donde se centraliza la configuración y registro; \textbf{schemas/} tiene los modelos Pydantic que hacen ejecutables los contratos; \textbf{services/} agrupa la lógica de instalación, registro, jobs, Docker, cola y métricas; \textbf{runners/} abstrae el mecanismo de ejecución; \textbf{storage/} conserva jobs en memoria durante la sesión; y \textbf{data/} mantiene información persistente mediante bind mounts.

\begin{center}
\begin{minipage}{0.96\textwidth}
\hrule
\vspace{0.10cm}
\scriptsize
\begin{Verbatim}[baselinestretch=0.72]
middleware/
|-- app/
|   |-- main.py                    (FastAPI: inicio, routers, error handlers)
|   |-- api/
|   |   |-- routes_health.py       (GET /health)
|   |   |-- routes_models.py       (CRUD modelos e instalación)
|   |   |-- routes_jobs.py         (POST /jobs/segment-video)
|   |   `-- routes_metrics.py      (GET /metrics)
|   |-- core/
|   |   |-- config.py              (Variables de entorno -> Settings)
|   |   `-- logging_config.py      (Logging estructurado)
|   |-- runners/
|   |   |-- base_runner.py         (Clase abstracta BaseRunner)
|   |   |-- docker_runner.py       (DockerRunner: runtime activo)
|   |   `-- runner_selector.py     (Selección de runner por manifest)
|   |-- schemas/                   (Contratos Pydantic)
|   |-- services/
|   |   |-- model_registry_service.py  (Instalar y gestionar modelos)
|   |   |-- job_service.py             (Orquestar inferencia completa)
|   |   |-- docker_service.py          (Wrapper del SDK Docker)
|   |   |-- docker_lifecycle_service.py (Build, start y comprobación de disponibilidad)
|   |   |-- job_queue.py               (Cola FIFO y concurrencia)
|   |   `-- metrics_service.py         (Contadores thread-safe)
|   `-- storage/
|       `-- memory_store.py        (Jobs en memoria por sesión)
|-- data/
|   |-- models_store/              (Registry y archivos de modelos)
|   `-- bootstrap_manifests/       (Manifests de respaldo)
|-- Dockerfile                     (Imagen del middleware)
|-- docker-compose.yml             (Orquestación del despliegue)
`-- requirements.txt               (Dependencias Python)
\end{Verbatim}
\vspace{0.03cm}
\hrule

\captionof{figure}{Estructura de directorios del middleware}
\label{lst:estructura_directorios}
\end{minipage}
\end{center}

Los archivos Java añadidos a ELAN se ubicaron en el paquete \textbf{mpi.eudico.client.annotator.recognizer.ai} como un subpaquete del módulo de reconocedores, de modo que el nuevo código no esté demasiado lejos del punto de extensión AVATecH, al mismo tiempo que distingue la integración de IA de los reconocedores anteriores. La Figura~\ref{lst:estructura_clases_elan} muestra esta organización en la que el cliente HTTP, el reconocedor, el panel, el cuadro de diálogo del modelo y los DTO principales están agrupados en el mismo grupo. La lista completa de clases y DTOs está disponible en el Anexo~IV.

\begin{center}
\begin{minipage}{0.96\textwidth}
\hrule
\vspace{0.10cm}
\scriptsize
\begin{Verbatim}[baselinestretch=0.72]
src/main/java/mpi/eudico/client/annotator/recognizer/ai/
|-- AIOrchestrationMiddlewareClient.java    (cliente HTTP del middleware)
|-- AIOrchestrationMiddlewareException.java (excepción con HTTP status y body)
|-- AIOrchestrationRecognizer.java          (reconocedor AVATecH)
|-- AIOrchestrationRecognizerPanel.java     (panel de configuración en ELAN)
|-- AIModelManagerDialog.java               (diálogo de gestión de modelos)
`-- dto/
    |-- ModelSummary.java                   (resumen de modelo para lista)
    |-- InstalledModelDetail.java           (detalle completo con UI config)
    |-- ModelUiConfig.java                  (configuración UI del modelo)
    |-- SegmentVideoRequest.java            (solicitud de inferencia)
    `-- SegmentVideoResponse.java           (respuesta con segmentos)
\end{Verbatim}
\vspace{0.03cm}
\hrule

\captionof{figure}{Estructura de clases Java integradas en ELAN}
\label{lst:estructura_clases_elan}
\end{minipage}
\end{center}


\subsection{Guía de replicación: creación de un nuevo backend de modelo}

Una propiedad central del sistema es que admite nuevos modelos sin modificar el código de ELAN ni del middleware. Para lograrlo, cada modelo empaqueta su propia variabilidad: dependencias, artefactos, configuración, forma de ejecución y contrato HTTP. Si el paquete ZIP cumple el contrato descrito en la Fase 2 y detallado en el Anexo~I, el middleware puede instalarlo, construir su imagen, iniciar el contenedor y mostrarlo en ELAN con el mismo flujo usado por los modelos ya validados.

El primer paso es implementar el backend como un servicio HTTP que exponga los endpoints requeridos. Puede usarse FastAPI u otro framework, porque el middleware no depende de la implementación interna del modelo; solo necesita que el contenedor cumpla el contrato. \textbf{GET /health} debe responder HTTP 200 únicamente cuando el pipeline esté listo para inferir, ya que un contenedor iniciado no siempre significa que pesos, vocabularios o recursos de visión artificial estén cargados. \textbf{POST /infer} debe aceptar el payload definido en la sección de contratos y responder con \textbf{output\_type}, \textbf{media\_info} y \textbf{segments}. Cada segmento debe incluir al menos \textbf{start\_ms}, \textbf{end\_ms} y \textbf{label}; los campos \textbf{confidence} y \textbf{predictions} pueden añadirse cuando el modelo entregue evidencia adicional.

El segundo paso es empaquetar el backend con sus pesos y configuración, colocando \textbf{manifest.json} en la raíz del paquete. Este manifest identifica el modelo, su versión, los artefactos requeridos y la forma de construir y ejecutar el backend. El \textbf{Dockerfile} debe instalar las dependencias necesarias y definir cómo inicia el servicio HTTP. Esta encapsulación permite aislar bibliotecas de IA potencialmente incompatibles con ELAN o con otros modelos, manteniendo el principio de contenedorización de la arquitectura \cite{docker2026overview}. Los artefactos, como pesos, vocabularios y archivos de configuración, se declaran con rutas relativas dentro del ZIP para que el paquete pueda instalarse en distintos equipos.

El tercer paso es comprimir la estructura de directorios en un ZIP, cuidando que \textbf{manifest.json} quede en la raíz y no dentro de un subdirectorio adicional. Esta condición permite que el instalador valide el contrato antes de extraer y construir el backend. En Windows PowerShell, el comando correcto es:

\begin{center}
\begin{minipage}{0.78\textwidth}
\hrule
\vspace{0.08cm}
\footnotesize
\begin{Verbatim}[baselinestretch=0.78]
cd carpeta_que_contiene_el_paquete\
Compress-Archive -Path nombre_paquete\* -DestinationPath modelo.zip
\end{Verbatim}
\vspace{0.03cm}
\hrule

\captionof{comando}{Comando PowerShell para comprimir un paquete de modelo}
\label{lst:comprimir_modelo_zip}
\end{minipage}
\end{center}

El cuarto paso es instalar el paquete mediante \textbf{POST /api/v1/models/install} o desde ELAN con el botón \textit{Instalar ZIP}. En ambos casos, el flujo técnico es el mismo: el middleware recibe el ZIP como \textbf{multipart/form-data}, localiza y valida el manifest, verifica artefactos, extrae el paquete, actualiza el registro, construye la imagen Docker y ejecuta las comprobaciones de arranque. Si todo es correcto, el modelo queda instalado, el contenedor se inicia y aparece disponible en el panel del reconocedor.

Este proceso concentra el esfuerzo de integración en el backend y en la preparación correcta del paquete. El desarrollador no necesita conocer el código interno de ELAN ni del middleware; debe respetar el contrato de \textbf{/health}, \textbf{/infer} y \textbf{manifest.json}. Con ello se conserva el desacoplamiento de la arquitectura: ELAN no conoce las bibliotecas del modelo, el middleware no interpreta la lógica interna del pipeline y el backend no manipula archivos EAF ni estructuras propias de ELAN.

\subsection{Uso ético de herramientas de inteligencia artificial generativa}

El presente Trabajo de Integración Curricular es de autoría propia. Durante su desarrollo se utilizaron de manera ética y responsable herramientas de inteligencia artificial generativa (ChatGPT y Claude) como apoyo técnico, en concordancia con las orientaciones internacionales sobre el uso de esta tecnología en contextos de investigación \cite{unesco2023guidance}.

Su aporte principal se dio en la etapa de ingeniería inversa del código fuente de ELAN 7.1. Dado que la herramienta carece de documentación interna detallada sobre su mecanismo de reconocedores, estas herramientas se emplearon para analizar el paquete \textbf{mpi.eudico.client.annotator.recognizer}, ubicar los puntos de extensión del protocolo AVATecH y orientar la forma de implementar las interfaces requeridas, como \textbf{Recognizer2}, \textbf{IncrementalRecognizer} y la comunicación mediante \textbf{RecognizerHost}. También fueron útiles para la interpretación de errores de integración durante las pruebas, para revisar fragmentos de código repetitivos del prototipo y para ayudarnos con algunas ideas claves de lo que era viable o no realizar.

Aunque la asistencia de IA me proporcionó el primer conocimiento de un código fuente grande y extranjero, fue necesario verificar cada sugerencia contra el código real de ELAN 7.1 y la especificación de AVATecH, adaptada al comportamiento observado por las pruebas de integración, y en algunos casos rechazada porque no estaba en la versión de la herramienta utilizada. Por el contrario, al leer código de terceros, desarrollar integraciones entre varias tecnologías y analizar los resultados generados por estas herramientas, desarrollé mis habilidades para leer el código de otras personas.

Las decisiones metodológicas, los ajustes técnicos, la validación de resultados y las conclusiones corresponden exclusivamente a mi autoría.


\section[RESULTADOS, CONCLUSIONES Y RECOMENDACIONES]{Resultados, conclusiones y recomendaciones}

\subsection{Resultados}

Los resultados de este capítulo están organizados de acuerdo con los tres diferentes objetivos específicos del trabajo, de modo que cada evidencia pueda vincularse a una decisión de diseño, una acción de implementación y un criterio de verificación concreto. Organizarlo de esta manera nos permite distinguir entre los productos construidos, los contratos técnicos que los respaldan y el comportamiento observado durante la validación final; por lo tanto, cada bloque describe los entregables que pudimos lograr, la forma en que se articulan con la arquitectura general y las métricas o evidencias con las que se verificó su funcionalidad.

% ── Objetivo 1 ──────────────────────────────────────────────────────────────
\subsubsection{Análisis del protocolo AVATecH e integración con ELAN}

El primer objetivo fue estudiar el mecanismo de extensión AVATecH de ELAN (es decir, la documentación oficial y la ingeniería inversa del código fuente) para definir los contratos de la interfaz REST y el esquema de instalación del modelo y cómo incorporar servicios de inferencia externos en el flujo de anotación. De su cumplimiento se desprenden dos resultados que se complementan entre sí. El primero son los contratos para la solicitud y respuesta de inferencia a través de HTTP (en el Anexo II) y el esquema manifest.json para instalar modelos (en el Anexo I). El segundo es su materialización en código: cinco clases principales y un conjunto de DTO implementados en el paquete \textbf{mpi.eudico.\allowbreak{}client.\allowbreak{}annotator.\allowbreak{}recognizer.\allowbreak{}ai}, organizados según la Figura~\ref{lst:estructura_clases_elan}; el inventario completo del paquete consta en el Anexo~IV. La pieza central de esta integración es el reconocedor \textbf{AIOrchestrationRecognizer}, que implementa la interfaz \textbf{Recognizer} de ELAN y hace que la herramienta delegue el procesamiento de video al middleware sin tocar la base del software; gracias a ello, ELAN conserva intacto su ciclo de trabajo mientras el procesamiento especializado queda encapsulado en un componente externo, en línea con el propósito con el que fue concebido el protocolo AVATecH \cite{auer2014avatech}.

En la Figura~\ref{fig:panel_elan} se aprecia el panel de configuración \textbf{AIOrchestrationRecognizerPanel} tal como aparece dentro de ELAN cuando el usuario selecciona el reconocedor \textit{AI Orchestration Middleware Bridge}. Dicho panel traduce en controles visibles parámetros que hasta entonces solo podían enviarse desde clientes externos ---dirección del servidor, video a analizar, modelo de IA instalado, nivel de anotación destino y formato de etiqueta---, de modo que la configuración necesaria para lanzar una inferencia quedó incorporada al mismo entorno donde el investigador revisa y edita sus anotaciones, sin pasos intermedios fuera de la herramienta.

\begin{figure}[H]
    \centering
    \includegraphics[width=0.95\textwidth]{Diagramas/Panel configuracion del reconocer integrado en ELAN.png}
    \caption{Panel de configuración del reconocedor integrado en ELAN}
    \label{fig:panel_elan}
\end{figure}

Por su parte, la Figura~\ref{fig:dialogo_modelos} recoge el diálogo \textbf{AIModelManagerDialog}, accesible desde el botón \textit{Gestionar modelos\ldots} del panel. Allí se concentran las operaciones de instalación, activación y desactivación de modelos sin que el usuario abandone el entorno de anotación: cada botón del diálogo corresponde a una operación HTTP del middleware que el investigador ya no necesita conocer, con lo cual la gestión del ciclo de vida de los modelos, resuelta internamente a través de endpoints REST, quedó expuesta al usuario final por medio de controles gráficos coherentes con la interfaz de ELAN.

\begin{figure}[H]
    \centering
    \includegraphics[width=0.95\textwidth]{Diagramas/Dialogo de gestion de modelos integrados en ELAN.png}
    \caption{Diálogo de gestión de modelos integrado en ELAN}
    \label{fig:dialogo_modelos}
\end{figure}

En cuanto al cliente HTTP \textbf{AIOrchestration\allowbreak{}MiddlewareClient}, implementado sobre \textbf{java.net.http.\allowbreak{}HttpClient}, este trabaja exclusivamente con HTTP/1.1 para asegurar la compatibilidad con Uvicorn y evitar que se negocie automáticamente HTTP/2 frente a un servidor que no lo ofrece por defecto \cite{java11_httpclient,uvicorn_docs}. Su método \textbf{installModel} construye a mano el cuerpo \textbf{multipart/form-data} que transporta los paquetes ZIP, sin recurrir a dependencias externas, mientras que \textbf{buildErrorMessage} convierte los códigos de error estructurados del middleware en mensajes legibles dentro de la interfaz gráfica de ELAN. La integración sigue el contrato AVATecH en los siguientes tres elementos: el reconocedor informa del progreso con RecognizerHost.setProgress(), entrega segmentos con insertSegmentationAsTier() y maneja la cancelación del usuario con stop(). Los pasos para compilar la clase con Maven, registrar la clase usando el mecanismo SPI de Java y cargarla en ELAN 7.1 se describen en el Anexo IV de manera que la integración puede reproducirse sobre el código fuente oficial de la herramienta.

% ── Objetivo 2 ──────────────────────────────────────────────────────────────
\subsubsection{Núcleo del orquestador: API REST, cola de jobs y gestión de contenedores Docker}

El segundo objetivo específico consistía en codificar el núcleo del orquestador con Python y hacer uso de la virtualización ligera mediante Docker, incorporando control de concurrencia y gestión del ciclo de vida de los contenedores para sostener la estabilidad del sistema durante inferencias simultáneas. En respuesta a este objetivo se desarrolló el servicio central del middleware que contiene la validación de solicitudes, el controlador (endpoints), la coordinación de jobs y la entrega de métricas operativas. Adoptar una API REST permitió que ELAN, Postman y cualquier otro cliente se comunicaran con el mismo contrato HTTP sin necesidad de que el lenguaje interno de cada componente se comunicara con el resto del servicio, lo cual está en línea con el principio de interfaces desacopladas propio de las arquitecturas basadas en recursos \cite{fielding2000rest,fastapi2026features}.

La API expone ocho endpoints organizados en cuatro grupos funcionales que se encuentran resumidos en la Tabla~\ref{tab:endpoints_api}. El servidor Uvicorn levanta el framework FastAPI y mantiene disponibles los endpoints \textbf{GET /health} y \textbf{GET /api/v1/metrics}, incluso mientras se está ejecutando una inferencia prolongada. A ello se suma que el resultado de cada job permanece consultable durante la sesión a través de \textbf{GET /api/v1/jobs/\{job\_id\}}, gracias a un almacén en memoria (\textbf{MemoryStore}) que conserva las respuestas completadas hasta el siguiente reinicio del servicio. Esto demuestra que el middleware no solo realiza inferencias, sino que también puede ser utilizado para monitorear el estado de un servicio mientras se utiliza el backend de IA.

Por otra parte, el componente \textbf{JobQueue} implementa una cola FIFO con límite configurable de concurrencia apoyándose en \textbf{collections.deque} y \textbf{threading.Event} que son primitivas estándar de Python empleadas para coordinar espera y notificación entre hilos \cite{python_threading}. El número de backends que pueden inferir a la vez lo decide el parámetro \textbf{MIDDLEWARE\_\allowbreak{}MAX\_\allowbreak{}CONCURRENT\_\allowbreak{}JOBS}, ajustable desde el \textbf{docker-compose.yml} descrito en el Anexo~III. Con el valor por defecto de~1, un solo contenedor carga pesos en memoria en cada instante, lo que reduce el riesgo de desbordamiento bajo cargas concurrentes; entretanto, mientras un job aguarda su turno, el resto de la API continúa respondiendo. La Tabla~\ref{tab:metricas_concurrencia} recoge los resultados de la prueba controlada con dos solicitudes simultáneas y deja ver que la segunda permanece en cola hasta que la primera libera el slot de ejecución.

\begin{table}[H]
\centering
\caption{Métricas del sistema durante la ejecución de dos inferencias simultáneas}
\label{tab:metricas_concurrencia}
\small
\renewcommand{\arraystretch}{1.2}
\begin{tabular}{@{}>{\raggedright\arraybackslash}p{4.5cm}>{\centering\arraybackslash}p{2.5cm}>{\raggedright\arraybackslash}p{5.5cm}@{}}
\toprule
\textbf{Métrica / Parámetro} & \textbf{Valor real} & \textbf{Interpretación} \\
\midrule
\texttt{total\_jobs} & 10 & Total de solicitudes procesadas en el histórico. \\
\texttt{completed\_jobs} & 7 & Inferencias finalizadas con éxito. \\
\texttt{active\_jobs} & 1 & Inferencia en ejecución activa en el contenedor. \\
\texttt{queued\_jobs} & 1 & Segunda inferencia en espera dentro de la cola. \\
\texttt{timeout\_jobs} & 1 & Trabajos que excedieron el tiempo límite establecido. \\
\texttt{failed\_jobs} & 0 & Trabajos finalizados con error crítico del sistema. \\
\texttt{average\_exec\_ms} & 33673.43 & Tiempo promedio de ejecución global [ms]. \\
\texttt{last\_exec\_ms} & 32428.00 & Tiempo de respuesta de la última inferencia [ms]. \\
\texttt{DOCKER\_TIMEOUT} & 1 & Contador de errores por expiración del contenedor. \\
\bottomrule
\end{tabular}
\end{table}

Los valores de la Tabla~\ref{tab:metricas_concurrencia} respaldan cuantitativamente tanto el aislamiento como el ordenamiento de peticiones del middleware. Al llegar dos solicitudes de inferencia en el mismo instante, los contadores internos pasaron a \texttt{active\_jobs = 1} y \texttt{queued\_jobs = 1}: el mecanismo de exclusión mutua retuvo la segunda petición en la cola FIFO mientras el contenedor procesaba la primera, con lo cual se evitó saturar los recursos del equipo. Los contadores, además, resultan coherentes entre sí, pues los 10 jobs del histórico se reparten entre 7 completados, 1 expirado por timeout y los 2 de la prueba en curso (1 activo y 1 en cola). El tiempo registrado en \texttt{last\_exec\_ms}, cercano a 32.4 segundos, aporta a su vez una línea base del rendimiento del modelo de reconocimiento integrado.

La Figura~\ref{fig:postman_health} corresponde al escenario E-01 de la validación final y presenta la respuesta del endpoint \textbf{GET /health} verificada desde Postman. Comprobar esta disponibilidad básica antecede, en el orden de las pruebas, a las operaciones más costosas del sistema, como la instalación de modelos o las inferencias sobre video.

\begin{figure}[H]
    \centering
    \includegraphics[width=0.95\textwidth]{Diagramas/Respuesta del endpoint GET -health.png}
    \caption{Respuesta del endpoint (GET /health) verificada desde Postman (escenario E-01)}
    \label{fig:postman_health}
\end{figure}

En la Figura~\ref{fig:postman_metrics} podemos ver la respuesta del endpoint \textbf{GET /api/v1/metrics} después de realizar dos inferencias exitosas y un tiempo de espera forzado. Los contadores \textbf{total\_jobs}, \textbf{completed\_jobs} y \textbf{timeout\_jobs} muestran que el middleware además de procesas solicitudes, conserva trazabilidad de cada inferencia. Esta información sirve como evidencia de observabilidad durante las validaciones finales. 


\begin{figure}[H]
    \centering
    \includegraphics[width=0.95\textwidth]{Diagramas/Metricas acumuladas del middleware tras la prueba de carga controlada.png}
    \caption{Métricas acumuladas del middleware tras la prueba de carga controlada (GET /api/v1/metrics)}
    \label{fig:postman_metrics}
\end{figure}

La gestión del ciclo de vida del contenedor, la segunda parte de este objetivo, se desarrolló en el componente \textbf{DockerRunner} y en el contrato de comunicación entre el middleware y el backend del modelo. Este resultado vincula la arquitectura de orquestación con el procesamiento real de video: el middleware no procesa directamente las señales; prepara el entorno, incorpora el backend de IA para validar su respuesta y adapta los segmentos al formato que ELAN puede insertar en la línea de tiempo.

En particular, el DockerRunner traduce la ruta del video desde el sistema host al punto de montaje interno del contenedor \textbf{/data/videos}, verifica la disponibilidad del backend mediante sondeos en GET /health con reintentos cada 250 ms y convierte la respuesta JSON del backend en objetos \textbf{TemporalSegment} validados por Pydantic. Conviene subrayar que el aporte del TIC no es el modelo de reconocimiento en sí, sino la capa que hace posible integrarlo: el middleware orquesta la ejecución sin conocer la arquitectura interna del modelo por lo que cualquier backend que respete el contrato del endpoint \textbf{/health} e \textbf{/infer}; y el formato del \textbf{manifest.json} puede ocupar ese lugar. Lo que corresponde al middleware es el post-procesamiento de las salidas, es decir, la validación y adaptación de la respuesta del backend en anotaciones semánticas compatibles con ELAN. Para la validación se empleó el backend \textbf{lsec\_bio\_gloss\_final\_v1}, cuyo desarrollo queda fuera del alcance de este trabajo dado que fue realizado por la persona encargada de la funcionalidad A. El pipeline interno que nos entregó fue adaptado para que funcione encapsulado en un contenedor Docker, con las dependencias aisladas tanto de ELAN como del middleware. El manifest completo con el que se instala este modelos se recoge en el Anexo~I.

La Figura~\ref{fig:postman_inferencia} muestra la respuesta de una inferencia exitosa realizada desde el cliente Postman hacia el endpoint \textbf{POST /api/v1/jobs/segment-video}, con segmentos detectados, marcas de tiempo, etiquetas de glosa, predicciones alternativas y trazabilidad de la ejecución. Su importancia radica en conectar el contrato REST con el resultado lingüístico esperado, confirmando que el backend devuelve suficiente información para crear anotaciones temporales compatibles con ELAN. Un ejemplo completo de esta respuesta, con la descripción campo por campo del objeto \textbf{trace} y de sus etapas, figura en el Anexo~II.

\begin{figure}[H]
    \centering
    \includegraphics[width=0.95\textwidth]{Diagramas/Respuesta de inferencia exitosa.png}
    \caption{Respuesta de inferencia exitosa desde Postman: segmentos con glosas y trazabilidad completa}
    \label{fig:postman_inferencia}
\end{figure}

Los tiempos reales observados en el campo \textbf{trace.stages} de la respuesta del middleware durante la inferencia de validación se resumen en la Tabla~\ref{tab:latencia_etapas}. Los valores corresponden a un video de LSEc con una duración de 8.51~segundos (510 fotogramas a 59.94~FPS) procesado en un equipo con CPU Intel Core i5 de 12.ª generación. Desagregar la latencia por etapas ayuda a distinguir el costo de validación, espera en cola, arranque del contenedor, comprobación de disponibilidad, inferencia y postprocesamiento; de esa desagregación se desprende que la sobrecarga introducida por la orquestación es marginal, pues la suma de validación, cola y postprocesamiento apenas alcanza unos pocos milisegundos, mientras que el tiempo restante transcurre dentro del backend y depende del modelo instalado, no del middleware. El arranque del contenedor registró 0~ms porque este ya se encontraba activo (\textit{warm start}), condición habitual después de la primera ejecución.

\begin{table}[H]
\centering
\caption{Latencia por etapa de inferencia (modelo lsec\_bio\_gloss\_final\_v1)}
\label{tab:latencia_etapas}
\small
\renewcommand{\arraystretch}{1.15}
\begin{tabular}{@{}>{\raggedright\arraybackslash}p{4.4cm}>{\centering\arraybackslash}p{3.0cm}>{\raggedright\arraybackslash}p{5.0cm}@{}}
\toprule
\textbf{Etapa} & \textbf{Latencia medida (ms)} & \textbf{Descripción} \\
\midrule
Validación del modelo          & 1        & Consulta al registro de modelos \\
Espera en cola (\texttt{QUEUED})   & 1        & Sin contención con 1 job activo \\
Arranque del contenedor         & 0        & Contenedor previamente activo (\textit{warm start}) \\
Comprobación de disponibilidad del backend        & 4        & Polling hasta disponibilidad \\
Inferencia en el backend        & 34\,181  & Procesamiento de 510 fotogramas \\
Postprocesamiento               & 1        & Validación y adaptación de respuesta \\
\textbf{Total (wall-clock)}    & \textbf{34\,198} & Desde \texttt{RECEIVED} hasta \texttt{COMPLETED} \\
\bottomrule
\end{tabular}
\end{table}

Como cierre del flujo, la Figura~\ref{fig:elan_anotaciones} presenta las anotaciones generadas por el sistema en la línea de tiempo de ELAN tras ejecutar el reconocedor desde el panel integrado, de acuerdo con los escenarios E-07 y E-08. Su valor probatorio está en que completa el ciclo: la inferencia ocurre fuera de ELAN, pero los resultados regresan como anotaciones temporales ubicadas en el nivel de anotación configurado. Las marcas de inicio y fin de cada anotación deben coincidir con los intervalos de señas visibles en el video procesado, ya que esa correspondencia confirma que la herramienta de anotación interpretó correctamente la segmentación temporal.

\begin{figure}[H]
    \centering
    \includegraphics[width=0.95\textwidth]{Diagramas/Anotaciones generadas automáticamente en ELAN por el reconocedor integrado.png}
    \caption{Anotaciones generadas automáticamente en ELAN por el reconocedor integrado (escenario E-07 y E-08)}
    \label{fig:elan_anotaciones}
\end{figure}

% ── Objetivo 3 ──────────────────────────────────────────────────────────────
\subsubsection{Validación técnica mediante pruebas funcionales}

El tercer objetivo específico, enfocado en validar el desempeño técnico del middleware mediante pruebas en escenarios de anotación de LSEc, se cubrió con quince escenarios de validación distribuidos en cuatro categorías: conectividad, gestión de modelos, inferencia y manejo de errores. La validación midió latencia de comunicación, consumo de recursos computacionales y recuperación ante fallos. Esta organización permitió comprobar tanto el comportamiento observable desde las interfaces visibles para el usuario como la sonsistencia interna de archivos, métricas y contenedores.

La Tabla~\ref{tab:escenarios_validacion} contiene los escenarios más representativos realizados en la validación final. Los identificadores E-01 a E-15 funcionan como etiquetas de trazabilidad que enlazan cada acción de prueba con su respuesta HTTP esperada, el resultado observado y la evidencia obtenida durante la ejecución de dicha prueba.

\begin{table}[H]
\centering
\caption{Resultados de escenarios de validación final (selección representativa)}
\label{tab:escenarios_validacion}
\scriptsize
\renewcommand{\arraystretch}{1.15}
\setlength{\tabcolsep}{4pt}

\begin{tabular}{
    >{\centering\arraybackslash}p{1.0cm}
    >{\raggedright\arraybackslash}p{3.2cm}
    >{\centering\arraybackslash}p{1.5cm}
    >{\centering\arraybackslash}p{1.6cm}
    >{\raggedright\arraybackslash}p{5.4cm}
}
\toprule
\textbf{ID} & \textbf{Escenario} & \textbf{HTTP esp.} & \textbf{Resultado} & \textbf{Observación} \\
\midrule

E-01 &
Comprobación de disponibilidad desde Postman &
200 &
$\checkmark$ &
Respuesta menor a 50 ms. \\

E-02 &
Comprobación de disponibilidad desde ELAN &
200 &
$\checkmark$ &
El panel muestra el estado de servidor disponible. \\

E-03 &
Instalación de modelo mediante ZIP &
200 &
$\checkmark$ &
Imagen Docker construida y registry actualizado. \\

E-04 &
Listado de modelos &
200 &
$\checkmark$ &
Modelo visible con estado \texttt{available}. \\

E-05 &
Activar o desactivar modelo &
200 &
$\checkmark$ &
La solicitud \texttt{PATCH} actualiza el estado en el registry. \\

E-06 &
Instalación duplicada &
409 &
$\checkmark$ &
Error controlado \texttt{MODEL\_ALREADY\_EXISTS}. \\

E-07 &
Inferencia exitosa desde ELAN &
200 &
$\checkmark$ &
Anotaciones generadas en el nivel de anotación \texttt{AUTO\_GLOSS\_SEGMENTS}. \\

E-08 &
Verificación visual de anotaciones &
N/A &
$\checkmark$ &
Segmentos alineados temporalmente en la línea de tiempo. \\

E-09 &
Inferencia exitosa desde Postman &
200 &
$\checkmark$ &
Campo \texttt{trace} con tiempos por etapa. \\

E-10 &
Video inexistente &
502 &
$\checkmark$ &
Error controlado \texttt{DOCKER\_INFERENCE\_ERROR}. \\

E-11 &
Modelo desactivado &
409 &
$\checkmark$ &
Error controlado \texttt{MODEL\_DISABLED}. \\

E-12 &
Middleware detenido desde ELAN &
N/A &
$\checkmark$ &
ELAN informa que el servicio no está disponible. \\

E-13 &
Timeout forzado con \texttt{timeout\_sec = 1} &
504 &
$\checkmark$ &
Error controlado \texttt{DOCKER\_TIMEOUT}. \\

E-14 &
Contenedor del modelo detenido &
200 &
$\checkmark$ &
DockerRunner reinicia el contenedor automáticamente. \\

E-15 &
Métricas después de fallos &
200 &
$\checkmark$ &
Contadores \texttt{timeout\_jobs} y \texttt{error\_counts} actualizados correctamente. \\

\bottomrule
\end{tabular}
\end{table}

Por otro lado, la Tabla~\ref{tab:consumo_recursos} condensa el consumo de recursos medido con \textbf{docker stats} durante la inferencia completa con el modelo \textbf{lsec\_bio\_gloss\_final\_v1}, procesando el mismo video de 8.51 segundos. La medición se realizó a nivel de contenedores porque el middleware y el backend consumen recursos de manera muy distinta: el primero mantiene un perfil bajo ($\sim$0.13\,\% de CPU y $\sim$52.45\,MiB de RAM), asociado estrictamente a la orquestación, la validación y las métricas, mientras que el gasto intensivo se concentra en el contenedor del backend, que utiliza $\sim$194.89\,\% de CPU (casi dos núcleos completos) y $\sim$1.4\,GiB de memoria para ejecutar el pipeline de inferencia visual \cite{docker_cli_stats}. Semejante asimetría respalda una decisión de diseño central del sistema: dado que el costo computacional vive en el backend, es posible detener o reemplazar modelos sin comprometer la disponibilidad del orquestador.

\begin{table}[H]
\centering
\caption{Consumo de recursos durante inferencia (medición real con \texttt{docker stats})}
\label{tab:consumo_recursos}
\small
\renewcommand{\arraystretch}{1.15}
\begin{tabular}{@{}>{\raggedright\arraybackslash}p{3.8cm}>{\centering\arraybackslash}p{2.8cm}>{\centering\arraybackslash}p{2.8cm}>{\raggedright\arraybackslash}p{3.2cm}@{}}
\toprule
\textbf{Contenedor} & \textbf{CPU (\%)} & \textbf{RAM} & \textbf{Observación} \\
\midrule
Middleware (\texttt{elan-ai-middleware}) & 0.13 & 52.45\,MiB & Rol de orquestación \\
Backend del modelo (\texttt{elan-ai-model-...}) & 194.89 & 1.4\,GiB & Pipeline de inferencia \\
\bottomrule
\end{tabular}
\end{table}

En el plano interno, los escenarios de caja blanca corroboraron la consistencia del estado después de las operaciones principales: el archivo \textbf{registry.json} reflejó la entrada correcta del manifest tras la instalación; el mecanismo de bootstrap manifests reconstruyó el registro de forma automática después de eliminar \textbf{registry.json} y reiniciar el middleware; la traducción de la ruta del video quedó asentada en los logs del backend bajo el prefijo \textbf{/data/videos}; y \textbf{docker inspect} verificó el montaje correcto del directorio de videos. Estas comprobaciones complementan los códigos HTTP porque examinan que el estado persistente y la infraestructura Docker permanezcan coherentes tras cada operación.

En conjunto, los tres bloques de resultados cubren el ciclo completo del sistema: el primero demuestra que ELAN puede delegar el procesamiento de videos sin perder su flujo de anotación y que los contratos definidos sostienen esa comunicación. El segundo, demuestra que el orquestador regula la carga, administra los contenedores y devuelve la inferencia convertida en anotaciones temporales utilizables por ELAN. Y el tercero, demuestra que el comportamiento se mantiene predecible tanto en operaciones normales como ante fallos. Los anexos I~a~IV complementan estos resultados con todo lo necesarios para reproducirlos nuevamente, es decir, presentan el contrato de empaquetado, un ejemplo íntegro de respuesta, la configuración de despliegue y el procedimiento de compilación e integración en ELAN.

\vspace{1\baselineskip}


\subsection{Conclusiones}

El presente trabajo confirmó que es posible incorporar servicios de inteligencia artificial al flujo de anotación de ELAN sin la necesidad de alterar de manera profunda su arquitectura. La solución desarrollada actúa como un software intermedio entre la herramienta de anotación y los modelos de reconocimiento, de tal modo que el investigador pueda conservar su entorno habitual de trabajo mientras accede a resultados automáticos que pueden ser revisados y corregidos. Por lo tanto, a partir de los objetivos planteados y de las validaciones realizadas, se establecen las siguientes conclusiones:

\begin{enumerate}
    \item El análisis de ELAN y de su mecanismo de extensión permitió comprender cómo debe incorporarse un proceso externo dentro de la línea de tiempo de anotación. Esta revisión fue fundamental para que la solución planteada no funcionara como una herramienta aislada, sino como un apoyo integrado al proceso que normalmente realiza un investigador. En lugar de reemplazar el criterio humano, el sistema entrega segmentos y etiquetas preliminares que pueden ser aceptadas, ajustadas o eliminadas dentro de ELAN. De esta manera, la integración respeta el uso lingüístico de la herramienta y mantiene la anotación final bajo responsabilidad del especialista.

    \item La arquitectura propuesta demostró que separar la anotación, la orquestación y la inferencia facilita la integración de tecnologías distintas sin comprometer la estabilidad del sistema principal. Por lo tanto, ELAN se mantiene como la herramienta de revisión y edición manual de los segmentos; el middleware se encarga de organizar las solicitudes, validar la información y coordinar la ejecución de las inferencias; mientras que los modelos procesan videos en entornos independientes. Esta división reduce el acoplamiento entre componentes evitando que las dependencias de los modelos o del middleware afecten a ELAN y proporciona una base más ordenada para incorporar nuevos modelos de reconocimiento en el futuro.
    
    \item En conclusión, el uso de una arquitectura local permitió proteger la confidencialidad de los videos y anotaciones del investigador, ya que la información no necesita salir de la estación de trabajo para ser procesada. Además, el control de ejecución de inferencias evita que varios procesos pesados compitan por los mismos recursos del sistema de forma no controlada. Este control es importante porque en escenarios de investigación, donde los equipos disponibles pueden no tener grandes capacidades de cómputo y donde la estabilidad de la herramienta de anotación es crucial, el sistema no colapsa los recursos disponibles dado que no permite ejecutar inferencias simultaneas.

    \item Las pruebas realizadas mostraron que el middleware responde de manera controlada ante situaciones esperadas y ante fallos comunes, como videos inexistentes, modelos no disponibles, tiempos de espera excedidos o contenedores detenidos. Las validaciones realizadas no se limitaron a comprobar respuestas exitosas, también permitió observar cómo el sistema conserva su estabilidad, informando errores y manteniendo la disponibilidad del servicio al poder realizar consultas de progreso y métricas. Esto confirma que la solución implementada no solo ejecuta inferencias, sino que también ofrece condiciones mínimas de trazabilidad y diagnóstico para un uso real.

    \item La principal contribución de este trabajo es una arquitectura local, extensible y desacoplada que acerca los modelos de inteligencia artificial al entorno cotidiano de anotación de LSEc. Es importante recalcar, que el sistema no pretende sustituir la revisión humana ni resolver por si solo el problema del reconocimiento automático,este sistema busca una vía práctica para que los resultados de los modelos lleguen al entorno de ELAN como anotaciones temporables editables. Con todo esto, se busca reducir tareas repetitivas, apoyar la exploración de material audiovisual y facilitar futuras investigaciones que busquen integrar nuevos modelos sin las necesidad de construir una nueva herramienta.
  
\end{enumerate}

\vspace{1\baselineskip}

\subsection{Recomendaciones}

A partir de los resultados obtenidos, el sistema puede seguir creciendo sin perder la idea central que guió este trabajo: mantener a ELAN como la herramienta principal de anotación y usar el middleware como un sistema de apoyo para integrar modelos externos de forma controlada. Por lo tanto, las siguientes recomendaciones proponen mejoras orientadas a fortalecer la estabilidad, facilitar el mantenimiento y aumentar la utilidad de la solución en proyectos reales de investigación o documentación de LSEc.

\begin{enumerate}
    \item Es importante mejorar la gestión de los recursos de ejecución para que el sistema pueda adaptarse mejor a equipos con diferentes capacidades. Una posible línea de trabajo es tener un planificador que evalúe si un modelo necesita CPU, GPU o una mayor cantidad de memoria antes de comenzar una inferencia. Así los modelos ligeros podrían ejecutarse con mayor flexibilidad, mientras que los modelos más complejos se mantendrían controlados para no saturar el equipo del investigador.

    \item El diseño cliente-servidor del middleware abre una amplia trayectoria de crecimiento dado que se podría mover el orquestador y los modelos a una máquina de alto rendimiento (Ejemplo: un servidor institucional con altas capacidades), que actúe como un servidor de inferencia compartido. Dado que la comunicación entre ELAN y el middleware es por HTTP, muchas estaciones de trabajo en el mismo laboratorio pueden enviar sus solicitudes a este servidor apuntando la URL del reconocedor a la dirección IP correspondiente, sin modificar el código en ninguno de los lados. Esta idea se puede replicar en caso de querer migrar el servicio a la nube, siempre y cuando se consideren los aspectos ligados a la soberanía de los datos.
    
    \item Es aconsejable fortalecer el almacenamiento local de la información de modelos instalados. Aunque el archivo actual cumple su función, una base de datos ligera permitiría registrar los cambios de forma más segura, facilitaría la consulta de versiones y reduciría el riesgo de inconsistencias si el servicio falla de manera inesperada. Esta mejora debería mantenerse simple y local, de modo que no complique la instalación de nuevos modelos ni contradiga el principio de soberanía de los datos. 

    \item Se recomienda ampliar el conjunto de modelos probados con el sistema desarrollado. Para ello, es importante realizar nuevos modelos de reconocimiento de LSEc e instalarlos cumpliendo con los parámetros solicitados. Esto con la intención de comprobar la flexibilidad real del sistema frente a modelos con distintos enfoques, por ejemplo aquellos orientados a señas específicas, glosas más amplias o información no manual como expresiones faciales y postura corporal.
    
    \item También sería importante mejorar la interacción entre el investigador y los resultados generados de forma automática por algún modelo. Una extensión futura del panel del reconocedor en ELAN podría presentar no solo la etiqueta principal, sino también alternativas de glosas en caso de que un modelo las entregue. Esto ayudaría a que la revisión sea más transparente y útil porque el especialista tendría más información para decidir si acepta una anotación, la corrige o simplemente la descarta antes de incorporarla definitivamente en la línea de tiempo.

    \item Finalmente, se recomienda complementar la validación técnica con una evaluación lingüística a cargo de especialistas en LSEc, que compare las anotaciones generadas automáticamente con anotaciones de referencia y considere no solo si la etiqueta propuesta es correcta, sino también si los tiempos de inicio y fin corresponden con el evento observado en el video. Sus resultados permitirían medir el aporte real del sistema en tareas de documentación lingüística y orientar mejoras tanto en los modelos como en la manera en que sus salidas se presentan al investigador.
\end{enumerate}

\section[REFERENCIAS BIBLIOGRÁFICAS]{Referencias bibliográficas}
% Durante la redacción, esta línea muestra todas las entradas del archivo .bib.
% Para la versión final IEEE, comenta o elimina \nocite{*} y deja únicamente las referencias citadas en el texto.
\nocite{*}
\printbibliography[heading=none]

% ============================================================
%  CAPÍTULO 5 — ANEXOS
%  Los anexos continúan la numeración arábiga de páginas del
%  cuerpo del documento (no se regresa a números romanos).
%  Cada anexo tiene un solo título (sin negrilla) y se desarrolla
%  en prosa continua, sin subtítulos internos. Las tablas y los
%  códigos se numeran por anexo: Tabla I.1, Código II.2, etc.
% ============================================================
\clearpage
\section[ANEXOS]{Anexos}

Este capítulo reúne el material complementario del trabajo. Los Anexos I a IV amplían el detalle técnico de la metodología y de los resultados: el contrato de empaquetado de modelos, un ejemplo íntegro de respuesta de inferencia, la configuración de despliegue del middleware y el procedimiento de compilación e integración del reconocedor en ELAN. Los Anexos V a VII recogen los enlaces a los repositorios donde se publicó el código fuente y a los manuales de uso y de desarrollo.

\section*{\normalfont ANEXO I. Esquema completo del archivo manifest.json}
\addcontentsline{toc}{section}{ANEXO I. Esquema completo del archivo manifest.json}
\label{anx:manifest_schema}

\renewcommand{\thetable}{I.\arabic{table}}
\renewcommand{\theHtable}{anexoI.table.\arabic{table}}
\setcounter{table}{0}
\renewcommand{\thecodigo}{I.\arabic{codigo}}
\renewcommand{\theHcodigo}{anexoI.codigo.\arabic{codigo}}
\setcounter{codigo}{0}
\renewcommand{\thecomando}{I.\arabic{comando}}
\renewcommand{\theHcomando}{anexoI.comando.\arabic{comando}}
\setcounter{comando}{0}
\renewcommand{\thefigure}{I.\arabic{figure}}
\renewcommand{\theHfigure}{anexoI.figure.\arabic{figure}}
\setcounter{figure}{0}

El archivo \textbf{manifest.json} funciona como el contrato de empaquetado de cada modelo. Debe ubicarse en la raíz del paquete ZIP para que el middleware pueda reconocerlo, validar su estructura e instalarlo de forma controlada. Durante la instalación, el servicio registra su contenido en \textbf{data/models\_store/registry.json} y guarda una copia en \textbf{data/bootstrap\_manifests/}; de esta manera, el registro puede reconstruirse automáticamente si se pierde o se reinicia el entorno.

El documento JSON se organiza en campos de primer nivel que describen la identidad del modelo, su forma de ejecución, sus artefactos y los valores que verá el usuario en ELAN: \textbf{model\_id} (identificador único dentro del registro), \textbf{name} (nombre legible mostrado en ELAN), \textbf{version} (versión semántica, por ejemplo \textbf{1.0.0}), \textbf{task} (tipo de tarea que realiza el modelo), \textbf{runtime} (modo de ejecución y tipo de runner), \textbf{artifacts} (mapa de rutas relativas a los artefactos del paquete), \textbf{input\_contract} y \textbf{output\_contract} (medio de entrada y tipo de salida), \textbf{container} (configuración de red del contenedor, obligatoria cuando se usa el runner \textbf{docker\_http}) y \textbf{ui} (valores predeterminados para el panel de configuración de ELAN). La Tabla~\ref{tab:manifest_root} detalla cada campo del objeto raíz, su tipo y su obligatoriedad.

\begin{table}[H]
\centering
\caption{Campos del objeto raíz del manifest.json}
\label{tab:manifest_root}
\scriptsize
\renewcommand{\arraystretch}{1.08}
\setlength{\tabcolsep}{3.5pt}

\begin{tabular}{
    >{\raggedright\arraybackslash}p{3.0cm}
    >{\raggedright\arraybackslash}p{1.6cm}
    >{\centering\arraybackslash}p{1.8cm}
    >{\raggedright\arraybackslash}p{6.3cm}
}
\toprule
\textbf{Campo} & \textbf{Tipo} & \textbf{Requerido} & \textbf{Descripción} \\
\midrule
\texttt{model\_id} & string & Sí & Identificador único; admite letras, dígitos, punto, guion y guion bajo (\texttt{[A-Za-z0-9\_.-]}). \\
\texttt{name} & string & Sí & Nombre descriptivo mostrado en la interfaz de ELAN. \\
\texttt{version} & string & Sí & Versión semántica del modelo, por ejemplo, \texttt{X.Y.Z}. \\
\texttt{task} & string & Sí & Tipo de tarea que resuelve el modelo. \\
\texttt{runtime} & object & Sí & Configuración del modo de ejecución. \\
\texttt{artifacts} & object & Sí & Mapa de rutas relativas a los artefactos del paquete. \\
\texttt{input\_contract} & object & Sí & Restricciones sobre el medio de entrada (tipo de medio, características). \\
\texttt{output\_contract} & object & Sí & Tipo de salida y clases esperadas. \\
\texttt{backend\_config} & object & No & Datos usados durante la instalación: nombre y tag de la imagen Docker, puerto y timeout de arranque inicial (180~s en el modelo de referencia). \\
\texttt{container} & object & Cond. & Obligatorio si \texttt{runtime.runner = docker\_http}. \\
\texttt{ui} & object & Sí & Valores predeterminados para el panel de ELAN; sus campos internos son opcionales. \\
\bottomrule
\end{tabular}
\end{table}

Respecto del campo \textbf{task}, el esquema \textbf{ModelManifest} del middleware acepta exactamente dos valores; cualquier otro provoca el rechazo del paquete durante la validación. El valor \textbf{video\_segmentation} indica que el modelo detecta regiones temporales en el video sin clasificar su contenido semántico, mientras que \textbf{video\_segmentation\_and\_gloss\_classification} indica que, además de detectar las regiones, el modelo asigna una glosa de lengua de señas a cada segmento; este último valor corresponde al modelo de referencia del trabajo.

El objeto \textbf{runtime} declara cómo debe ejecutarse el backend. Sus tres campos se resumen en la Tabla~\ref{tab:manifest_runtime}: el modo de ejecución (actualmente solo se soporta \textbf{docker}), el tipo de contenedor (reservado para extensiones futuras) y el identificador del runner del middleware encargado de iniciar el contenedor y llamar a los endpoints del backend.

\begin{table}[H]
\centering
\caption{Campos del objeto (runtime)}
\label{tab:manifest_runtime}
\scriptsize
\renewcommand{\arraystretch}{1.08}
\setlength{\tabcolsep}{3.5pt}

\begin{tabular}{
    >{\raggedright\arraybackslash}p{2.5cm}
    >{\raggedright\arraybackslash}p{1.7cm}
    >{\raggedright\arraybackslash}p{2.7cm}
    >{\raggedright\arraybackslash}p{6.5cm}
}
\toprule
\textbf{Campo} & \textbf{Tipo} & \textbf{Valor} & \textbf{Descripción} \\
\midrule
\texttt{mode} & string & \texttt{docker} & Modo de ejecución del backend. Actualmente, solo se soporta \texttt{docker}. \\
\texttt{framework} & string & \texttt{container} & Tipo de contenedor. Se mantiene reservado para extensiones futuras. \\
\texttt{runner} & string & \texttt{docker\_http} & Identifica el runner del middleware encargado de iniciar el contenedor y llamar a los endpoints \texttt{/health} e \texttt{/infer}. \\
\bottomrule
\end{tabular}
\end{table}

En cuanto al objeto \textbf{artifacts}, el middleware no interpreta el contenido de los artefactos; su función es conservar las rutas declaradas para que el \textbf{Dockerfile} del paquete pueda construir la imagen del backend con los archivos correctos. Las claves del mapa actúan como nombres semánticos definidos por el autor del modelo, mientras que los valores apuntan a rutas relativas dentro del paquete. Durante la instalación se comprueba que cada archivo declarado exista dentro del ZIP y que ninguna ruta sea absoluta ni contenga \textbf{../}; si alguna condición falla, el paquete se rechaza antes de construir la imagen. La Tabla~\ref{tab:manifest_artifacts} enumera los artefactos del modelo de referencia y su propósito.

\begin{table}[H]
\centering
\caption{Artefactos del modelo de referencia (lsec\_bio\_gloss\_final\_v1)}
\label{tab:manifest_artifacts}
\scriptsize
\renewcommand{\arraystretch}{1.08}
\setlength{\tabcolsep}{4pt}

\begin{tabular}{
    >{\raggedright\arraybackslash}p{3.1cm}
    >{\raggedright\arraybackslash}p{9.8cm}
}
\toprule
\textbf{Clave} & \textbf{Ruta relativa y propósito} \\
\midrule
\texttt{weights\_bio} & \texttt{weights/best\_bio\_segmenter\_v2.pt}: pesos del modelo BiLSTM de segmentación BIO. \\
\texttt{weights\_gloss} & \texttt{weights/best\_keypoint\_transformer\_v11.pt}: pesos del clasificador Transformer. \\
\texttt{vocab} & \texttt{vocab/gloss\_vocab\_top20.csv}: vocabulario de 20 glosas LSEc. \\
\texttt{config} & \texttt{config/pipeline\_config.json}: hiperparámetros del pipeline. \\
\texttt{dockerfile} & \texttt{backend/Dockerfile}: instrucciones de construcción de la imagen Docker. \\
\texttt{requirements} & \texttt{backend/requirements.txt}: dependencias Python del backend. \\
\texttt{app\_main} & \texttt{backend/app/main.py}: punto de entrada FastAPI del servicio de inferencia. \\
\bottomrule
\end{tabular}
\end{table}

El objeto \textbf{container} es obligatorio cuando \textbf{runtime.runner = docker\_http}, porque reúne los datos que el \textbf{DockerRunner} necesita para iniciar el backend, comprobar su disponibilidad y enviarle las solicitudes de inferencia. Sus campos se describen en la Tabla~\ref{tab:manifest_container}.

\begin{table}[H]
\centering
\caption{Campos del objeto (container)}
\label{tab:manifest_container}
\scriptsize
\renewcommand{\arraystretch}{1.08}
\setlength{\tabcolsep}{3.5pt}

\begin{tabular}{
    >{\raggedright\arraybackslash}p{3.3cm}
    >{\raggedright\arraybackslash}p{1.5cm}
    >{\raggedright\arraybackslash}p{2.4cm}
    >{\raggedright\arraybackslash}p{6.0cm}
}
\toprule
\textbf{Campo} & \textbf{Tipo} & \textbf{Valor típico} & \textbf{Descripción} \\
\midrule
\texttt{internal\_port} & integer & \texttt{8080} & Puerto expuesto por el backend dentro del contenedor. \\
\texttt{host\_port} & integer & \textit{opcional} & Puerto en el host para depuración; se recomienda omitirlo en producción. \\
\texttt{health\_path} & string & \texttt{/health} & Ruta del endpoint de comprobación de disponibilidad. \\
\texttt{infer\_path} & string & \texttt{/infer} & Ruta del endpoint de inferencia. \\
\texttt{startup\_timeout\_sec} & integer & \texttt{30} & Tiempo máximo, en segundos, que el middleware espera hasta que \texttt{/health} responda con estado 200 después de iniciar el contenedor. \\
\texttt{service\_url} & string & \textit{opcional} & URL base del servicio cuando el contenedor se ejecuta fuera del control del SDK. Si está presente, el runner omite la gestión del ciclo de vida. \\
\texttt{name} & string & \textit{opcional} & Nombre fijo del contenedor. Si se omite, Docker asigna un nombre automáticamente. \\
\texttt{environment} & object & \textit{opcional} & Variables de entorno pasadas al contenedor en tiempo de ejecución. \\
\texttt{volumes} & object & \textit{opcional} & Bind mounts adicionales del modelo. El middleware agrega automáticamente el volumen de videos. \\
\bottomrule
\end{tabular}
\end{table}

Por último, los valores del objeto \textbf{ui} se cargan automáticamente en el panel \textbf{AIOrchestrationRecognizerPanel} cuando el usuario selecciona un modelo en ELAN. Con ello se evita repetir configuraciones básicas en cada ejecución y se mantiene una experiencia más cercana al flujo habitual de anotación. Sus campos se resumen en la Tabla~\ref{tab:manifest_ui}.

\begin{table}[H]
\centering
\caption{Campos del objeto (ui)}
\label{tab:manifest_ui}
\scriptsize
\renewcommand{\arraystretch}{1.08}
\setlength{\tabcolsep}{4pt}

\begin{tabular}{
    >{\raggedright\arraybackslash}p{3.5cm}
    >{\raggedright\arraybackslash}p{1.7cm}
    >{\raggedright\arraybackslash}p{8.0cm}
}
\toprule
\textbf{Campo} & \textbf{Tipo} & \textbf{Descripción} \\
\midrule
\texttt{default\_label} & string & Etiqueta usada cuando \texttt{label\_mode = simple}. También actúa como valor de respaldo si el modelo no devuelve una glosa. \\
\texttt{default\_target\_tier} & string & Nombre del nivel de anotación de ELAN donde se insertarán las anotaciones. \\
\texttt{label\_mode} & string & Modo de generación de etiquetas. El valor \texttt{gloss\_top1} usa la glosa de mayor probabilidad, mientras que \texttt{simple} usa siempre \texttt{default\_label}. \\
\texttt{supports\_threshold} & boolean & Indica si el panel muestra un control de umbral de confianza. Actualmente, no es utilizado por el backend de referencia. \\
\bottomrule
\end{tabular}
\end{table}

Para cerrar el anexo, el Código~\ref{fig:manifest_ejemplo} presenta el \textbf{manifest.json} completo del modelo de referencia utilizado en la validación del sistema. Al incluirse en la raíz del paquete ZIP, este archivo se convierte en la fuente de información que guía la instalación, el registro del modelo y la configuración inicial mostrada en ELAN. Nótese la diferencia entre los dos timeouts de arranque: \textbf{backend\_config} usa 180 segundos porque durante la instalación el backend carga los pesos por primera vez, mientras que \textbf{container} usa 30 segundos porque en cada inferencia el contenedor ya debería estar preparado.

\begin{center}
\begin{minipage}{0.96\textwidth}
\hrule
\vspace{0.06cm}
\scriptsize
\begin{Verbatim}[baselinestretch=0.68]
{
  "model_id": "lsec_bio_gloss_final_v1",
  "name": "LSEC BIO Gloss Pipeline - Implementacion Final Tesis",
  "version": "1.0.0",
  "task": "video_segmentation_and_gloss_classification",
  "runtime": {
    "mode": "docker",
    "framework": "container",
    "runner": "docker_http"
  },
  "artifacts": {
    "weights_bio": "weights/best_bio_segmenter_v2.pt",
    "weights_gloss": "weights/best_keypoint_transformer_v11.pt",
    "vocab": "vocab/gloss_vocab_top20.csv",
    "config": "config/pipeline_config.json",
    "dockerfile": "backend/Dockerfile",
    "requirements": "backend/requirements.txt",
    "app_main": "backend/app/main.py"
  },
  "input_contract": {
    "media_type": "video",
    "feature_type": "keypoints",
    "input_dim": 356
  },
  "output_contract": {
    "type": "segments_with_gloss",
    "classes": ["O", "B", "I"],
    "top_k": 5
  },
  "ui": {
    "default_label": "LSEC_REGION",
    "default_target_tier": "AUTO_GLOSS_SEGMENTS",
    "label_mode": "gloss_top1",
    "supports_threshold": false
  },
  "backend_config": {
    "docker_image_name": "lsec-bio-gloss-final",
    "docker_image_tag": "1.0.0",
    "container_port": 8080,
    "health_path": "/health",
    "infer_path": "/infer",
    "startup_timeout_sec": 180
  },
  "container": {
    "internal_port": 8080,
    "health_path": "/health",
    "infer_path": "/infer",
    "startup_timeout_sec": 30
  }
}
\end{Verbatim}
\vspace{0.03cm}
\hrule

\captionof{codigo}{Manifest del modelo (lsec\_bio\_gloss\_final\_v1) v1.0.0}
\label{fig:manifest_ejemplo}
\end{minipage}
\end{center}

% ============================================================
%  Anexo II — Ejemplo completo de respuesta JSON de inferencia
% ============================================================

\section*{\normalfont ANEXO II. Ejemplo completo de respuesta JSON de inferencia}
\addcontentsline{toc}{section}{ANEXO II. Ejemplo completo de respuesta JSON de inferencia}
\label{anx:respuesta_json}

\renewcommand{\thetable}{II.\arabic{table}}
\renewcommand{\theHtable}{anexoII.table.\arabic{table}}
\setcounter{table}{0}
\renewcommand{\thecodigo}{II.\arabic{codigo}}
\renewcommand{\theHcodigo}{anexoII.codigo.\arabic{codigo}}
\setcounter{codigo}{0}
\renewcommand{\thecomando}{II.\arabic{comando}}
\renewcommand{\theHcomando}{anexoII.comando.\arabic{comando}}
\setcounter{comando}{0}
\renewcommand{\thefigure}{II.\arabic{figure}}
\renewcommand{\theHfigure}{anexoII.figure.\arabic{figure}}
\setcounter{figure}{0}

Este anexo muestra una respuesta exitosa del endpoint \textbf{POST /api/v1/jobs/segment-video} del middleware. El ejemplo conserva los campos definidos en el esquema \textbf{SegmentVideoResponse} del archivo \textbf{middleware/app/schemas/jobs.py} y corresponde a una inferencia del modelo \textbf{lsec\_bio\_gloss\_final\_v1} v1.0.0 sobre un video corto de prueba de la LSEc. La intención es dejar claro cómo viajan los resultados desde el backend de IA hasta el formato que luego puede transformarse en anotaciones de ELAN.

El intercambio comienza con la solicitud que el cliente envía al middleware, reproducida en el Código~\ref{fig:request_inferencia}.

\begin{center}
\begin{minipage}{0.96\textwidth}
\hrule
\vspace{0.06cm}
\scriptsize
\begin{Verbatim}[baselinestretch=0.68]
POST http://127.0.0.1:8000/api/v1/jobs/segment-video
Content-Type: application/json

{
  "job_id": "job-elan-20260515143022512",
  "media": {
    "path": "/data/videos/lsec_validacion_001.mp4"
  },
  "annotation": {
    "target_tier": "AUTO_GLOSS_SEGMENTS",
    "default_label": "LSEC_REGION",
    "label_mode": "gloss_top1"
  },
  "model": {
    "model_id": "lsec_bio_gloss_final_v1",
    "version": "1.0.0"
  },
  "execution": {
    "device_preference": "auto",
    "runner": "auto",
    "timeout_sec": 300
  }
}
\end{Verbatim}
\vspace{0.03cm}
\hrule

\captionof{codigo}{Request de inferencia enviado al endpoint (POST /api/v1/jobs/segment-video)}
\label{fig:request_inferencia}
\end{minipage}
\end{center}

La respuesta correspondiente, con estado HTTP 200, incluye cinco campos de primer nivel: \textbf{job\_id}, \textbf{status}, \textbf{media\_info}, \textbf{segments} y \textbf{trace}. Los cuatro primeros identifican el trabajo, su estado, el video procesado y los segmentos obtenidos. El campo \textbf{trace}, por su parte, agrega el historial de estados, los tiempos de cada etapa del pipeline y los metadatos necesarios para observar la ejecución del contenedor Docker. Por brevedad, el Código~\ref{fig:respuesta_inferencia} muestra solo el primero de los cuatro segmentos detectados (\textbf{n\_detected\_segments}); los restantes siguen exactamente la misma estructura.

\begin{center}
\begin{minipage}{0.9\textwidth}
\hrule
\vspace{0.09cm}
\tiny
\begin{Verbatim}[baselinestretch=0.92]
HTTP/1.1 200 OK
Content-Type: application/json

{
  "job_id": "job-elan-20260515143022512",
  "status": "COMPLETED",

  "media_info": {
    "fps": 25.0,
    "duration_ms": 9840,
    "total_frames": 246
  },

  "segments": [
    {
      "segment_id": 1,
      "start_ms": 480,
      "end_ms": 1320,
      "start_frame": 12,
      "end_frame": 33,
      "duration_frames": 21,
      "label": "YO",
      "confidence": 0.9231,
      "predictions": [
        {"rank": 1, "gloss_id": 0,  "gloss": "YO",      "probability": 0.9231},
        {"rank": 2, "gloss_id": 6,  "gloss": "PERSONA", "probability": 0.0412},
        {"rank": 3, "gloss_id": 12, "gloss": "HOMBRE",  "probability": 0.0189},
        {"rank": 4, "gloss_id": 2,  "gloss": "QUERER",  "probability": 0.0091},
        {"rank": 5, "gloss_id": 9,  "gloss": "PENSAR",  "probability": 0.0077}
      ]
    }
  ],
  "trace": {
    "runner": "docker_http",
    "device": "container",
    "model_id": "lsec_bio_gloss_final_v1",
    "model_version": "1.0.0",
    "output_type": "segments_with_gloss",
    "exec_ms": 18472,
    "stages": {
      "validation_ms": 12,
      "queue_ms": 0,
      "container_start_ms": 3210,
      "healthcheck_ms": 842,
      "docker_inference_ms": 14356,
      "inference_ms": 14356,
      "keypoint_extraction_ms": 7821,
      "bio_model_load_ms": 390,
      "gloss_model_load_ms": 218,
      "vocab_load_ms": 44,
      "bio_inference_ms": 4103,
      "bio_postprocessing_ms": 312,
      "gloss_classification_ms": 1468,
      "postprocessing_ms": 52,
      "total_ms": 18472
    },
    "state_history": [
      "RECEIVED",
      "VALIDATING",
      "PREPROCESSING",
      "QUEUED",
      "RUNNING",
      "POSTPROCESSING",
      "COMPLETED"
    ],
    "fps": 25.0,
    "total_frames": 246,
    "sampled_frames": 246,
    "windows_count": 30,
    "original_width": 640,
    "original_height": 480,
    "n_detected_segments": 4,
    "docker_mode": "sdk_managed",
    "docker_image": "lsec-bio-gloss-final:1.0.0",
    "container_id": "a3f7b91c2d84e056f1a2b3c4d5e6f789abcd1234ef015678cd9012ab34d63456",
    "container_start_ms": 3210,
    "healthcheck_ms": 842,
    "docker_inference_ms": 14356,
    "total_ms": 18472
  }
}
\end{Verbatim}
\vspace{0.03cm}
\hrule

\captionof{codigo}{Respuesta completa de inferencia exitosa HTTP 200}
\label{fig:respuesta_inferencia}
\end{minipage}
\end{center}

Para interpretar la respuesta, la Tabla~\ref{tab:trace_campos} describe cada campo del objeto \textbf{trace}, que concentra la información de trazabilidad de la ejecución.

\begin{table}[H]
\centering
\caption{Campos del objeto \texttt{trace} en la respuesta de inferencia}
\label{tab:trace_campos}
\scriptsize
\renewcommand{\arraystretch}{1.08}
\setlength{\tabcolsep}{4pt}

\begin{tabular}{
    >{\raggedright\arraybackslash}p{4.2cm}
    >{\raggedright\arraybackslash}p{8.8cm}
}
\toprule
\textbf{Campo} & \textbf{Descripción} \\
\midrule
\texttt{runner} & Identificador del runner del middleware, por ejemplo, \texttt{docker\_http}. \\
\texttt{device} & Tipo de dispositivo de ejecución reportado por el runner, por ejemplo, \texttt{container}. \\
\texttt{model\_id} & Identificador del modelo que procesó el job. \\
\texttt{model\_version} & Versión del modelo. \\
\texttt{output\_type} & Tipo de salida producida, por ejemplo, \texttt{segments\_with\_gloss}. \\
\texttt{exec\_ms} & Tiempo total de ejecución en el middleware, expresado en milisegundos. \\
\texttt{stages} & Desglose de latencia por etapa del procesamiento. \\
\texttt{state\_history} & Lista ordenada de estados que atravesó el job. \\
\texttt{fps / total\_frames} & Metadatos del video procesado, propagados desde el backend. \\
\texttt{n\_detected\_segments} & Número de segmentos temporales detectados. \\
\texttt{docker\_mode} & Modo de gestión del contenedor, por ejemplo, \texttt{sdk\_managed} o \texttt{compose\_service}. \\
\texttt{docker\_image} & Imagen Docker utilizada por el backend de inferencia. \\
\texttt{container\_id} & Identificador del contenedor Docker. \\
\texttt{container\_start\_ms} & Tiempo empleado en iniciar el contenedor, expresado en milisegundos. \\
\texttt{healthcheck\_ms} & Tiempo hasta que \textbf{GET /health} respondió HTTP 200, expresado en milisegundos. \\
\texttt{docker\_inference\_ms} & Tiempo de la llamada \textbf{POST /infer} dentro del contenedor, expresado en milisegundos. \\
\bottomrule
\end{tabular}
\end{table}

Dentro de \textbf{trace}, el objeto \textbf{stages} corresponde al esquema \textbf{StageMetrics} definido en \textbf{middleware/app/schemas/metrics.py}. Su utilidad principal es separar la latencia total en dos niveles: las etapas propias del orquestador y las etapas internas del pipeline de IA. Esta separación facilita identificar si el tiempo de respuesta se consume en validación, espera, comunicación con Docker o procesamiento del modelo. El significado de cada campo se presenta en la Tabla~\ref{tab:stages_campos}.

\begin{table}[H]
\centering
\caption{Significado de los campos del objeto (stages)}
\label{tab:stages_campos}
\scriptsize
\renewcommand{\arraystretch}{1.08}
\setlength{\tabcolsep}{4pt}

\begin{tabular}{
    >{\raggedright\arraybackslash}p{4.8cm}
    >{\raggedright\arraybackslash}p{8.2cm}
}
\toprule
\textbf{Campo} & \textbf{Descripción} \\
\midrule
\texttt{validation\_ms} & Validación del cuerpo JSON del request mediante Pydantic. \\
\texttt{queue\_ms} & Tiempo de espera en la cola FIFO de jobs. \\
\texttt{container\_start\_ms} & Arranque del contenedor Docker mediante el SDK. \\
\texttt{healthcheck\_ms} & Espera activa hasta que \texttt{/health} responde 200. \\
\texttt{docker\_inference\_ms} & Tiempo total de la llamada \textbf{POST /infer} al backend. \\
\texttt{keypoint\_extraction\_ms} & Extracción de keypoints con MediaPipe Holistic. \\
\texttt{bio\_model\_load\_ms} & Carga del modelo BiLSTM a memoria. \\
\texttt{gloss\_model\_load\_ms} & Carga del Transformer de clasificación de glosas. \\
\texttt{vocab\_load\_ms} & Lectura del vocabulario CSV de glosas. \\
\texttt{bio\_inference\_ms} & Inferencia del segmentador BIO mediante ventana deslizante. \\
\texttt{bio\_postprocessing\_ms} & Suavizado y filtrado de etiquetas BIO. \\
\texttt{gloss\_classification\_ms} & Clasificación Transformer de cada segmento detectado. \\
\texttt{postprocessing\_ms} & Construcción de la respuesta final en el middleware. \\
\texttt{total\_ms} & Tiempo total de extremo a extremo. \\
\bottomrule
\end{tabular}
\end{table}

% ============================================================
%  Anexo III — Configuración completa del docker-compose.yml
% ============================================================

\section*{\normalfont ANEXO III. Configuración del docker-compose.yml del middleware}
\addcontentsline{toc}{section}{ANEXO III. Configuración del docker-compose.yml del middleware}
\label{anx:docker_compose}

\renewcommand{\thetable}{III.\arabic{table}}
\renewcommand{\theHtable}{anexoIII.table.\arabic{table}}
\setcounter{table}{0}
\renewcommand{\thecodigo}{III.\arabic{codigo}}
\renewcommand{\theHcodigo}{anexoIII.codigo.\arabic{codigo}}
\setcounter{codigo}{0}
\renewcommand{\thecomando}{III.\arabic{comando}}
\renewcommand{\theHcomando}{anexoIII.comando.\arabic{comando}}
\setcounter{comando}{0}
\renewcommand{\thefigure}{III.\arabic{figure}}
\renewcommand{\theHfigure}{anexoIII.figure.\arabic{figure}}
\setcounter{figure}{0}

El archivo \textbf{docker-compose.yml} reúne la configuración necesaria para ejecutar el middleware como servicio local y conectarlo con los contenedores de modelos de IA. Es el mecanismo recomendado para el despliegue, porque concentra en un solo lugar el ciclo de vida del contenedor, las variables de entorno, los bind mounts de persistencia y la red compartida \textbf{elan-ai-shared}. El Código~\ref{fig:docker_compose} reproduce el archivo completo utilizado durante el desarrollo y la validación.

\begin{center}
\begin{minipage}{0.96\textwidth}
\hrule
\vspace{0.06cm}
\scriptsize
\begin{Verbatim}[baselinestretch=0.68]
services:
  middleware:
    build:
      context: .
    image: elan-ai-middleware:0.1.0
    container_name: elan-ai-middleware
    ports:
      - "127.0.0.1:8000:8000"
    environment:
      MIDDLEWARE_HOST: 0.0.0.0
      MIDDLEWARE_PORT: 8000
      MIDDLEWARE_LOG_LEVEL: INFO
      MIDDLEWARE_RUNTIME_PROFILE: final
      MIDDLEWARE_MODELS_STORE_DIR: /app/data/models_store
      # Bootstrap: on every startup the middleware scans this directory
      # and auto-registers any manifest.json found there.
      # Combined with the bind mount below, installed model manifests
      # survive even a full volume reset.
      MIDDLEWARE_BOOTSTRAP_MANIFESTS_DIR: /app/data/bootstrap_manifests
      # Directory on the HOST where videos are stored.
      # Mounted read-only at /data/videos inside every model container.
      # Use forward slashes even on Windows (Docker Desktop translates them).
      MIDDLEWARE_VIDEOS_DIR: "C:/Users/imbaq/OneDrive/Desktop"
      # Shared Docker network used for container-to-container communication.
      # Model containers are attached to this network so the middleware
      # can reach them by container name (no host-port mapping required).
      MIDDLEWARE_DOCKER_NETWORK: elan-ai-shared
    volumes:
      # BIND MOUNTS: live on the host filesystem and are NEVER removed
      # by "docker compose down -v", "docker volume prune" or Docker Desktop.
      - ./data/models_store:/app/data/models_store
      - ./data/bootstrap_manifests:/app/data/bootstrap_manifests
      # Mount the Docker socket so the middleware can build/run model containers.
      - /var/run/docker.sock:/var/run/docker.sock
    networks:
      - shared
    restart: unless-stopped

# No named volumes: everything is on the host filesystem via bind mounts.

networks:
  shared:
    # Fixed name so model containers can join by this exact name
    # regardless of the docker-compose project prefix.
    name: elan-ai-shared
\end{Verbatim}
\vspace{0.03cm}
\hrule

\captionof{codigo}{Archivo (docker-compose.yml) del middleware}
\label{fig:docker_compose}
\end{minipage}
\end{center}

Las variables de entorno del servicio se describen en la Tabla~\ref{tab:env_vars}. Las variables que no aparecen declaradas en el archivo, como \textbf{MIDDLEWARE\_MAX\_CONCURRENT\_JOBS}, toman su valor por defecto desde \textbf{app/core/config.py}; para modificarlas basta con añadirlas a la sección \textbf{environment} del servicio y reiniciar el contenedor.

\begin{table}[H]
\centering
\caption{Variables de entorno del servicio middleware}
\label{tab:env_vars}
\tiny
\renewcommand{\arraystretch}{1.08}
\setlength{\tabcolsep}{4pt}

\begin{tabular}{
    >{\raggedright\arraybackslash}p{4.6cm}
    >{\raggedright\arraybackslash}p{3.2cm}
    >{\raggedright\arraybackslash}p{5.6cm}
}
\toprule
\textbf{Variable} & \textbf{Valor por defecto} & \textbf{Descripción} \\
\midrule
\texttt{MIDDLEWARE\_HOST} &
\texttt{0.0.0.0} &
Interfaz en que Uvicorn acepta conexiones dentro del contenedor. \\

\texttt{MIDDLEWARE\_PORT} &
\texttt{8000} &
Puerto de escucha de Uvicorn. \\

\texttt{MIDDLEWARE\_LOG\_LEVEL} &
\texttt{INFO} &
Nivel de log de la aplicación FastAPI. \\

\texttt{MIDDLEWARE\_RUNTIME\_PROFILE} &
\texttt{final} &
Perfil de ejecución. En producción se utiliza \texttt{final}. \\

\texttt{MIDDLEWARE\_MODELS\_STORE\_DIR} &
\path{/app/data/models_store} &
Ruta al directorio del registro de modelos y archivos de instalación. \\

\texttt{MIDDLEWARE\_BOOTSTRAP\_MANIFESTS\_DIR} &
\path{/app/data/bootstrap_manifests} &
Directorio escaneado al iniciar para reconstruir el registro automáticamente. \\

\texttt{MIDDLEWARE\_VIDEOS\_DIR} &
\textit{requerido} &
Ruta del host donde residen los videos. Se monta como \texttt{/data/videos} en los contenedores de modelo. \\

\texttt{MIDDLEWARE\_DOCKER\_NETWORK} &
\texttt{elan-ai-shared} &
Nombre de la red Docker compartida. Los contenedores de modelo se conectan a esta red para ser alcanzables por su nombre. \\

\texttt{MIDDLEWARE\_MAX\_CONCURRENT\_JOBS} &
\texttt{1} &
Número máximo de inferencias simultáneas. Controla el acceso concurrente a GPU, CPU o memoria del host. \\
\bottomrule
\end{tabular}
\end{table}

La configuración utiliza únicamente bind mounts, es decir, directorios del equipo anfitrión montados dentro del contenedor. Esta decisión evita depender de volúmenes Docker con nombre y facilita revisar los archivos directamente desde el sistema operativo. Además, los datos de modelos y manifests se conservan incluso después de detener el servicio o limpiar volúmenes desde Docker. Los tres montajes utilizados y su propósito se detallan en la Tabla~\ref{tab:bind_mounts}.

\begin{table}[H]
\centering
\caption{Bind mounts del servicio middleware}
\label{tab:bind_mounts}
\tiny
\renewcommand{\arraystretch}{1.08}
\setlength{\tabcolsep}{3.5pt}

\begin{tabular}{
    >{\raggedright\arraybackslash}p{4.3cm}
    >{\raggedright\arraybackslash}p{4.3cm}
    >{\raggedright\arraybackslash}p{4.8cm}
}
\toprule
\textbf{Ruta host} & \textbf{Ruta contenedor} & \textbf{Propósito} \\
\midrule
\texttt{./data/models\_store} &
\texttt{/app/data/models\_store} &
Registro de modelos instalados, archivo \texttt{registry.json} y archivos del paquete de cada modelo. \\

\texttt{./data/bootstrap\_manifests} &
\path{/app/data/bootstrap_manifests} &
Copias de los manifests de modelos para reconstrucción automática del registro al reiniciar. \\

\texttt{/var/run/docker.sock} &
\texttt{/var/run/docker.sock} &
Socket del daemon Docker. Permite al middleware construir imágenes y gestionar contenedores de modelos mediante el SDK de Python. \\
\bottomrule
\end{tabular}
\end{table}

En cuanto a la red, \textbf{elan-ai-shared} conserva un nombre fijo, independiente del prefijo que Docker Compose asigna al proyecto. Esto permite que los contenedores de modelo, creados dinámicamente por el SDK de Docker durante la inferencia, se unan a la misma red sin configuraciones adicionales. Una vez conectados, el middleware puede comunicarse con cada backend por su nombre de contenedor mediante la resolución DNS interna de Docker, sin exponer puertos innecesarios en el host. Los comandos más frecuentes durante la operación del servicio se resumen en el Comando~\ref{fig:compose_cmds}.

\begin{center}
\begin{minipage}{0.90\textwidth}
\hrule
\vspace{0.06cm}
\scriptsize
\begin{Verbatim}[baselinestretch=0.72]
# Construir la imagen e iniciar el middleware en segundo plano
docker compose up --build -d

# Ver logs en tiempo real
docker compose logs -f middleware

# Detener el middleware; los datos persisten en ./data/
docker compose down

# Detener y eliminar la imagen; los datos persisten en ./data/
docker compose down --rmi local

# Verificar que el middleware responde
curl http://127.0.0.1:8000/health

# Consultar modelos registrados
curl http://127.0.0.1:8000/api/v1/models

# Consultar metricas de operacion
curl http://127.0.0.1:8000/api/v1/metrics
\end{Verbatim}
\vspace{0.03cm}
\hrule

\captionof{comando}{Comandos Docker Compose más frecuentes durante la operación}
\label{fig:compose_cmds}
\end{minipage}
\end{center}

Respecto de la seguridad y la red, el puerto 8000 del middleware se expone únicamente en la interfaz de loopback del host, mediante \textbf{127.0.0.1:8000}. Con esta decisión, el servicio queda disponible para ELAN en la misma máquina, pero no para otros equipos de la red local. Esta restricción encaja con el escenario previsto del TIC, en el que el investigador ejecuta ELAN, el middleware y Docker en su propio equipo de trabajo. Si se quisiera que otras estaciones con ELAN consumieran el mismo middleware, como se plantea en las recomendaciones, bastaría con cambiar la publicación del puerto de \textbf{"127.0.0.1:8000:8000"} a \textbf{"8000:8000"} y apuntar la URL del panel del reconocedor a la dirección IP del servidor; en ese caso se recomienda añadir mecanismos de autenticación y cifrado si la red no es de confianza, además de ajustar \textbf{MIDDLEWARE\_MAX\_CONCURRENT\_JOBS} a la capacidad del equipo.

El socket de Docker montado en el servicio entrega al contenedor del middleware un nivel de control amplio sobre el daemon Docker del host. Para un entorno local y controlado, esta decisión simplifica la construcción y ejecución de modelos; sin embargo, en escenarios compartidos o multiusuario conviene restringir ese acceso mediante un proxy de Docker, como \textbf{Tecnativa Docker Socket Proxy}, o migrar a la API REST de Docker Engine con autenticación TLS.

% ============================================================
%  Anexo IV — Instrucciones de compilación e integración en ELAN 7.1
% ============================================================

\section*{\normalfont ANEXO IV. Compilación e integración del reconocedor en ELAN 7.1}
\addcontentsline{toc}{section}{ANEXO IV. Compilación e integración del reconocedor en ELAN 7.1}
\label{anx:compilacion_elan}

\renewcommand{\thetable}{IV.\arabic{table}}
\renewcommand{\theHtable}{anexoIV.table.\arabic{table}}
\setcounter{table}{0}
\renewcommand{\thecodigo}{IV.\arabic{codigo}}
\renewcommand{\theHcodigo}{anexoIV.codigo.\arabic{codigo}}
\setcounter{codigo}{0}
\renewcommand{\thecomando}{IV.\arabic{comando}}
\renewcommand{\theHcomando}{anexoIV.comando.\arabic{comando}}
\setcounter{comando}{0}
\renewcommand{\thefigure}{IV.\arabic{figure}}
\renewcommand{\theHfigure}{anexoIV.figure.\arabic{figure}}
\setcounter{figure}{0}

Este anexo describe cómo compilar el código fuente modificado de ELAN 7.1 e integrar el reconocedor de IA en una instalación funcional de la herramienta. La modificación incorpora el paquete \textbf{mpi.eudico.client.annotator.recognizer.ai} al árbol fuente de ELAN, manteniendo intactas las clases originales. Esta decisión reduce el riesgo de afectar funcionalidades existentes y facilita retirar o actualizar el reconocedor en el futuro.

Antes de compilar deben cumplirse los prerrequisitos de la Tabla~\ref{tab:prerrequisitos}. En particular, la variable de entorno \textbf{JAVA\_HOME} debe apuntar al JDK 21 antes de ejecutar Maven; esta verificación evita errores de compilación asociados a versiones de Java distintas de las declaradas en el proyecto.

\begin{table}[H]
\centering
\caption{Prerrequisitos para compilar ELAN 7.1 con el reconocedor de IA}
\label{tab:prerrequisitos}
\scriptsize
\renewcommand{\arraystretch}{1.08}
\setlength{\tabcolsep}{4pt}

\begin{tabular}{
    >{\raggedright\arraybackslash}p{3.3cm}
    >{\raggedright\arraybackslash}p{2.4cm}
    >{\raggedright\arraybackslash}p{7.4cm}
}
\toprule
\textbf{Herramienta} & \textbf{Versión mínima} & \textbf{Notas} \\
\midrule
JDK OpenJDK/Oracle &
21 &
El archivo \textbf{pom.xml} de ELAN especifica \textbf{source} y \textbf{target} en nivel 21. \\

Apache Maven &
3.8 &
Gestor de dependencias y ciclo de vida de construcción. \\

Git &
2.x &
Permite clonar o actualizar el repositorio fuente de ELAN. \\

Docker Desktop &
4.x &
Requerido en tiempo de ejecución para el middleware y los backends de modelos. \\
\bottomrule
\end{tabular}
\end{table}

Las clases Java añadidas se concentran en un solo paquete dentro del repositorio de ELAN. La Figura~\ref{fig:paquete_java_ai} muestra la ubicación del reconocedor, su panel de configuración, el cliente HTTP y los DTO utilizados para intercambiar datos con el middleware.

\begin{center}
\begin{minipage}{0.96\textwidth}
\hrule
\vspace{0.06cm}
\scriptsize
\begin{Verbatim}[baselinestretch=0.68]
elan-7.1/src/main/java/mpi/eudico/client/annotator/recognizer/ai/
|
+-- AIOrchestrationRecognizer.java          [implementa Recognizer AVATecH]
+-- AIOrchestrationRecognizerPanel.java     [panel de configuracion Swing]
+-- AIOrchestrationMiddlewareClient.java    [cliente HTTP del middleware]
+-- AIOrchestrationMiddlewareException.java [excepcion tipada]
+-- AIModelManagerDialog.java               [dialogo de gestion de modelos]
|
+-- dto/
    +-- SegmentVideoRequest.java            [DTO de request de inferencia]
    +-- SegmentVideoResponse.java           [DTO de respuesta de inferencia]
    +-- Segment.java                        [segmento temporal]
    +-- Prediction.java                     [glosa candidata con probabilidad]
    +-- MediaInfo.java                      [metadatos del video]
    +-- Trace.java                          [informacion de traza del job]
    +-- ModelSummary.java                   [resumen de modelo instalado]
    +-- InstalledModelDetail.java           [detalle completo de modelo]
    +-- ModelUiConfig.java                  [configuracion UI del modelo]
    +-- MediaPayload.java                   [sub-DTO del campo media]
    +-- AnnotationPayload.java              [sub-DTO del campo annotation]
    +-- ModelPayload.java                   [sub-DTO del campo model]
    +-- ExecutionPayload.java               [sub-DTO del campo execution]
    +-- ParametersPayload.java              [sub-DTO del campo parameters]
\end{Verbatim}
\vspace{0.03cm}
\hrule

\captionof{figure}{Paquete Java del reconocedor de IA dentro del árbol fuente de ELAN}
\label{fig:paquete_java_ai}
\end{minipage}
\end{center}

Sobre las dependencias, el cliente HTTP \textbf{AIOrchestrationMiddlewareClient} usa clases disponibles en el JDK 11+, en particular \textbf{java.net.http.HttpClient}, por lo que las llamadas al middleware no requieren librerías externas. Para el parseo de las respuestas JSON se utiliza \textbf{org.json}, dependencia que ya forma parte del archivo \textbf{pom.xml} de ELAN 7.1. En consecuencia, la integración no agrega nuevas dependencias al proyecto.

ELAN descubre los reconocedores disponibles mediante el mecanismo Service Provider Interface (SPI) de Java. Por ello, el nuevo reconocedor debe declararse en el archivo \textbf{src/main/resources/META-INF/services/mpi.eudico.client.annotator.recognizer.api.Recognizer} con la línea del Código~\ref{fig:spi_recognizer}, añadida al final del archivo existente sin eliminar otras implementaciones registradas. Al iniciar la función \textit{Run Recognizer}, ELAN lee este archivo y carga todos los reconocedores declarados.

\begin{center}
\begin{minipage}{0.92\textwidth}
\hrule
\vspace{0.06cm}
\scriptsize
\begin{Verbatim}[baselinestretch=0.72]
mpi.eudico.client.annotator.recognizer.ai.AIOrchestrationRecognizer
\end{Verbatim}
\vspace{0.03cm}
\hrule

\captionof{codigo}{Línea añadida al archivo SPI para registrar el reconocedor de IA}
\label{fig:spi_recognizer}
\end{minipage}
\end{center}

La compilación sigue el ciclo de vida estándar de Maven. Los comandos del Comando~\ref{fig:maven_build} deben ejecutarse desde el directorio raíz del proyecto \textbf{elan-7.1/}, de modo que Maven encuentre el archivo \textbf{pom.xml}, resuelva las dependencias y genere el JAR final. Si la compilación termina correctamente, Maven reporta \textbf{BUILD SUCCESS} y deja el archivo generado en \textbf{target/elan-7.1.jar}.

\begin{center}
\begin{minipage}{0.88\textwidth}
\hrule
\vspace{0.06cm}
\scriptsize
\begin{Verbatim}[baselinestretch=0.72]
# 1. Verificar version de Java; debe ser 21
java -version

# 2. Limpiar compilaciones anteriores
mvn clean

# 3. Compilar el arbol fuente completo
mvn compile

# 4. Crear el JAR sin lanzar ELAN
mvn package -Dexec.phase=install

# El JAR resultante queda en:
#   elan-7.1/target/elan-7.1.jar

# Las dependencias se copian en:
#   elan-7.1/target/lib/
\end{Verbatim}
\vspace{0.03cm}
\hrule

\captionof{comando}{Proceso de compilación y empaquetado con Maven}
\label{fig:maven_build}
\end{minipage}
\end{center}

Durante el desarrollo, ELAN puede ejecutarse directamente desde Maven con los comandos del Comando~\ref{fig:maven_run}. Esta opción permite probar el reconocedor integrado sin reemplazar todavía el JAR de una instalación existente.

\begin{center}
\begin{minipage}{0.74\textwidth}
\hrule
\vspace{0.06cm}
\scriptsize
\begin{Verbatim}[baselinestretch=0.72]
# Compilar y ejecutar ELAN directamente
mvn test

# Ejecutar sin recompilar, si ya se compiló
mvn exec:exec
\end{Verbatim}
\vspace{0.03cm}
\hrule

\captionof{comando}{Ejecución de ELAN durante el desarrollo sin construir el JAR}
\label{fig:maven_run}
\end{minipage}
\end{center}

Para distribuir ELAN con el reconocedor de IA integrado, se puede reemplazar el JAR original de una instalación existente siguiendo el Comando~\ref{fig:reemplazar_jar}. Antes de hacerlo, conviene guardar una copia de seguridad para poder restaurar la herramienta si el nuevo paquete no inicia correctamente. Debe considerarse, además, que si la instalación de ELAN utiliza un \textit{launcher} independiente, como ocurre con algunos instaladores de MPI, el nombre del JAR puede variar; por esta razón, antes de reemplazarlo se debe revisar qué archivo ejecuta realmente el script de lanzamiento o el acceso directo de ELAN.

\begin{center}
\begin{minipage}{0.92\textwidth}
\hrule
\vspace{0.06cm}
\scriptsize
\begin{Verbatim}[baselinestretch=0.72]
# Ruta de instalacion tipica de ELAN en Windows
# C:\Program Files\ELAN_7-1\

# Hacer copia de seguridad del JAR original
copy "C:\Program Files\ELAN_7-1\ELAN.jar" `
     "C:\Program Files\ELAN_7-1\ELAN.jar.bak"

# Copiar el nuevo JAR compilado
copy "elan-7.1\target\elan-7.1.jar" `
     "C:\Program Files\ELAN_7-1\ELAN.jar"
\end{Verbatim}
\vspace{0.03cm}
\hrule

\captionof{comando}{Sustitución del JAR en una instalación existente de ELAN en Windows}
\label{fig:reemplazar_jar}
\end{minipage}
\end{center}

Con ELAN ejecutándose desde el JAR modificado, la integración puede comprobarse siguiendo el flujo normal de uso del reconocedor:

\begin{enumerate}[noitemsep, topsep=2pt, parsep=0pt, partopsep=0pt]
    \item Abrir un archivo EAF en ELAN o crear uno nuevo.
    \item Navegar a \textbf{Recognizer} $\rightarrow$ \textbf{Run Recognizer\ldots}
    \item En el desplegable de reconocedores, verificar que aparece la entrada \textbf{AI Orchestration Middleware Bridge}.
    \item Seleccionar dicha entrada. El panel \textbf{AIOrchestrationRecognizerPanel} debe cargarse en la sección inferior del diálogo.
    \item Asegurarse de que el middleware está en ejecución mediante \textbf{docker compose up -d}.
    \item Pulsar el botón \textbf{Probar conexión} para verificar la conectividad con \textbf{GET /health}.
    \item Desde el botón \textbf{Gestionar modelos\ldots}, instalar el modelo \textbf{lsec\_bio\_gloss\_final\_v1} subiendo el paquete ZIP.
    \item Seleccionar el modelo instalado en el desplegable de modelos.
    \item Pulsar \textbf{Start} en el diálogo de reconocedores para iniciar la inferencia.
\end{enumerate}

Si la integración funciona correctamente, las anotaciones generadas por el modelo aparecerán al finalizar el proceso en un nuevo nivel de anotación llamado \textbf{AUTO\_GLOSS\_SEGMENTS}. Este resultado confirma que la inferencia se ejecutó fuera de ELAN, pero regresó a la herramienta como segmentos temporales editables por el investigador. Para finalizar, la Tabla~\ref{tab:troubleshooting} reúne los problemas más frecuentes observados durante la compilación y la integración, junto con su solución.

\begin{table}[H]
\centering
\caption{Problemas frecuentes durante la compilación e integración}
\label{tab:troubleshooting}
\scriptsize
\renewcommand{\arraystretch}{1.08}
\setlength{\tabcolsep}{4pt}

\begin{tabular}{
    >{\raggedright\arraybackslash}p{4.7cm}
    >{\raggedright\arraybackslash}p{8.4cm}
}
\toprule
\textbf{Problema} & \textbf{Solución} \\
\midrule
\textbf{JAVA\_HOME} no configurado &
Definir \textbf{JAVA\_HOME} apuntando al directorio raíz del JDK 21. \\

El reconocedor no aparece en ELAN &
Verificar que la línea SPI se añadió correctamente al archivo \textbf{META-INF/services/...Recognizer}. \\

\textbf{BUILD FAILURE: cannot find symbol} &
Verificar que todos los archivos del paquete \textbf{ai} y \textbf{ai/dto} están presentes en el árbol fuente. \\

ELAN no inicia tras reemplazar el JAR &
Restaurar \textbf{ELAN.jar.bak} y verificar que el nuevo JAR se compiló sin errores. \\

Panel muestra error de conexión &
Verificar que el middleware está en ejecución mediante \textbf{curl http://127.0.0.1:8000/health}. \\

Timeout durante la instalación del modelo &
El proceso de construcción de Docker puede tardar entre 5 y 10 minutos durante la primera instalación. El timeout del cliente está configurado en 300 segundos. \\

\textbf{422 Unprocessable Entity} al instalar &
El boundary de \textbf{multipart/form-data} no debe comenzar con \textbf{--}. Verificar la versión compilada de \textbf{AIOrchestrationMiddlewareClient.installModel()}. \\
\bottomrule
\end{tabular}
\end{table}

% ============================================================
%  Anexo V — Repositorio del código fuente del middleware
% ============================================================

\section*{\normalfont ANEXO V. Repositorio del código fuente del middleware}
\addcontentsline{toc}{section}{ANEXO V. Repositorio del código fuente del middleware}
\label{anx:repo_middleware}

El código fuente completo del middleware de orquestación desarrollado en este trabajo se encuentra publicado en el siguiente repositorio de GitHub:

% TODO: reemplazar por la URL real del repositorio
\begin{center}
\url{https://github.com/TPanchoX/Middleware.git}
\end{center}

El repositorio contiene la aplicación FastAPI organizada según la estructura de directorios presentada en la metodología (Figura~\ref{lst:estructura_directorios}), junto con el \textbf{Dockerfile} del servicio, el archivo \textbf{docker-compose.yml} reproducido en el Anexo~III y las dependencias de Python declaradas en \textbf{requirements.txt}. Con este material, cualquier lector puede clonar el proyecto y levantar el middleware en su propio equipo siguiendo los comandos del apartado de despliegue y configuración del entorno.

% ============================================================
%  Anexo VI — Repositorio del código fuente de ELAN 7.1 modificado
% ============================================================

\section*{\normalfont ANEXO VI. Repositorio del código fuente de ELAN 7.1 modificado}
\addcontentsline{toc}{section}{ANEXO VI. Repositorio del código fuente de ELAN 7.1 modificado}
\label{anx:repo_elan}

El árbol fuente de ELAN 7.1 con el reconocedor de IA incorporado está disponible en el siguiente repositorio de GitHub:

% TODO: reemplazar por la URL real del repositorio
\begin{center}
\url{https://github.com/TPanchoX/ELAN-7.1-TIC.git}
\end{center}

Este repositorio parte del código fuente oficial de ELAN 7.1 y agrega el paquete \textbf{mpi.eudico.client.annotator.recognizer.ai}, que reúne el reconocedor, su panel de configuración, el cliente HTTP y los DTO de comunicación con el middleware; las clases originales de la herramienta se mantienen intactas. El procedimiento para compilar este código con Maven, registrar el reconocedor mediante SPI y sustituir el JAR en una instalación existente es el descrito en el Anexo~IV.

% ============================================================
%  Anexo VII — Manuales de uso y de desarrollo
% ============================================================

\section*{\normalfont ANEXO VII. Manuales de uso y de desarrollo del sistema}
\addcontentsline{toc}{section}{ANEXO VII. Manuales de uso y de desarrollo del sistema}
\label{anx:manuales}

Los manuales elaborados como parte de este trabajo se encuentran disponibles en la siguiente carpeta compartida de OneDrive:

% TODO: reemplazar por la URL real de la carpeta compartida
\begin{center}
\url{https://epnecuador-my.sharepoint.com/:f:/g/personal/francisco_imbaquinga_epn_edu_ec/IgD8h8hnYtAoRKZSkhMi-MT7AcRfGOtbaeIAXFgJOTQBQXU?e=hxU3Qd}
\end{center}

La carpeta contiene dos documentos. El primero es el manual de usuario dirigido a investigadores, que explica paso a paso cómo poner en marcha el middleware, instalar modelos desde ELAN y ejecutar el reconocedor para obtener anotaciones automáticas. El segundo es el manual técnico dirigido a desarrolladores, que detalla cómo construir y empaquetar nuevos backends de modelos compatibles con el contrato de instalación del sistema. Ambos manuales complementan la información de los Anexos I a IV con un enfoque práctico y con capturas del uso real de la herramienta.

\end{document}
